% ================================================================
%  MÉMOIRE DE FIN D'ÉTUDES (PFE) — MASTER 2 MATHÉMATIQUES APPLIQUÉES
%  Université Hassan Premier — FST Settat
% ----------------------------------------------------------------
%  Sujet  : Apprentissage par Renforcement Profond pour
%                 l'Optimisation Combinatoire
%  Auteur  : Abdellah Kahlaoui
%  Année   : 2025–2026
% ================================================================



\documentclass[12pt, a4paper, twoside, openright]{book}

\usepackage[utf8]{inputenc}
\usepackage[T1]{fontenc}
\usepackage[numbers, sort&compress, square]{natbib}
\usepackage[french]{babel}
\usepackage{csquotes}
\usepackage{mathpazo}
\usepackage[scaled=0.85]{helvet}
\usepackage{courier}
\usepackage{microtype}

\usepackage[
a4paper,
top    = 2.5cm,
bottom = 2.5cm,
inner  = 3.2cm,
outer  = 2.2cm,
headheight = 15pt
]{geometry}

\usepackage[dvipsnames, table, x11names]{xcolor}

\definecolor{ThBlue}  {RGB}{0,  51, 102}
\definecolor{ThRed}   {RGB}{180,  0,   0}
\definecolor{ThGray}  {RGB}{245,245,245}
\definecolor{ThGreen} {RGB}{ 30,130,  30}
\definecolor{ThSilver}{RGB}{100,100,110}

\usepackage{amsmath, amssymb, mathtools, amsthm}
\usepackage{bm}
\usepackage{bbm}
\usepackage{dsfont}
\usepackage{stmaryrd}
\usepackage{esint}

\usepackage[most, listings]{tcolorbox}
\tcbuselibrary{theorems, breakable, skins, xparse}

\tcbset{
	ThPlain/.style={
		enhanced, breakable,
		colback=ThBlue!5!white, colframe=ThBlue,
		coltitle=white, fonttitle=\bfseries\upshape,
		attach boxed title to top left={yshift=-2mm, xshift=4mm},
		boxed title style={colback=ThBlue, colframe=ThBlue,
			rounded corners, sharp corners=east},
		separator sign none,
		description font=\normalfont\itshape,
		top=6pt, bottom=6pt,
	},
	ThDef/.style={
		enhanced, breakable,
		colback=ThRed!5!white, colframe=ThRed,
		coltitle=white, fonttitle=\bfseries\upshape,
		attach boxed title to top left={yshift=-2mm, xshift=4mm},
		boxed title style={colback=ThRed, colframe=ThRed,
			rounded corners, sharp corners=east},
		separator sign none,
		description font=\normalfont\itshape,
		top=6pt, bottom=6pt,
	},
	ThRem/.style={
		enhanced, breakable,
		colback=ThSilver!8!white, colframe=ThSilver,
		coltitle=ThSilver!30!black, fonttitle=\bfseries\upshape,
		attach boxed title to top left={yshift=-2mm, xshift=4mm},
		boxed title style={colback=ThSilver!15!white, colframe=ThSilver,
			rounded corners, sharp corners=east},
		separator sign none,
		description font=\normalfont\itshape,
		top=6pt, bottom=6pt,
	},
	ThEx/.style={
		enhanced, breakable,
		colback=ThGreen!5!white, colframe=ThGreen!70!black,
		coltitle=white, fonttitle=\bfseries\upshape,
		attach boxed title to top left={yshift=-2mm, xshift=4mm},
		boxed title style={colback=ThGreen!80!black, colframe=ThGreen!80!black,
			rounded corners, sharp corners=east},
		separator sign none,
		description font=\normalfont\itshape,
		top=6pt, bottom=6pt,
	},
}

\newcounter{thmcount}[chapter]
\renewcommand{\thethmcount}{\arabic{chapter}.\arabic{thmcount}}

\newtcbtheorem[use counter=thmcount]{theorem}    {Th\'eor\`eme}  {ThPlain}{thm}
\newtcbtheorem[use counter=thmcount]{lemma}      {Lemme}         {ThPlain}{lem}
\newtcbtheorem[use counter=thmcount]{proposition}{Proposition}   {ThPlain}{prop}
\newtcbtheorem[use counter=thmcount]{corollary}  {Corollaire}    {ThPlain}{cor}
\newtcbtheorem[use counter=thmcount]{definition} {D\'efinition}  {ThDef} {def}
\newtcbtheorem[use counter=thmcount]{assumption} {Hypoth\`ese}   {ThDef} {hyp}
\newtcbtheorem[use counter=thmcount]{remark}     {Remarque}      {ThRem} {rem}
\newtcbtheorem[use counter=thmcount]{example}    {Exemple}       {ThEx}  {ex}

\renewcommand{\qedsymbol}{$\blacksquare$}

\renewcommand{\thechapter}      {\Roman{chapter}}
\renewcommand{\thesection}      {\arabic{chapter}.\arabic{section}}
\renewcommand{\thesubsection}   {\arabic{chapter}.\arabic{section}.\arabic{subsection}}
\renewcommand{\thesubsubsection}%
{\arabic{chapter}.\arabic{section}.\arabic{subsection}.\arabic{subsubsection}}
\renewcommand{\theequation}     {\arabic{chapter}.\arabic{equation}}
\renewcommand{\thefigure}       {\arabic{chapter}.\arabic{figure}}
\renewcommand{\thetable}        {\arabic{chapter}.\arabic{table}}

\usepackage{titlesec}

\titleformat{\chapter}[display]
{\normalfont\bfseries\color{ThBlue}}
{\hfill\fontsize{72}{72}\selectfont\color{ThBlue!18}\thechapter}
{-8pt}
{{\color{ThBlue}\rule{\linewidth}{1.6pt}}\\[4pt]\Huge\bfseries\color{ThBlue}}
[{\vspace{4pt}\color{ThBlue}\rule{\linewidth}{0.4pt}}]

\titlespacing*{\chapter}{0pt}{-30pt}{36pt}

\titleformat{\section}
{\normalfont\Large\bfseries\color{ThBlue}}{\thesection}{1em}{}

\titleformat{\subsection}
{\normalfont\large\bfseries\color{ThBlue!75!black}}{\thesubsection}{1em}{}

\titleformat{\subsubsection}
{\normalfont\normalsize\bfseries\color{ThBlue!60!black}}
{\thesubsubsection}{1em}{}

\usepackage{fancyhdr}
\pagestyle{fancy}
\fancyhf{}
\fancyhead[LE]{\small\color{ThSilver}\nouppercase{\leftmark}}
\fancyhead[RO]{\small\color{ThSilver}\nouppercase{\rightmark}}
\fancyfoot[C]{\small\thepage}
\renewcommand{\headrulewidth}{0.4pt}
\renewcommand{\footrulewidth}{0pt}
\fancypagestyle{plain}{
	\fancyhf{}
	\fancyfoot[C]{\small\thepage}
	\renewcommand{\headrulewidth}{0pt}
}

\usepackage{graphicx}
\usepackage{tikz}
\usepackage{pgfplots}
\pgfplotsset{compat=1.18}

\usetikzlibrary{
	arrows.meta, positioning, shapes.geometric, calc,
	decorations.pathreplacing, decorations.markings,
	matrix, fit, backgrounds, shadows,
	automata, chains, trees
}

\tikzset{
	every node/.style = {font=\small},
	block/.style = {
		draw=ThBlue, fill=ThBlue!8, rounded corners=3pt,
		minimum width=2.6cm, minimum height=1.0cm,
		text centered, align=center, font=\small\bfseries
	},
	arrow/.style     = {-Stealth, thick, color=ThBlue},
	dasharrow/.style = {-Stealth, dashed, thick, color=ThSilver},
}

\usepackage{booktabs}
\usepackage{tabularx}
\usepackage{multirow}
\usepackage{array}
\usepackage{longtable}
\usepackage{rotating}

\usepackage[ruled, vlined, linesnumbered, french]{algorithm2e}
\SetAlgorithmName{Algorithme}{algorithme}{Liste des Algorithmes}
\SetKwInput{KwEntree}{Entr\'ee}
\SetKwInput{KwSortie}{Sortie}
\SetKwInput{KwInit}{Initialisation}
\SetKwComment{Comment}{$\triangleright$\ }{}
\SetAlFnt{\small}
\SetAlCapFnt{\small\bfseries}

\usepackage{listings}
\usepackage{upquote}
\lstset{
	language         = Python,
	basicstyle       = \ttfamily\small,
	keywordstyle     = \color{ThBlue}\bfseries,
	commentstyle     = \color{ThGreen}\itshape,
	stringstyle      = \color{ThRed},
	numberstyle      = \tiny\color{ThSilver},
	numbers          = left,
	stepnumber       = 1,
	backgroundcolor  = \color{ThGray},
	frame            = single,
	framerule        = 0.5pt,
	rulecolor        = \color{ThBlue!30},
	breaklines       = true,
	breakatwhitespace= true,
	showstringspaces = false,
	tabsize          = 4,
	captionpos       = b,
	xleftmargin      = 18pt,
	framexleftmargin = 18pt,
	morekeywords     = {as, with, yield, lambda, None, True, False,
		self, cls, import, from, pass, assert},
}

\usepackage{caption}
\usepackage{subcaption}
\captionsetup{
	font      = small,
	labelfont = {bf, color=ThBlue},
	labelsep  = period,
	margin    = 10pt,
}

\usepackage[intoc]{nomencl}
\makenomenclature
\renewcommand{\nomname}{Nomenclature et Notations}
\usepackage{etoolbox}
\renewcommand\nomgroup[1]{%
	\item[\bfseries\color{ThBlue}%
	\ifstrequal{#1}{A}{Ensembles et Espaces}{%
		\ifstrequal{#1}{B}{Vecteurs et Matrices}{%
			\ifstrequal{#1}{C}{Op\'erateurs et Fonctions}{%
				\ifstrequal{#1}{D}{Apprentissage par Renforcement}{%
					\ifstrequal{#1}{E}{Optimisation Combinatoire}{%
						\ifstrequal{#1}{Z}{Abr\'eviations}{}%
	}}}}}]%
}

\usepackage[
colorlinks        = true,
linkcolor         = ThBlue,
citecolor         = ThRed,
urlcolor          = ThBlue!70!black,
bookmarks         = true,
bookmarksnumbered = true,
pdfencoding       = auto,
pdftitle          = {Apprentissage par renforcement profond pour l'optimisation combinatoire},
pdfauthor         = {Abdellah Kahlaoui},
pdfsubject        = {Memoire PFE Master 2 Mathematiques Appliquees},
]{hyperref}

\hypersetup{
	pdfkeywords = {Deep Reinforcement Learning, Optimisation Combinatoire,
		TSP, Pointer Network, REINFORCE, heuristiques}
}

\usepackage[french]{cleveref}
\crefname{theorem}    {th\'eor\`eme}  {th\'eor\`emes}
\crefname{lemma}      {lemme}         {lemmes}
\crefname{proposition}{proposition}   {propositions}
\crefname{corollary}  {corollaire}    {corollaires}
\crefname{definition} {d\'efinition}  {d\'efinitions}
\crefname{assumption} {hypoth\`ese}   {hypoth\`eses}
\crefname{remark}     {remarque}      {remarques}
\crefname{example}    {exemple}       {exemples}
\crefname{equation}   {\'equation}   {\'equations}
\crefname{figure}     {figure}        {figures}
\crefname{table}      {tableau}       {tableaux}
\crefname{algorithm}  {algorithme}    {algorithmes}


\usepackage{setspace}
\onehalfspacing
\usepackage{emptypage}
\usepackage{float}
\usepackage{enumitem}
\usepackage{epigraph}
\setlength{\epigraphwidth}{0.6\textwidth}
\usepackage{url}



% ── Commandes mathematiques personnalisees ──────────────────────
\newcommand{\R}{\mathbb{R}}
\newcommand{\N}{\mathbb{N}}
\newcommand{\Z}{\mathbb{Z}}
\newcommand{\E}{\mathbb{E}}
\newcommand{\Prob}{\mathbb{P}}
\newcommand{\calS}{\mathcal{S}}
\newcommand{\calA}{\mathcal{A}}
\newcommand{\calR}{\mathcal{R}}
\newcommand{\calT}{\mathcal{T}}
\newcommand{\calL}{\mathcal{L}}
\newcommand{\calN}{\mathcal{N}}
\newcommand{\norm}[1]{\left\lVert #1 \right\rVert}
\newcommand{\abs}[1]{\left\lvert #1 \right\rvert}
\newcommand{\inner}[2]{\left\langle #1,\, #2 \right\rangle}
\DeclareMathOperator*{\argmax}{arg\,max}
\DeclareMathOperator*{\argmin}{arg\,min}
\newcommand{\ind}[1]{\mathbbm{1}\!\left[#1\right]}

\bibliographystyle{plainnat}

% ================================================================
%  DEBUT DU DOCUMENT
% ================================================================
\begin{document}
	
	% ------------------------------------------------------------
	%  MATIERE PRELIMINAIRE
	% ------------------------------------------------------------
	\frontmatter
	\pagestyle{empty}
	
	% ============================================================
	%  PAGE DE COUVERTURE
	% ============================================================
	\begin{titlepage}
		\begin{tikzpicture}[remember picture, overlay]
			\fill[ThBlue]
			(current page.north west) rectangle
			([yshift=-2.5cm] current page.north east);
			\fill[ThRed]
			([yshift=-2.5cm] current page.north west) rectangle
			([yshift=-2.7cm] current page.north east);
			\fill[ThBlue]
			(current page.south west) rectangle
			([yshift=1.5cm] current page.south east);
		\end{tikzpicture}
		
		
		\begin{center}
			
			\begin{tikzpicture}[remember picture, overlay]
				\node[anchor=north west, xshift=2.0cm, yshift=0.2cm] at (current page.north west)
				{\includegraphics[height=2.5cm]{logos/fst_settat.png}};
				\node[anchor=north east, xshift=-2.0cm, yshift=-0.25cm] at (current page.north east)
				{\includegraphics[height=1.8cm]{logos/hassan1er.png}};
			\end{tikzpicture}
			
			
			{\color{black}\large\bfseries
				UNIVERSIT\'E HASSAN I\textsuperscript{er} --- SETTAT}\\[2pt]
			{\color{black}\normalsize
				Facult\'e des Sciences et Techniques de Settat}\\[2pt]
			{\color{black}\small D\'epartement de Math\'ematiques}
			
			\vspace{0.4cm}
			{\color{ThBlue!25}\rule{0.45\linewidth}{0.5pt}}
			
			\vspace{0.4cm}
			{\large\color{ThBlue}\itshape M\'emoire de Fin d'\'Etudes}\\[3pt]
			{\normalsize\color{ThBlue}
				Master Sciences et Techniques --- Math\'ematiques Appliqu\'ees}\\[1pt]
			{\small\color{ThSilver} Ann\'ee Universitaire 2025--2026}
			
			\vspace{0.6cm}
			\begin{tcolorbox}[
				enhanced, breakable,
				colback=ThBlue!6, colframe=ThBlue,
				boxrule=1.2pt, arc=4pt,
				left=12pt, right=12pt, top=10pt, bottom=10pt,
				width=0.88\linewidth,
				]
				\centering
				{\LARGE\bfseries\color{ThBlue}
					Apprentissage par Renforcement \\[4pt]
					Profond pour l'Optimisation Combinatoire :\\[4pt]
					\'Etude Algorithmique du Probl\`eme\\[4pt]
					du Voyageur de Commerce}
			\end{tcolorbox}
			
			\vspace{0.8cm}
			{\large\bfseries R\'ealis\'e par :}\\[3pt]
			{\Large\color{ThBlue}\bfseries Abdellah KAHLAOUI}
			
			\vspace{0.6cm}
			{\large\bfseries Encadr\'e par :}\\[3pt]
			{\large [Nom à compléter]}
			
			\vspace{0.6cm}
			{\color{ThBlue!25}\rule{0.45\linewidth}{0.5pt}}
			
			\vspace{0.5cm}
			{\large\bfseries Membres du Jury :}\\[5pt]
			\begin{tabular}{ll}
				Pr. [Nom à compléter] & \textit{Pr\'esident}\\[2pt]
				Pr. [Nom à compléter] & \textit{Encadrant}\\[2pt]
				Pr. [Nom à compléter] & \textit{Examinateur}\\[2pt]
				Pr. [Nom à compléter] & \textit{Examinateur}\\
			\end{tabular}
			
			\vfill
			\begin{tikzpicture}[remember picture, overlay]
				\node[white, font=\small] at
				([yshift=0.75cm] current page.south)
				{Soutenu le [date à compléter] $\cdot$ Settat $\cdot$ Maroc};
			\end{tikzpicture}
		\end{center}
	\end{titlepage}
	\cleardoublepage
	
	% ============================================================
	%  DEDICACE
	% ============================================================
	\chapter*{D\'edicace}
	\addcontentsline{toc}{chapter}{D\'edicace}
	\thispagestyle{empty}
	\vspace*{4cm}
	\begin{flushright}
		\large\itshape
		À ma famille, pour son soutien et ses encouragements.\par
		\vspace{0.4cm}
		\normalsize\upshape
	\end{flushright}
	\newpage
	
	% ============================================================
	%  REMERCIEMENTS
	% ============================================================
	\chapter*{Remerciements}
	\addcontentsline{toc}{chapter}{Remerciements}
	\pagestyle{plain}
	Je tiens à remercier mon encadrant pour son accompagnement.
	\cleardoublepage
	
% ============================================================
%  RESUME
% ============================================================
\chapter*{R\'esum\'e}
\addcontentsline{toc}{chapter}{R\'esum\'e}

\begin{tcolorbox}[
enhanced,
colback=ThBlue!5,
colframe=ThBlue,
boxsep=2mm,
left=2mm,
right=2mm,
top=1mm,
bottom=1mm,
fontupper=\small,
]
Ce mémoire étudie le problème du voyageur de commerce (TSP) à travers
une approche fondée sur l'apprentissage par renforcement profond (DRL).
La construction d'une tournée est formulée comme un processus de décision
séquentielle. L'approche repose sur un Pointer Network doté d'un
encodeur-décodeur LSTM, entraîné à l'aide de REINFORCE avec une baseline
apprise. Trois entraînements indépendants, initialisés avec les graines
\texttt{20261013}, \texttt{20261014} et \texttt{20261015}, sont conduits
avec le même budget de 13\,000 mises à jour. Pour chaque entraînement,
le checkpoint est sélectionné sur un jeu de validation indépendant de
100 instances à 20 villes. Les trois modèles sont ensuite évalués sur les
mêmes 100 instances inédites pour chacune des tailles 20, 50 et 100 villes.
Les résultats sont comparés à trois références classiques : le plus proche
voisin (PPV), PPV amélioré par 2-opt et l'optimisation par colonies de
fourmis (ACO).

Les résultats moyens sur les trois graines dépendent de la taille des
instances. Pour 20 villes, ACO obtient la plus faible longueur moyenne
(3,8780), devant PPV+2-opt (3,9424) et DRL meilleur-sur-512 (3,9584).
Pour ce dernier mode, les trois intervalles bootstrap appariés par modèle
comparant DRL à PPV+2-opt croisent zéro ; cela ne prouve pas leur
équivalence. Pour 50 villes, les deux modes DRL surpassent PPV seul, mais
restent derrière PPV+2-opt (5,9712) et ACO (6,1568). Pour 100 villes,
les deux modes DRL sont dépassés par les trois références classiques ;
PPV+2-opt obtient le meilleur coût observé sur les 100 instances.

La variabilité observée entre les trois entraînements reste descriptive :
l'écart-type inter-graines de DRL meilleur-sur-512 varie de 0,0054 à
0,0293 selon la taille, mais trois graines ne suffisent pas à caractériser
précisément la distribution des performances de tous les entraînements.
Les conclusions portent donc sur l'architecture, les hyperparamètres,
le budget et la distribution d'instances étudiés. Les écarts sont calculés
par rapport au meilleur coût observé parmi les neuf configurations
comparées sur chaque instance, et non par rapport à un optimum exact connu.

\medskip
\textbf{Mots-cl\'es :} optimisation combinatoire, problème du voyageur
de commerce, apprentissage par renforcement profond, REINFORCE,
Pointer Network, heuristiques, robustesse inter-graines.
\end{tcolorbox}


% ============================================================
%  ABSTRACT
% ============================================================
\clearpage
\chapter*{Abstract}
\addcontentsline{toc}{chapter}{Abstract}

\begin{tcolorbox}[
enhanced,
colback=ThRed!4,
colframe=ThRed,
boxsep=2mm,
left=2mm,
right=2mm,
top=1mm,
bottom=1mm,
fontupper=\small,
]
This thesis studies the Traveling Salesperson Problem (TSP) through a
deep reinforcement learning (DRL) approach. Tour construction is
formulated as a sequential decision-making process. The method uses a
Pointer Network with an LSTM encoder--decoder, trained using REINFORCE
with a learned baseline. Three independent training runs, initialized
with seeds \texttt{20261013}, \texttt{20261014}, and \texttt{20261015},
use the same budget of 13,000 optimizer updates. A checkpoint is selected
independently for each run on a validation set of 100 instances with 20
cities. The three selected models are evaluated on the same 100 previously
unseen instances for each of 20, 50, and 100 cities. They are compared with
three classical references: nearest neighbor (NN), NN improved by 2-opt,
and ant colony optimization (ACO).

The mean results across the three training seeds depend on instance size.
For 20 cities, ACO obtains the lowest mean tour length (3.8780), followed
by NN+2-opt (3.9424) and DRL best-of-512 (3.9584). For this DRL decoding
mode, all three model-conditional paired-bootstrap intervals against
NN+2-opt cross zero; this is not evidence of equivalence. For 50 cities,
both DRL modes outperform plain NN but remain behind NN+2-opt (5.9712)
and ACO (6.1568). For 100 cities, both DRL modes are outperformed by all
three classical methods, and NN+2-opt has the lowest observed cost on all
100 test instances.

The observed variability across the three training runs is descriptive:
the between-seed standard deviation for DRL best-of-512 ranges from
0.0054 to 0.0293 across the tested sizes, but three seeds are insufficient
to characterize precisely the performance distribution over all possible
training runs. Conclusions are therefore limited to the architecture,
hyperparameters, training budget, baselines, and instance distribution
studied. Gaps are measured against the best observed result among the nine
compared configurations on each instance, not against a known exact optimum.

\medskip
\textbf{Keywords:} combinatorial optimization, Traveling Salesperson
Problem, deep reinforcement learning, REINFORCE, Pointer Network,
heuristics, between-seed robustness.
\end{tcolorbox}

\cleardoublepage
	
	% ============================================================
	%  TABLES
	% ============================================================
	\pagestyle{fancy}
	\tableofcontents
	\cleardoublepage
	
	\listoffigures
	\addcontentsline{toc}{chapter}{\listfigurename}
	\cleardoublepage
	
	\listoftables
	\addcontentsline{toc}{chapter}{\listtablename}
	\cleardoublepage
	
	\addcontentsline{toc}{chapter}{Liste des Algorithmes}
	\cleardoublepage
	
	% ============================================================
	%  NOMENCLATURE
	% ============================================================
	\nomenclature[A]{$\R$}{Corps des r\'eels}
	\nomenclature[A]{$\N$}{Entiers naturels}
	\nomenclature[A]{$\calS$}{Espace d'\'etats (MDP)}
	\nomenclature[A]{$\calA$}{Espace d'actions (MDP)}
	\nomenclature[B]{$\bm{x}$}{Vecteur colonne (minuscule gras)}
	\nomenclature[B]{$\mathbf{W}$}{Matrice (majuscule gras)}
	\nomenclature[C]{$\E[\cdot]$}{Esp\'erance math\'ematique}
	\nomenclature[C]{$\Prob(\cdot)$}{Mesure de probabilit\'e}
	\nomenclature[D]{$\pi_\theta$}{Politique param\'etr\'ee par $\theta$}
	\nomenclature[D]{$V^\pi(s)$}{Fonction valeur d'\'etat sous $\pi$}
	\nomenclature[D]{$Q^\pi(s,a)$}{Fonction valeur action-\'etat sous $\pi$}
	\nomenclature[D]{$\gamma$}{Facteur d'actualisation ($\gamma\in[0,1]$)}
	\nomenclature[D]{$r_t$}{R\'ecompense \`a l'instant $t$}
	\nomenclature[E]{$G=(V,E)$}{Graphe avec sommets $V$ et ar\^etes $E$}
	\nomenclature[E]{$c_{ij}$}{Co\^ut de l'arc $(i,j)$}
	\nomenclature[Z]{RL}{Reinforcement Learning}
	\nomenclature[Z]{DRL}{Deep Reinforcement Learning}
	\nomenclature[Z]{MDP}{Processus de D\'ecision Markovien}
	\nomenclature[Z]{TSP}{Travelling Salesman Problem}
	\nomenclature[Z]{GNN}{Graph Neural Network}
	\nomenclature[Z]{PPO}{Proximal Policy Optimization}
	\nomenclature[Z]{DQN}{Deep Q-Network}
	\nomenclature[Z]{ILP}{Integer Linear Programming}
	\nomenclature[Z]{ACO}{Ant Colony Optimization}
	\nomenclature[Z]{LKH}{Lin-Kernighan-Helsgaun}
	
	\printnomenclature[2.5cm]
	\cleardoublepage
	
	% ------------------------------------------------------------
	%  CORPS DU MEMOIRE
	% ------------------------------------------------------------
	\mainmatter
	\pagestyle{fancy}
	
	% ============================================================
	%  INTRODUCTION GENERALE
	% ============================================================
	\chapter*{Introduction G\'en\'erale}
	\addcontentsline{toc}{chapter}{Introduction G\'en\'erale}
	\markboth{Introduction G\'en\'erale}{Introduction G\'en\'erale}
	
\section*{Contexte et Motivation}

Le problème du voyageur de commerce (TSP) et, plus généralement, les problèmes de routage, sont des défis majeurs de la recherche logistique et opérationnelle d'aujourd'hui. Pendant longtemps, la résolution de ces problèmes NP-difficiles a été abordée soit par des approches exactes, qui se heurtent à des temps de calcul prohibitifs en pratique, soit par des heuristiques manuelles nécessitant une forte expertise.

L'apprentissage par renforcement profond (DRL) propose une autre voie : apprendre, à partir d'instances d'entraînement, une politique qui construit progressivement une solution. Des travaux fondateurs sur les réseaux Pointer et l'optimisation combinatoire par renforcement ont montré que cette approche peut produire des heuristiques apprises compétitives dans certains cadres expérimentaux~\cite{vinyals2015pointer,bello2016neural}. Ces performances dépendent toutefois de l'architecture, du protocole d'entraînement, de la taille des instances et du budget d'inférence ; elles ne constituent pas une garantie générale de qualité ni de rapidité.

Cependant, la littérature récente appelle à une analyse plus approfondie. De nombreuses revues comparatives mettent en évidence le principal inconvénient des solveurs basés sur le DRL : la généralisation. Un modèle entraîné sur de petits graphes peine souvent à conserver sa performance sur des instances à grande échelle, non vues pendant l'entraînement. De plus, ils n'offrent aucune garantie théorique d'optimalité.

C'est cette tension entre potentiel et limites qui motive le présent mémoire. L'étude porte sur la formulation du TSP comme problème de décision séquentielle et sur l'évaluation expérimentale d'un Pointer Network entraîné par REINFORCE avec une baseline apprise. Les performances de cette configuration sont comparées à celles d'heuristiques classiques, en qualité de tournée et en temps d'inférence, puis analysées à plusieurs tailles d'instances. PPO~\cite{schulman2017proximal} et DQN sont abordés uniquement dans le cadre théorique général lorsqu'ils sont pertinents ; ils ne font pas partie des méthodes implémentées et comparées dans l'expérience présentée.

\section*{Probl\'ematique}

Compte tenu de l'état actuel des recherches, l'application de l'apprentissage par renforcement profond à la recherche opérationnelle soulève plusieurs défis. Bien qu'un modèle neuronal entraîné puisse construire une solution sans relancer une recherche locale complète, son temps d'inférence et la qualité des tournées dépendent de l'architecture, du décodage et de la taille de l'instance. Transformer un problème combinatoire en un processus d'apprentissage par essais et erreurs n'est pas immédiat.

Ce changement de paradigme mène à la problématique centrale de cette thèse, qui s'articule autour des trois défis suivants :

\begin{enumerate}
	\item \textbf{Défi de modélisation :} Comment définir mathématiquement un problème d'optimisation combinatoire statique tel que le TSP (généralement modélisé par la programmation linéaire entière) comme un processus de décision séquentiel (MDP) respectant les contraintes du graphe ?
	\item \textbf{Défi algorithmique :} Comment un Pointer Network à encodeur-décodeur LSTM peut-il paramétrer une politique de construction de tournée, et comment l'algorithme REINFORCE avec baseline apprise peut-il l'entraîner sans disposer de tournées optimales étiquetées ?
	\item \textbf{Défi d'évaluation et de limites :} Quel compromis entre qualité de tournée et temps d'inférence observe-t-on pour cette configuration DRL par rapport au plus proche voisin, à 2-opt et à ACO ? Comment les performances évoluent-elles lorsque le nombre de villes dépasse la taille utilisée à l'entraînement ?
\end{enumerate}

Pour répondre à ces questions, il est nécessaire de mener une étude algorithmique rigoureuse, mêlant les mathématiques de l'optimisation discrète aux dernières avancées de l'intelligence artificielle.
	
	\section*{Objectifs du Projet}
	\begin{enumerate}[leftmargin=*]
		\item Formaliser le TSP euclidien symétrique et rappeler les enjeux de complexité des méthodes exactes et approchées.
		\item Décrire le cadre MDP et les principes de l'optimisation de politique utilisés pour construire une tournée valide.
		\item Adapter et documenter une implémentation du Pointer Network entraînée par REINFORCE avec baseline apprise.
		\item Comparer les résultats de cette configuration à ceux du plus proche voisin, de NN + 2-opt et d'ACO sur des instances de 20, 50 et 100 villes, en rapportant la qualité de solution et le temps d'inférence.
		\item Analyser avec prudence la généralisation et les limites du protocole, notamment la taille réduite du jeu de test, le nombre de pas d'entraînement et l'variabilité observée sur trois graines indépendantes.
	\end{enumerate}
	
	\section*{Organisation du M\'emoire}
	
	Pour répondre à cette problématique de façon progressive, ce mémoire s'articule autour de quatre chapitres principaux, encadrés par une introduction et une conclusion générales :
	
	\begin{itemize}
		\item \textbf{Le Chapitre I} pose le cadre mathématique fondamental de l'optimisation combinatoire. Dans ce chapitre, nous introduisons la théorie et les concepts de base de la complexité algorithmique, ainsi qu'un aperçu des méthodes classiques de résolution.
		
		\item \textbf{Le Chapitre II} propose une étude détaillée du problème du voyageur de commerce (TSP). Nous y définissons mathématiquement le problème, analysons ses propriétés et examinons les méthodes de résolution de pointe. Nous en montrons également les limites, particulièrement en termes de passage à l'échelle.
		
		\item \textbf{Le Chapitre III} étudie l'apprentissage par renforcement à travers le prisme théorique des processus de décision markoviens (MDP). Nous y présentons les algorithmes les plus avancés en apprentissage par renforcement profond (DRL) et discutons du changement de paradigme nécessaire pour formuler le TSP comme un processus de décision séquentielle.
		
		\item \textbf{Le Chapitre IV} présente l'implémentation du Pointer Network, le protocole d'entraînement et les résultats comparatifs ; il discute notamment les écarts observés et les limites de généralisation.
	\end{itemize}
	
	La conclusion générale synthétise les apports effectifs du travail, les limites de l'évaluation et les prolongements expérimentaux possibles.
	
	% ============================================================
	%  CHAPITRE I : FONDEMENTS DE L'OPTIMISATION COMBINATOIRE
	% ============================================================
	\chapter[Fondements de l'optimisation combinatoire]{Fondements de l'optimisation \\ combinatoire}
	\label{chap:fondements}
	
	L'optimisation combinatoire constitue l'un des piliers fondamentaux des mathématiques appliquées et de la recherche opérationnelle. Elle consiste à choisir la meilleure solution possible parmi un ensemble fini ou dénombrable d'options, en se basant sur la minimisation ou la maximisation d'une fonction objectif spécifique. 
	
	Si, au premier abord, il peut sembler aisé de faire un choix parmi un nombre fini de possibilités, c'est précisément cette nature << discrète >> qui engendre l'explosion combinatoire et rend la résolution de ces problèmes si coûteuse en temps de calcul.
	
	% ─────────────────────────────────────────────────────────────
	\section{Cadre mathématique}
	\label{sec:cadre_math}
	% ─────────────────────────────────────────────────────────────
	
	\begin{definition}{Problème d'optimisation combinatoire}{prob_optim_comb}
		Un problème d'optimisation combinatoire est un triplet
		$\mathcal{P} = (\mathcal{S}, f, \mathcal{C})$, où :
		\begin{itemize}
			\item $\mathcal{S}$ désigne l'\textbf{espace des solutions}
			(ou espace de recherche), constitué de l'ensemble de toutes
			les configurations discrètes envisageables,
			\item $f : \mathcal{S} \rightarrow \R$ est la
			\textbf{fonction objectif}, que l'on cherche à minimiser
			(ou maximiser),
			\item $\mathcal{C}$ représente un ensemble de
			\textbf{contraintes} définissant la sous-région admissible
			$\mathcal{F} \subseteq \mathcal{S}$, appelée ensemble des
			solutions réalisables.
		\end{itemize}
	\end{definition}
	\newpage
	L'objectif est de déterminer une solution $s^* \in \mathcal{F}$ telle que :
	$$
	s^* = \argmin_{s \in \mathcal{F}} f(s)
	$$
	ou, dans le cas d'un problème de maximisation,
	$s^* = \argmax_{s \in \mathcal{F}} f(s)$.
	
	\textbf{Espace de solutions discrètes.}
	À la différence de l'optimisation continue, où $\mathcal{S} \subseteq \R^n$
	est un ensemble connexe explorable par des méthodes de gradient, l'espace
	de solutions d'un problème combinatoire est construit à partir d'un ensemble
	fini d'éléments de base. Soit $E = \{e_1, e_2, \ldots, e_n\}$ un ensemble
	fini. L'espace de solutions est une famille particulière de sous-ensembles,
	de permutations, ou de structures construits sur $E$. \\
	Pour un problème où l'on doit déterminer l'ordre optimal de $n$ éléments (comme le choix d'un trajet entre $n$ villes), le nombre total de solutions possibles correspond à $n!$. Ce nombre grandit de manière si fulgurante à mesure qu'on ajoute des éléments qu'il devient rapidement impossible pour un ordinateur d'énumérer et de tester toutes les combinaisons : c'est ce qu'on appelle l'explosion combinatoire. 	
	
	\textbf{Fonction objectif.} La fonction objectif $f$ représente le critère que l'on cherche à optimiser. Mathématiquement, cela se traduit par la somme $f(s) = \sum_{e \in s} w(e)$, où $w(e)$ représente le coût individuel (ou le poids) de l'élément $e$.	
	
	\textbf{Contraintes.} Dans tout problème d'optimisation, il existe des règles strictes à respecter (comme un budget à ne pas dépasser ou l'obligation de visiter chaque ville). Ces contraintes agissent comme un filtre : elles éliminent les combinaisons impossibles pour ne conserver que les solutions « valides » (ou réalisables). L'ensemble de ces solutions valides forme ce que l'on appelle la région admissible $\mathcal{F}$. Mathématiquement, cette région est définie par un système d'inégalités (des limites à ne pas franchir, notées $g$) et d'égalités (des conditions strictes, notées $h$) :
	$$
	\mathcal{F} = \left\{ s \in \mathcal{S} \;\middle|\; g_i(s) \leq 0, \; i = 1, \ldots, p \quad \text{et} \quad h_j(s) = 0, \; j = 1, \ldots, q \right\}
	$$
	\begin{figure}[H]
		\centering
		\begin{tikzpicture}[scale=0.5, font=\small]
			
			% 1. Définition du contour de la zone admissible
			\def\regionpath{(0,0) -- (7,0) -- (9,5) -- (0,2.5) -- cycle}
			
			% 2. Remplissage de la région avec ta couleur ThGreen (éclaircie pour voir les points)
			\fill[ThGreen!40] \regionpath;
			
			% 3. Dessin des axes
			\draw[-latex, thick] (-1,0) -- (11,0) node[right] {$x$};
			\draw[-latex, thick] (0,-1) -- (0,6.5) node[above] {$y$};
			
			% 4. Lignes des contraintes
			% Contrainte 1 (pente douce)
			\draw[thick] (-1, 2.22) -- (11, 5.55) node[pos=0.45, above, sloped] {Contrainte 1};
			
			% Contrainte 2 (pente raide)
			\draw[thick] (6.6, -1) -- (9.6, 6.5) node[pos=0.6, below, sloped] {Contrainte 2};
			
			% 5. Conditions de non-négativité
			\node[rotate=90] at (-0.4, 1.2) {$x > 0$};
			\node at (1.5, -0.4) {$y > 0$};
			
			% 6. Génération propre des points entiers (la fonction clip garantit qu'ils restent dedans)
			\begin{scope}
				\clip \regionpath;
				\foreach \x in {1,2,...,9} {
					\foreach \y in {1,2,...,5} {
						\fill[black] (\x,\y) circle (2pt);
					}
				}
			\end{scope}
			
			% 7. Texte explicatif
			\draw[latex-, thick] (7.5, 2.5) -- (9.5, 1.5) node[right, align=left] {Points entiers dans \\ la région admissible};
			
		\end{tikzpicture}
		\caption{Représentation visuelle d'un problème d'optimisation en nombres entiers. Contrairement à l'optimisation continue, la région admissible $\mathcal{F}$ est restreinte aux points discrets qui respectent toutes les contraintes.}
		\label{fig:points_entiers_tikz}
	\end{figure}
	
	
	% ─────────────────────────────────────────────────────────────
	\section{Notions de complexité}
	\label{sec:notions_complexite}
	% ─────────────────────────────────────────────────────────────
	
	La théorie de la complexité computationnelle fournit un cadre formel
	permettant de classifier les problèmes selon les ressources algorithmiques
	nécessaires à leur résolution.
	
	\textbf{La classe $\mathcal{P}$.}
	La classe $\mathcal{P}$ comprend tous les problèmes de décision qui peuvent
	être résolus par un algorithme déterministe en temps polynomial par rapport
	à la taille de l'entrée. En d'autres termes, un problème appartient à
	$\mathcal{P}$ s'il existe un algorithme qui, pour toute entrée de taille
	$n$, produit une réponse en un temps de la forme $\mathcal{O}(n^k)$, où
	$k$ est une constante fixée :
	$$
	\mathcal{P} = \left\{ L \;\middle|\; \exists \text{ algorithme
		déterministe résolvant } L \text{ en }
	\mathcal{O}(n^k), \; k \in \N \right\}
	$$
	Les problèmes dans $\mathcal{P}$ sont considérés comme \textit{faciles}
	au sens calculatoire : leur temps de résolution reste raisonnable même
	lorsque la taille de l'entrée augmente. Le tri d'une liste, la recherche
	dans un graphe, et la détermination de la connectivité d'un graphe en
	sont des exemples classiques.
	
	\textbf{La classe $\mathcal{NP}$.}
	La classe $\mathcal{NP}$ (Non-déterministe Polynomiale) contient tous les
	problèmes de décision pour lesquels une solution candidate, si elle existe,
	peut être \textit{vérifiée} en temps polynomial par un algorithme
	déterministe. Cela ne signifie pas que le problème peut nécessairement
	être \textit{résolu} en temps polynomial, mais que si l'on dispose d'un
	\textbf{certificat} (une preuve de solution), sa validité peut être
	confirmée efficacement :
	$$
	\mathcal{NP} = \left\{ L \;\middle|\; \exists \text{ vérifieur
		polynomial } V \text{ tel que }
	x \in L \Leftrightarrow \exists\, c,\; V(x, c) = 1 \right\}
	$$
	
	\begin{example}{Le problème SAT}{exemple_sat}
		Le problème de la satisfiabilité booléenne (SAT) est un exemple
		canonique de problème $\mathcal{NP}$. Étant donnée une expression
		booléenne, il s'agit de déterminer s'il existe une affectation des
		variables qui la rend vraie. Trouver cette affectation est difficile,
		mais si on nous en propose une, la vérifier se fait en temps linéaire.
		C'est précisément la dissymétrie entre \textit{résoudre} et
		\textit{vérifier} qui définit $\mathcal{NP}$.
	\end{example}
	
	Il est important de noter que $\mathcal{P} \subseteq \mathcal{NP}$ :
	tout problème soluble en temps polynomial peut également avoir sa solution
	vérifiée en temps polynomial. La question fondamentale --- encore ouverte
	à ce jour --- est de savoir si cette inclusion est stricte, c'est-à-dire
	si $\mathcal{P} \neq \mathcal{NP}$.
	
	\textbf{La $\mathcal{NP}$-complétude.}
	La classe $\mathcal{NP}$-complet est un sous-ensemble de $\mathcal{NP}$
	regroupant les problèmes les plus \textit{difficiles} de cette classe.
	Un problème $\mathcal{NP}$-complet possède la propriété remarquable
	suivante : si l'on trouvait un algorithme polynomial pour en résoudre un
	seul, alors \textit{tous} les problèmes de $\mathcal{NP}$ pourraient être
	résolus en temps polynomial.
	
	La démonstration qu'un problème $A$ est $\mathcal{NP}$-complet repose sur
	deux conditions à établir simultanément :
	\begin{enumerate}
		\item $A \in \mathcal{NP}$ : une solution candidate est vérifiable en
		temps polynomial.
		\item Tout problème de $\mathcal{NP}$ se réduit polynomialement à $A$ :
		résoudre $A$ suffit à résoudre n'importe quel problème de la classe.
	\end{enumerate}
	
	\begin{remark}{Historique de la NP-complétude}{historique_npc}
		Le premier problème démontré $\mathcal{NP}$-complet est SAT (théorème
		de Cook-Levin, 1971~\cite{cook1971complexity}). Depuis, une vaste famille de problèmes a été
		prouvée $\mathcal{NP}$-complète par réduction en chaîne : la coloration
		de graphe, le problème du sac à dos, et --- objet central de ce mémoire
		--- la version de décision du Problème du Voyageur de Commerce.
	\end{remark}
	
	\textbf{La $\mathcal{NP}$-difficulté.}
	Un problème est dit $\mathcal{NP}$-difficile s'il est au moins aussi
	difficile que tous les problèmes $\mathcal{NP}$-complets, sans
	nécessairement appartenir à $\mathcal{NP}$ lui-même. C'est le cas de la
	\textit{version optimisation} du TSP : trouver la tournée de coût minimum
	n'est pas un problème de décision (il n'y a pas de réponse oui/non), mais
	il est au moins aussi difficile que sa version de décision, qui elle est
	$\mathcal{NP}$-complète.
	
	La Figure~\ref{fig:hierarchie_complexite} illustre cette hiérarchie.
	
	$$
	\begin{aligned}
	\mathcal{P} &\subseteq \mathcal{NP},\\
	\mathcal{NP}\text{-complet}
	&= \mathcal{NP}\cap\mathcal{NP}\text{-difficile}.
	\end{aligned}
	$$
	
	% ── Figure hiérarchie ──────────────────────────────────────
	\begin{figure}[H]
		\centering
		\begin{tikzpicture}[scale=0.6, every node/.style={font=\sffamily}]
			
			% Définition des couleurs (inspirées de ton image, douces et académiques)
			\definecolor{colorP}{RGB}{180, 255, 180}
			\definecolor{colorNP}{RGB}{160, 240, 160}
			\definecolor{colorNPC}{RGB}{110, 200, 130}
			\definecolor{colorNPH}{RGB}{150, 190, 180}
			
			% 1. Grande bulle NP (En bas)
			\filldraw[fill=colorNP, draw=black!40, thick] (0, -1.5) ellipse (4.5cm and 3.5cm);
			
			% 2. Zone d'intersection (NP-Complet) calculée par clipping
			\begin{scope}
				\clip (0, -1.5) ellipse (4.5cm and 3.5cm);
				\filldraw[fill=colorNPC, draw=black!40, thick] (0, 2) ellipse (4.5cm and 3cm);
			\end{scope}
			
			% 3. Grande bulle NP-Hard (En haut, avec opacité pour laisser voir l'intersection)
			\draw[black!40, thick] (0, 2) ellipse (4.5cm and 3cm);
			% On remplit uniquement la partie hors NP (visuellement)
			\fill[colorNPH, fill opacity=0.7] (0, 2) ellipse (4.5cm and 3cm);
			
			% 4. Petite bulle P (À l'intérieur de NP)
			\filldraw[fill=colorP, draw=black!40, thick] (0, -3.2) ellipse (2cm and 1.2cm);
			
			% 5. Textes et Exemples dans les bulles
			\node[align=center] at (0, -3.2) {\textbf{\Large P}\\ \small Tri, pgcd};
			\node[align=center] at (-2.5, -1.5) {\textbf{\Large NP}\\ \small Factorisation};
			\node[align=center] at (0, 0.4) {\textbf{\Large NP-Complet}\\ \small TSP (Décision), SAT};
			\node[align=center] at (0, 3.2) {\textbf{\Large NP-Difficile}\\ \small \textbf{TSP (Optimisation)}};
			
			% 6. Flèche de difficulté (Dégradé vertical)
			\shade[top color=green!70!black, bottom color=colorP] (6, -4) rectangle (6.4, 3.5);
			\fill[green!70!black] (5.7, 3.5) -- (6.7, 3.5) -- (6.2, 4.5) -- cycle; % Pointe de la flèche
			
			% 7. Textes de l'axe de difficulté
			\node[right] at (6.8, 4) {\textbf{Le plus difficile}};
			\node[right] at (6.8, 1) {\textbf{Difficile}};
			\node[right] at (6.8, -1.5) {\textbf{Moyen}};
			\node[right] at (6.8, -4) {\textbf{Facile}};
			
		\end{tikzpicture}
		\caption{Schéma des relations entre classes de complexité. La séparation représentée entre $\mathcal{P}$ et les problèmes $\mathcal{NP}$-complets illustre le cas hypothétique $\mathcal{P}\subsetneq\mathcal{NP}$ ; la question $\mathcal{P}=\mathcal{NP}$ demeure ouverte. La version optimisation du TSP est $\mathcal{NP}$-difficile.}
		\label{fig:hierarchie_complexite}
	\end{figure}
	
	% ──────────────────────────────────────────────────────────
	
	\textbf{Application au TSP.}
	La version de décision du TSP est $\mathcal{NP}$-complète, et sa version
	d'optimisation est $\mathcal{NP}$-difficile.	
	% ──────────────────────────────────────────────────────────
	L'algorithme exact de référence (Bellman-Held-Karp) résout le TSP
	en $\mathcal{O}(n^2 2^n)$ en temps et $\mathcal{O}(n\,2^n)$ en espace,
	illustrant la barrière exponentielle caractéristique des problèmes
	$\mathcal{NP}$-difficiles. Une analyse quantitative de cette barrière
	est conduite au Chapitre~\ref{chap:tsp}, Section~\ref{subsec:tsp_np}.
	% ──────────────────────────────────────────────────────────
	
	
	
	
	% ─────────────────────────────────────────────────────────────
	\section{Méthodes de résolution}
	\label{sec:methodes_res}
	% ─────────────────────────────────────────────────────────────
	
	Face à la difficulté intrinsèque des problèmes combinatoires, trois
	familles de méthodes ont été développées, offrant des compromis distincts
	entre garantie d'optimalité et temps de calcul. La
	Figure~\ref{fig:taxonomie_methodes} présente leur organisation.
	
	% ── Figure taxonomie ───────────────────────────────────────
	% ─────────────────────────────────────────────────────────────
	% FIGURE 2 : Taxonomie des méthodes de résolution
	% ─────────────────────────────────────────────────────────────
	\begin{figure}[H]
		\centering
		\resizebox{\textwidth}{!}{%
			\begin{tikzpicture}[
				node distance = 0.6cm and 1.2cm,
				box/.style  = {draw=ThBlue, fill=ThBlue!8, rounded corners=4pt,
					minimum width=3.2cm, minimum height=0.85cm,
					align=center, font=\small\bfseries, text=ThBlue},
				leaf/.style = {draw=ThBlue!60, fill=ThBlue!4, rounded corners=4pt,
					minimum width=2.9cm, minimum height=0.75cm,
					align=center, font=\small, text=ThBlue!80!black},
				meta/.style = {draw=ThRed!60, fill=ThRed!4, rounded corners=4pt,
					minimum width=2.9cm, minimum height=0.75cm,
					align=center, font=\small, text=ThRed!80!black},
				arr/.style  = {-Stealth, thick, color=ThBlue!70}
				]
				
				% Racine
				\node[box, minimum width=5.0cm] (root)
				{Méthodes de résolution};
				
				% Niveau 1 — trois branches
				\node[box, below left=1.2cm and 3.8cm of root]  (exact)
				{Méthodes exactes};
				\node[box, below=1.2cm of root]                 (heur)
				{Heuristiques};
				\node[box, below right=1.2cm and 3.8cm of root]  (meta)
				{Métaheuristiques};
				
				% Feuilles exactes
				\node[leaf, below=0.6cm of exact, xshift=-1.7cm] (bhk)
				{Bellman-\\Held-Karp};
				\node[leaf, below=0.6cm of exact, xshift= 1.7cm] (bnb)
				{Branch-and-Cut\\(Concorde)};
				
				% Feuilles heuristiques
				\node[leaf, below=0.6cm of heur, xshift=-1.6cm] (glouton)
				{Algorithmes\\gloutons};
				\node[leaf, below=0.6cm of heur, xshift= 1.6cm] (local)
				{Recherche\\locale};
				
				% Feuilles métaheuristiques
				\node[meta, below=0.6cm of meta, xshift=-1.7cm] (sa)
				{Recuit\\simulé};
				\node[meta, below=0.6cm of meta, xshift= 1.7cm] (ga)
				{Algorithmes\\génétiques};
				
				% Arcs
				\draw[arr] (root) -- (exact);
				\draw[arr] (root) -- (heur);
				\draw[arr] (root) -- (meta);
				\draw[arr] (exact) -- (bhk);
				\draw[arr] (exact) -- (bnb);
				\draw[arr] (heur)  -- (glouton);
				\draw[arr] (heur)  -- (local);
				\draw[arr] (meta)  -- (sa);
				\draw[arr] (meta)  -- (ga);
				
				% Complexités
				\node[font=\tiny, color=ThRed, below=0.05cm of bhk]
				{$\mathcal{O}(n^2 2^n)$};
				\node[font=\tiny, color=ThRed, below=0.05cm of bnb]
				{exponentiel};
				\node[font=\tiny, color=ThGreen!60!black, below=0.05cm of glouton]
				{$\mathcal{O}(n^2)$};
				\node[font=\tiny, color=ThGreen!60!black, below=0.05cm of local]
				{$\mathcal{O}(n^2)$/iter};
				\node[font=\tiny, color=ThGreen!60!black, below=0.05cm of sa]
				{paramétrique};
				\node[font=\tiny, color=ThGreen!60!black, below=0.05cm of ga]
				{paramétrique};
				
			\end{tikzpicture}%
		}
		\caption{Taxonomie des méthodes de résolution des problèmes d'optimisation
			combinatoire. Les heuristiques produisent une solution en temps
			polynomial sans garantie d'optimalité. Les métaheuristiques
			(en rouge) ajoutent des mécanismes d'échappement aux optima locaux.}
		\label{fig:taxonomie_methodes}
	\end{figure}
	% ──────────────────────────────────────────────────────────
	
	\subsection{Méthodes exactes}
	\label{subsec:methodes_exactes}
	
	Les méthodes exactes garantissent la solution optimale globale, au prix
	d'une complexité exponentielle dans le pire des cas.
	
	\textbf{Programmation linéaire en nombres entiers (PLNE).}
	Un problème de PLNE s'écrit :
	$$
	\min_{x \in \Z^n} c^\top x \quad \text{sous} \quad
	Ax \leq b, \; x \geq 0
	$$
	La relaxation continue (LP) fournit une borne inférieure sur la valeur
	optimale, constituant la base des méthodes de décomposition.
	
	\textbf{Méthode de séparation et évaluation (Branch-and-Cut).}
	Cette méthode explore implicitement l'espace des solutions en le
	partitionnant récursivement (séparation) et en calculant des bornes
	(évaluation) pour élaguer les branches non prometteuses.
	
	\subsection{Heuristiques et métaheuristiques}
	\label{subsec:heuristiques_meta}
	
	Les heuristiques proposent un compromis pragmatique : des solutions de
	haute qualité en temps maîtrisé, sans garantie d'optimalité.
	
	\textbf{Algorithmes gloutons.}
	Ils construisent une solution par des choix locaux optimaux séquentiels.
	Rapides en $\mathcal{O}(n^2)$, ils demeurent myopes et incapables de
	corriger de mauvais choix initiaux.
	
	\textbf{Recherche locale.}
	Partant d'une solution $s_0$, la méthode itère vers un voisin améliorant
	jusqu'à convergence vers un optimum local $s^*$ :
	$$
	f(s^*) \leq f(s), \quad \forall s \in \mathcal{N}(s^*)
	$$
	
	
	\textbf{Recuit simulé.}
	Cette métaheuristique accepte des mouvements dégradants avec la
	probabilité de Metropolis :
	$$
	P(\text{accepter } s') = \begin{cases}
		1 & \text{si } \Delta f < 0 \\
		\exp\!\left(-\dfrac{\Delta f}{T}\right) & \text{sinon}
	\end{cases}
	$$
	où $T > 0$ est une température décroissante au fil des itérations.
	
	\textbf{Algorithmes génétiques.}
	Métaheuristiques populationnelles imitant la sélection naturelle :
	$$
	\mathcal{P}_{t+1} =
	\text{Mutation}\Big(\text{Croisement}
	\big(\text{Sélection}(\mathcal{P}_t)\big)\Big)
	$$
	\\
	\\
	\textbf{Bilan et transition.}
	En effet, ces approches classiques se heurtent à deux limites fondamentales : elles requièrent une conception manuelle experte et abordent chaque nouvelle instance sans jamais tirer parti de l'expérience acquise. La réponse à ces obstacles se trouve dans le domaine de l'\textit{Apprentissage par Renforcement Profond} (DRL). Ce dernier vise à apprendre une politique de décision paramétrée $\pi_\theta$ sur un vaste corpus d'entraînement, afin de l'appliquer de manière quasi instantanée lors de la phase d'inférence.
	
	Ces trois familles de méthodes constituent le socle algorithmique de notre mémoire. Présentées ici sous leur forme générique, elles seront appliquées au Problème du Voyageur de Commerce dans le Chapitre~\ref{chap:tsp}. Nous y détaillerons leurs formulations spécifiques, leurs performances empiriques, ainsi que leurs limites respectives dans ce contexte précis.

	
	
	% ============================================================
	%  CHAPITRE II : LE PROBLEME DU VOYAGEUR DE COMMERCE
	% ============================================================
	\chapter{Le probl\`eme du voyageur de commerce}
	\label{chap:tsp}
	
Le Problème du Voyageur de Commerce (TSP, de l'anglais \textit{Traveling Salesman Problem}) est probablement le problème ayant reçu le plus d'attention dans toute l'histoire de l'optimisation combinatoire. Son énoncé est d'une grande simplicité : étant donné un ensemble de villes et les distances qui les séparent, l'objectif est de trouver le plus court itinéraire passant par chaque ville exactement une fois avant de revenir à son point de départ.
	
	% ─────────────────────────────────────────────────────────────
	\section{D\'efinition et mod\'elisation}
	\label{sec:tsp_definition}
	% ─────────────────────────────────────────────────────────────
	
	\subsection{Cadre graph-th\'eorique}
	\label{subsec:cadre_graphe}
	
	Le TSP s'inscrit naturellement dans le cadre de la théorie des graphes.
	Soit $G = (V, E, c)$ un graphe complet pondéré, où :
	\begin{itemize}
		\item $V = \{v_1, v_2, \ldots, v_n\}$ est l'ensemble des $n$ sommets
		(les villes),
		\item $E = \{\{i,j\} \mid i,j \in V,\; i \neq j\}$ est l'ensemble des
		arêtes (les liaisons entre villes),
		\item $c : E \to \R^+$ est une fonction de coût associant à chaque arête
		$\{i,j\}$ un poids $c_{ij} \geq 0$.
	\end{itemize}
	
	% ── Macro interne : dessine K_n (définie à l'extérieur pour servir partout) ──
	\newcommand{\grapheComplet}[2]{%
		% #1 = n (nombre de sommets)
		% #2 = angle de départ
		\begin{tikzpicture}[scale=0.6]
			% Arêtes (dessinées en premier pour passer sous les nœuds)
			\foreach \i in {1,...,#1} {
				\foreach \j in {1,...,#1} {
					\ifnum\i<\j
					\draw[ThBlue!40, line width=0.5pt]
					({#2+(\i-1)*360/#1}:1.1cm) --
					({#2+(\j-1)*360/#1}:1.1cm);
					\fi
				}
			}
			% Sommets
			\foreach \i in {1,...,#1} {
				\node[circle, fill=ThBlue, draw=white,
				line width=0.8pt, inner sep=1.5pt,
				label={[font=\tiny, color=ThBlue, inner sep=1pt]
					{#2+(\i-1)*360/#1}:$v_{\i}$}]
				at ({#2+(\i-1)*360/#1}:1.1cm) {};
			}
		\end{tikzpicture}%
	}
	
	% ──────────────────────────────────────────────────────────────
	% PREMIÈRE FIGURE : n=3 et n=4
	% ──────────────────────────────────────────────────────────────
	\begin{figure}[H]
		\centering
		\begin{subfigure}[b]{0.45\textwidth}
			\centering
			\grapheComplet{3}{90}
			\vspace{2pt}
			\caption*{\footnotesize $n = 3$\\[2pt]
				$|E| = \frac{3 \times 2}{2} = 3$}
		\end{subfigure}
		\hfill
		\begin{subfigure}[b]{0.45\textwidth}
			\centering
			\grapheComplet{4}{45}
			\vspace{2pt}
			\caption*{\footnotesize $n = 4$\\[2pt]
				$|E| = \frac{4 \times 3}{2} = 6$}
		\end{subfigure}
		\vspace{2pt}
		\caption{Graphes complets $K_n$ pour $n \in \{3, 4\}$. Chaque paire de sommets est reliée par exactement une arête, et le nombre total d'arêtes est donné par la formule $|E| = \frac{n(n-1)}{2}$.}
		\label{fig:graphes_complets_3_4}
	\end{figure}
	
	% Tu peux écrire du texte ici entre les deux figures si tu le souhaites !
	
	% ──────────────────────────────────────────────────────────────
	% DEUXIÈME FIGURE : n=5 et n=6
	% ──────────────────────────────────────────────────────────────
	\begin{figure}[H]
		\centering
		\begin{subfigure}[b]{0.45\textwidth}
			\centering
			\grapheComplet{5}{90}
			\vspace{2pt}
			\caption*{\footnotesize $n = 5$\\[2pt]
				$|E| = \frac{5 \times 4}{2} = 10$}
		\end{subfigure}
		\hfill
		\begin{subfigure}[b]{0.45\textwidth}
			\centering
			\grapheComplet{6}{90}
			\vspace{2pt}
			\caption*{\footnotesize $n = 6$\\[2pt]
				$|E| = \frac{6 \times 5}{2} = 15$}
		\end{subfigure}
		\vspace{2pt}
		\caption{Graphes complets $K_n$ pour $n \in \{5, 6\}$. L'ajout de sommets montre visuellement la croissance quadratique du nombre de liaisons. Pour une instance industrielle de $n = 100$ villes, cela représente déjà $4\,950$ arêtes à explorer.}
		\label{fig:graphes_complets_5_6}
	\end{figure}
	% ──────────────────────────────────────────────────────────────
	
	
	
	Un graphe complet sur $n$ sommets, noté $K_n$, contient exactement
	$|E| = \dfrac{n(n-1)}{2}$ arêtes. Comme l'illustrent les
	Figures~\ref{fig:graphes_complets_3_4} et \ref{fig:graphes_complets_5_6}, cette quantité croît rapidement :
	dès $n = 20$ villes, le graphe contient $190$ arêtes, et pour $n = 100$,
	on atteint $4\,950$ arêtes. C'est sur ce graphe complet que le TSP
	cherche le cycle hamiltonien de coût minimum.
	
	\begin{definition}{Cycle hamiltonien}{cycle_hamiltonien}
		Soit $G = (V, E)$ un graphe. Un \textbf{cycle hamiltonien} est un cycle
		qui visite chaque sommet de $V$ exactement une fois avant de revenir au
		sommet de départ. Formellement, c'est une permutation
		$\sigma \in \mathfrak{S}_n$ des sommets de $V$ définissant le circuit
		ordonné $(v_{\sigma(1)},\, v_{\sigma(2)},\, \ldots,\, v_{\sigma(n)},\,
		v_{\sigma(1)})$.
	\end{definition}
	\begin{figure}[H]
		\centering
		
		% --- SOUS-FIGURE 1 : Permutation 1 ---
		\begin{subfigure}[b]{0.48\textwidth}
			\centering
			\begin{tikzpicture}[
				scale=0.55, % Échelle légèrement réduite pour tenir côte à côte
				vertex/.style={circle, draw=ThBlue, fill=ThBlue!10, thick, minimum size=9mm, font=\small\bfseries, text=ThBlue!80!black},
				backgroundEdge/.style={thin, color=black!15},
				hamiltonianEdge/.style={-Stealth, ultra thick, color=ThGreen!80!black}, % Couleur verte pour le différencier
				weight/.style={midway, fill=white, font=\footnotesize, inner sep=1pt, text=black!70}
				]
				
				\node[vertex] (v1) at (90:3cm) {$v_1$};
				\node[vertex] (v2) at (18:3cm) {$v_2$};
				\node[vertex] (v3) at (-54:3cm) {$v_3$};
				\node[vertex] (v4) at (-126:3cm) {$v_4$};
				\node[vertex] (v5) at (162:3cm) {$v_5$};
				
				% Arêtes de fond
				\foreach \i in {1,...,4} {
					\pgfmathsetmacro{\next}{\i+1}
					\foreach \j in {\next,...,5} {
						\draw[backgroundEdge] (v\i) -- (v\j);
					}
				}
				
				% Cycle Hamiltonien 1 : (1, 2, 3, 4, 5)
				\draw[hamiltonianEdge] (v1) -- (v2) node[weight, text=ThGreen!80!black] {$c_{12}$};
				\draw[hamiltonianEdge] (v2) -- (v3) node[weight, text=ThGreen!80!black] {$c_{23}$};
				\draw[hamiltonianEdge] (v3) -- (v4) node[weight, text=ThGreen!80!black] {$c_{34}$};
				\draw[hamiltonianEdge] (v4) -- (v5) node[weight, text=ThGreen!80!black] {$c_{45}$};
				\draw[hamiltonianEdge] (v5) -- (v1) node[weight, text=ThGreen!80!black] {$c_{51}$};
				
			\end{tikzpicture}
			\caption{Tournée $\sigma_1 = (1, 2, 3, 4, 5)$}
		\end{subfigure}
		\hfill
		% --- SOUS-FIGURE 2 : Permutation 2 ---
		\begin{subfigure}[b]{0.48\textwidth}
			\centering
			\begin{tikzpicture}[
				scale=0.55, % Même échelle
				vertex/.style={circle, draw=ThBlue, fill=ThBlue!10, thick, minimum size=9mm, font=\small\bfseries, text=ThBlue!80!black},
				backgroundEdge/.style={thin, color=black!15},
				hamiltonianEdge/.style={-Stealth, ultra thick, color=ThRed}, % Couleur rouge (ton graphe initial)
				weight/.style={midway, fill=white, font=\footnotesize, inner sep=1pt, text=black!70}
				]
				
				\node[vertex] (v1) at (90:3cm) {$v_1$};
				\node[vertex] (v2) at (18:3cm) {$v_2$};
				\node[vertex] (v3) at (-54:3cm) {$v_3$};
				\node[vertex] (v4) at (-126:3cm) {$v_4$};
				\node[vertex] (v5) at (162:3cm) {$v_5$};
				
				% Arêtes de fond
				\foreach \i in {1,...,4} {
					\pgfmathsetmacro{\next}{\i+1}
					\foreach \j in {\next,...,5} {
						\draw[backgroundEdge] (v\i) -- (v\j);
					}
				}
				
				% Cycle Hamiltonien 2 : (1, 3, 4, 2, 5)
				\draw[hamiltonianEdge] (v1) -- (v3) node[weight, text=ThRed] {$c_{13}$};
				\draw[hamiltonianEdge] (v3) -- (v4) node[weight, text=ThRed] {$c_{34}$};
				\draw[hamiltonianEdge] (v4) -- (v2) node[weight, text=ThRed] {$c_{42}$};
				\draw[hamiltonianEdge] (v2) -- (v5) node[weight, text=ThRed] {$c_{25}$};
				\draw[hamiltonianEdge] (v5) -- (v1) node[weight, text=ThRed] {$c_{51}$};
				
			\end{tikzpicture}
			\caption{Tournée $\sigma_2 = (1, 3, 4, 2, 5)$}
		\end{subfigure}
		\caption{Comparaison de deux cycles hamiltoniens différents sur le même graphe $K_5$. Chaque permutation $\sigma$ génère une tournée différente avec un coût total $\mathcal{C}(\sigma)$ distinct. L'objectif de l'optimisation combinatoire est d'isoler la permutation minimisant ce coût parmi les $(n-1)!/2$ possibilités.}	
		\label{fig:comparaison_cycles_hamiltoniens}
	\end{figure}
	
	Le coût d'un cycle hamiltonien $\sigma$ est alors la somme des coûts des
	arêtes qui le composent :
	$$
	C(\sigma) = \sum_{k=1}^{n-1} c_{\sigma(k),\,\sigma(k+1)}
	+ c_{\sigma(n),\,\sigma(1)}
	= \sum_{k=1}^{n} c_{\sigma(k),\,\sigma(k \bmod n + 1)}
	$$
	
	L'objectif du TSP est de déterminer le cycle hamiltonien de coût minimum :
	$$
	\sigma^* = \argmin_{\sigma \in \mathfrak{S}_n} C(\sigma)
	$$
	
	\begin{remark}{Variantes du TSP}{variantes_tsp}
		Deux variantes principales sont classiquement distinguées :
		\begin{itemize}
			\item \textbf{TSP symétrique (sTSP) :} $c_{ij} = c_{ji}$ pour tout
			$(i,j) \in V \times V$. C'est la variante la plus courante,
			correspondant à des distances euclidiennes ou routières.
			\item \textbf{TSP asymétrique (aTSP) :} $c_{ij} \neq c_{ji}$ en
			général, modélisant des situations où le coût de déplacement
			dépend du sens de parcours (trafic, sens uniques).
		\end{itemize}
		Sauf mention explicite, ce mémoire se concentre sur le \textbf{sTSP
			euclidien}, où les villes sont des points de $\R^2$ et $c_{ij}$ est la
		distance euclidienne standard.
	\end{remark}
	
	% ─────────────────────────────────────────────────────────────
	\subsection{Formulation math\'ematique}
	\label{subsec:formulation_math}
	% ─────────────────────────────────────────────────────────────
	
	La richesse du TSP tient en partie à l'existence de plusieurs formulations
	équivalentes en Programmation Linéaire en Nombres Entiers (PLNE). Nous
	présentons ici la formulation de référence, dite \textbf{DFJ} (Dantzig,
	Fulkerson et Johnson~\cite{dantzig1954solution}), qui constitue la base des solveurs exacts modernes.
	
	\subsubsection*{Variables de d\'ecision}
	
	On associe à chaque arête $\{i,j\} \in E$ une variable binaire :
	$$
	x_{ij} = \begin{cases}
		1 & \text{si l'arc } (i,j) \text{ est emprunté dans la tournée,} \\
		0 & \text{sinon.}
	\end{cases}
	$$
	Par symétrie, on pose $x_{ij} = x_{ji}$ et on ne considère qu'une variable
	par paire non ordonnée $\{i,j\}$.
	
	\subsubsection*{Formulation DFJ}
	
	Le TSP symétrique s'écrit alors comme le programme entier suivant :
	
	\begin{equation}
		\min \sum_{\{i,j\} \in E} c_{ij}\, x_{ij}
		\label{eq:tsp_obj}
	\end{equation}
	
	sous les contraintes :
	
	\begin{align}
		\sum_{j \in V \setminus \{i\}} x_{ij} &= 2
		&& \forall\, i \in V
		\label{eq:tsp_degree} \\[6pt]
		\sum_{\substack{i \in S,\, j \notin S}} x_{ij} &\geq 2
		&& \forall\, S \subsetneq V,\; 2 \leq |S| \leq n-2
		\label{eq:tsp_sec} \\[6pt]
		x_{ij} &\in \{0, 1\}
		&& \forall\, \{i,j\} \in E
		\label{eq:tsp_bin}
	\end{align}
	
	Chaque contrainte joue un rôle précis dans la structure de la solution :
	
	\begin{itemize}
		\item \textbf{Contraintes de degré} \eqref{eq:tsp_degree} : chaque ville
		est reliée à exactement deux arêtes de la tournée, garantissant
		qu'elle est visitée une fois en entrée et une fois en sortie.
		
		\item \textbf{Contraintes d'élimination de sous-tours (SEC)} 
		\eqref{eq:tsp_sec} : toute coupe séparant un sous-ensemble propre
		$S$ de villes du reste du graphe doit être traversée par au moins
		deux arêtes. Ces contraintes empêchent la formation de cycles
		disconnectés ne couvrant pas l'ensemble $V$.
		
		\item \textbf{Contraintes d'intégralité} \eqref{eq:tsp_bin} : les
		variables de décision sont binaires.
	\end{itemize}
	
	\begin{remark}{Explosion du nombre de contraintes SEC}{explosion_sec}
Le nombre de contraintes d'élimination de sous-tours est de l'ordre de $2^n$, ce qui rend impossible leur énumération exhaustive dès que $n$ prend des valeurs importantes. En pratique, des solveurs exacts tels que Concorde les génèrent dynamiquement via la technique des plans coupants (\textit{cutting planes}), en n'ajoutant à chaque étape que les contraintes violées par la solution courante. C'est précisément le c\oe{}ur de la méthode de \textit{Branch-and-Cut} abordée au Chapitre~\ref{chap:fondements}.
	\end{remark}
	
	
	
	% ─────────────────────────────────────────────────────────────
	\section{Complexit\'e et propri\'et\'es}
	\label{sec:tsp_complexite}
	% ─────────────────────────────────────────────────────────────
	
	\subsection{Taille de l'espace de recherche}
	\label{subsec:tsp_espace}
	
	La difficulté du TSP s'apprécie d'abord par la taille de son espace de
	recherche. Pour un graphe complet de $n$ sommets, le nombre total de cycles
	hamiltoniens distincts est :
	$$
	|\mathfrak{S}_n^{\text{TSP}}| = \frac{(n-1)!}{2}
	$$
	Le facteur $(n-1)!$ provient du fait que l'on fixe une ville de départ
	(le cycle est invariant par rotation), et la division par $2$ élimine les
	doublons dus à la symétrie (parcourir un cycle dans un sens ou dans l'autre
	donne la même tournée). Cette croissance factorielle rend toute énumération
	exhaustive computationnellement prohibitive dès que $n$ dépasse quelques
	dizaines, comme l'illustre le tableau suivant.
	
	\begin{table}[htbp	]
		\centering
		\caption{Croissance de l'espace de recherche du TSP en fonction de $n$.}
		\label{tab:tsp_espace}
		\begin{tabular}{rr}
			\toprule
			\textbf{Nombre de villes} $n$ & \textbf{Nombre de tournées}
			$\dfrac{(n-1)!}{2}$ \\
			\midrule
			5   & $12$ \\
			10  & $181\,440$ \\
			20  & $\approx 6{,}08 \times 10^{16}$ \\
			50  & $\approx 1{,}52 \times 10^{62}$ \\
			100 & $\approx 4{,}67 \times 10^{156}$ \\
			\bottomrule
		\end{tabular}
	\end{table}	
	\subsection{Statut de complexit\'e}
	\label{subsec:tsp_np}	
	En s'appuyant sur les notions de complexité introduites au
	Chapitre~\ref{chap:fondements}, on peut établir précisément le statut
	computationnel du TSP.	
	\begin{theorem}{NP-difficulté du TSP}{tsp_np_dur}
		La version d'optimisation du TSP est $\mathcal{NP}$-difficile. Sa
		version de décision --- \og Existe-t-il un cycle hamiltonien de coût
		inférieur à un seuil $K$ donné ? \fg{} --- est $\mathcal{NP}$-complète.
	\end{theorem}	
	La $\mathcal{NP}$-complétude de la version de décision se démontre par
	réduction polynomiale depuis le problème du cycle hamiltonien, lui-même
	$\mathcal{NP}$-complet. La vérification d'un certificat (une tournée
	candidate) se fait en temps $\mathcal{O}(n)$, plaçant bien le problème dans
	$\mathcal{NP}$.
	
	L'algorithme exact de référence, la programmation dynamique de
	Bellman-Held-Karp, résout le TSP en $\mathcal{O}(n^2 2^n)$ en temps et
	$\mathcal{O}(n\, 2^n)$ en espace. Bien que cette complexité soit
	exponentiellement meilleure que la force brute en $\mathcal{O}(n!)$, elle
	demeure intraitable pour des instances industrielles ($n \geq 50$).
	
	Pour quantifier concrètement cette barrière, considérons un processeur
	moderne capable d'exécuter $10^9$ opérations par seconde.
	Le Tableau~\ref{tab:bhk_temps} confronte les temps de calcul théoriques
	de Bellman-Held-Karp à la réalité physique des ressources disponibles.
	
	\begin{table}[H]
		\centering
		\caption{Temps de calcul théorique de Bellman-Held-Karp
			($\mathcal{O}(n^2 2^n)$) et mémoire requise
			($\mathcal{O}(n\, 2^n)$) en fonction de $n$.}
		\label{tab:bhk_temps}
		\small
		\begin{tabular}{rrrr}
			\toprule
			$n$ & Opérations & Temps estimé & Mémoire requise \\
			\midrule
			10  & $\approx 10^5$           & $< 1$ ms                    & $\sim$ 80 Ko \\
			20  & $\approx 4\times10^8$    & $\sim$ 0{,}4 s              & $\sim$ 160 Mo \\
			30  & $\approx 10^{11}$        & $\sim$ 97 s                 & $\sim$ 256 Go \\
			50  & $\approx 3\times10^{18}$ & $\sim$ 90 ans               & hors d'atteinte \\
			100 & $\approx 10^{34}$        & $\gg$ âge de l'univers      & hors d'atteinte \\
			\bottomrule
		\end{tabular}
	\end{table}
	
	\begin{figure}[H]
		\centering
		\includegraphics[width=\textwidth]{images/complexite_comparative.pdf}
		\caption{Comparaison des complexités algorithmiques en échelle
			logarithmique. Bellman-Held-Karp $\mathcal{O}(n^2 2^n)$
			est exponentiellement plus efficace que la force brute
			$\mathcal{O}(n!)$, mais les deux restent exponentiels et
			deviennent inéxécutables dès $n \geq 30$ sur une machine standard.}
		\label{fig:complexite_comparative}
	\end{figure}
	
	Deux obstacles distincts apparaissent. D'une part, la \textbf{barrière
		temporelle} : dès $n = 50$, le temps de calcul dépasse toute échelle
	pratique, indépendamment de la puissance de calcul disponible. D'autre
	part, la \textbf{barrière mémorielle} : BHK requiert le stockage d'un
	tableau de $n \times 2^n$ entrées. Pour $n = 30$, cela représente
	environ 256 Go de RAM --- hors de portée d'une machine standard --- avant
	même que la contrainte temporelle ne devienne déterminante.
	
	C'est cette double barrière, temporelle et mémorielle, qui justifie
	structurellement le recours aux méthodes approchées présentées en
	Section~\ref{sec:tsp_methodes}, et, au-delà, à l'apprentissage par
	renforcement profond étudié aux Chapitres~\ref{chap:apprentissage_renf}
	et~\ref{chap:impl_resultats}.
	
	
	\subsection{Inapproximabilit\'e et cas particuliers}
	\label{subsec:tsp_approx}
	
	La $\mathcal{NP}$-difficulté du TSP général s'accompagne d'un résultat
	d'inapproximabilité fort.
	
	\begin{theorem}{Inapproximabilité du TSP général}{tsp_inapprox}
		Pour tout $\rho \geq 1$, il n'existe pas d'algorithme polynomial
		garantissant une solution de coût inférieur à $\rho \cdot C(\sigma^*)$
		pour le TSP général (sans hypothèse sur les coûts), sauf si
		$\mathcal{P} = \mathcal{NP}$.
	\end{theorem}
	
	Ce résultat négatif justifie l'importance de la restriction au
	\textbf{TSP métrique}, dans lequel la fonction de coût satisfait
	l'inégalité triangulaire :
	$$
	c_{ij} \leq c_{ik} + c_{kj}, \quad \forall\, i, j, k \in V
	$$
	Cette hypothèse, vérifiée par les distances euclidiennes et routières, ouvre
	la voie aux algorithmes d'approximation. En particulier, l'algorithme de
	Christofides-Serdyukov~\cite{christofides1976worst} garantit une solution de coût au plus
	$\frac{3}{2} \cdot C(\sigma^*)$ en temps polynomial pour le TSP métrique,
	résultat qui demeure à ce jour la meilleure garantie théorique connue pour
	cette classe.
	
	\begin{remark}{Pertinence pour le DRL}{pertinence_drl}
		L'absence de garantie d'approximation polynomiale pour le TSP général,
		et la borne $\frac{3}{2}$ pour le cas métrique, constituent les
		références théoriques incontournables pour évaluer les approches
		d'apprentissage par renforcement profond. Un agent DRL ne disposant
		d'aucune garantie théorique d'optimalité, son évaluation empirique ---
		via un \textbf{écart relatif à une solution de référence} (ou à la meilleure
		méthode observée lorsqu'aucun optimum exact n'est connu) --- prend tout son
		sens dans ce contexte. Dans le chapitre expérimental, faute d'optimum exact disponible pour toutes les tailles testées, nous rapportons un écart relatif à la meilleure méthode observée sur chaque instance ; cette mesure comparative ne doit pas être confondue avec un écart à l'optimum global.
	\end{remark}
	
	% ─────────────────────────────────────────────────────────────
	\section{M\'ethodes classiques de r\'esolution}
	\label{sec:tsp_methodes}
	% ─────────────────────────────────────────────────────────────
Les trois classes d'algorithmes présentées dans le Chapitre~\ref{chap:fondements} --- les algorithmes exacts, les algorithmes heuristiques et les algorithmes métaheuristiques --- s'appliquent naturellement au TSP. Cependant, lorsque l'on considère comment ces algorithmes s'appliquent au problème en question, il devient clair que chacune de ces classes présente des caractéristiques spécifiques en fonction de la structure combinatoire du graphe et des contraintes hamiltoniennes.

Face à la $\mathcal{NP}$-difficulté du TSP, ces trois classes d'approches se sont développées pour offrir différentes stratégies de résolution.
	% ─────────────────────────────────────────────────────────────
	\subsection{M\'ethodes exactes}
	\label{subsec:tsp_exactes}
	% ─────────────────────────────────────────────────────────────
	
	Les méthodes exactes garantissent de trouver la solution optimale globale.
	Pour le TSP, elles reposent toutes sur la formulation PLNE présentée à la
	Section~\ref{subsec:formulation_math}.
	
	\textbf{Programmation dynamique de Bellman-Held-Karp.}
	C'est l'algorithme exact de référence. Soit $d(S, i)$ le coût minimal d'un
	chemin partant de la ville $1$, visitant exactement les villes de
	$S \subseteq V$ et se terminant en $i \in S$. La récurrence est :
	$$
	d(S, i) = \min_{j \in S \setminus \{i\}} \left[ d(S \setminus \{i\},\, j)
	+ c_{ji} \right]
	$$
	avec la condition initiale $d(\{1\}, 1) = 0$. Le coût optimal de la tournée
	est alors :
	$$
	C(\sigma^*) = \min_{i \in V \setminus \{1\}}
	\left[ d(V, i) + c_{i1} \right]
	$$
	Cet algorithme résout le TSP en $\mathcal{O}(n^2 2^n)$ en temps et
	$\mathcal{O}(n\, 2^n)$ en espace, ce qui le rend intraitable pour
	$n \geq 50$ dans un contexte industriel.
	
	\textbf{Branch-and-Cut.}
	Les solveurs exacts modernes, dont Concorde~\cite{applegate2006traveling} constitue l'implémentation de
	référence, reposent sur une combinaison de Branch-and-Bound et de génération
	dynamique de contraintes SEC (voir Remarque~\ref{rem:explosion_sec}). À
	chaque n\oe{}ud de l'arbre de recherche, la relaxation linéaire est résolue.
	Si la solution obtenue contient des sous-tours, les contraintes SEC
	correspondantes sont ajoutées (\textit{cutting plane}), resserrant la borne
	inférieure avant toute branche. Ce mécanisme permet à Concorde de résoudre
	des instances du TSP à plusieurs milliers de villes, mais son temps de calcul
	reste non borné polynomialement dans le pire des cas.
	
	\begin{remark}{Limites pratiques des m\'ethodes exactes}{limites_exactes}
		En pratique, les méthodes exactes deviennent inopérantes au-delà de
		quelques centaines à quelques milliers de villes selon la structure de
		l'instance. Les applications industrielles (logistique de grande échelle,
		ordonnancement) traitent couramment des instances à $10^4$ à $10^6$
		villes, rendant le recours aux méthodes approchées inévitable.
	\end{remark}
	
	% ─────────────────────────────────────────────────────────────
\subsection{Heuristiques}
\label{subsec:tsp_heuristiques}
% ─────────────────────────────────────────────────────────────

Les heuristiques construisent ou améliorent une solution en temps polynomial,
sans garantie d'optimalité globale.

\textbf{Heuristique du plus proche voisin (Nearest Neighbor).}
C'est l'heuristique constructive la plus simple. À partir d'une ville de
départ $v_1$, l'algorithme sélectionne à chaque étape la ville non encore
visitée la plus proche de la ville courante.
La complexité de cet algorithme est $\mathcal{O}(n^2)$. En pratique,
il produit des tournées dont le coût excède l'optimal de 20 à 25\% en
moyenne sur des instances euclidiennes aléatoires. Sa simplicité en fait
la baseline constructive naturelle pour toute comparaison algorithmique.

\textbf{Heuristique 2-opt.}
La méthode 2-opt est une heuristique d'amélioration par recherche locale.
Partant d'une tournée initiale $\sigma$, elle explore le voisinage $\mathcal{N}_{2\text{-opt}}(\sigma)$ défini par tous les échanges possibles de deux arêtes :
une paire d'arêtes $\{i, i+1\}$ et $\{j, j+1\}$ est supprimée et remplacée
par $\{i, j\}$ et $\{i+1, j+1\}$, inversant le segment intermédiaire.

Chaque itération a une complexité en $\mathcal{O}(n^2)$. L'algorithme
converge vers un optimum local au sens 2-opt, dont la qualité dépend
fortement de la solution initiale. En pratique, 2-opt appliqué après
Nearest Neighbor réduit le coût de la tournée de 5 à 10\% supplémentaires.

\begin{algorithm}[htbp]
	\caption{Plus Proche Voisin (Nearest Neighbor)}
	\label{alg:nn}
	\KwEntree{Graphe complet pondéré $G = (V, E, c)$, ville de départ $v_1$}
	\KwSortie{Tournée $\sigma$}
	$\text{visite} \leftarrow \{v_1\}$ \;
	$v_{\text{courant}} \leftarrow v_1$ \;
	\While{$|\text{visite}| < n$}{
		$v_{\text{next}} \leftarrow \argmin_{v \in V \setminus \text{visite}}\,
		c_{v_{\text{courant}},\, v}$ \;
		Ajouter $v_{\text{next}}$ à $\sigma$ \;
		$v_{\text{courant}} \leftarrow v_{\text{next}}$ \;
		$\text{visite} \leftarrow \text{visite} \cup \{v_{\text{next}}\}$ \;
	}
	Fermer la tournée : ajouter l'arête $(v_{\text{courant}}, v_1)$ \;
	\Return $\sigma$
\end{algorithm}

\begin{algorithm}[htbp]
	\caption{Amélioration 2-opt}
	\label{alg:2opt}
	\KwEntree{Tournée initiale $\sigma$}
	\KwSortie{Optimum local $\sigma^*$}
	$\text{amélioration} \leftarrow \texttt{Vrai}$ \;
	\While{$\text{amélioration}$}{
		$\text{amélioration} \leftarrow \texttt{Faux}$ \;
		\For{$i \leftarrow 1$ \KwTo $n-1$}{
			\For{$j \leftarrow i+1$ \KwTo $n$}{
				\If{$c_{i,i+1} + c_{j,j+1} > c_{i,j} + c_{i+1,j+1}$}{
					Inverser le segment $[\sigma(i+1), \ldots, \sigma(j)]$ \;
					$\text{amélioration} \leftarrow \texttt{Vrai}$ \;
				}
			}
		}
	}
	\Return $\sigma$
\end{algorithm}
	% ─────────────────────────────────────────────────────────────
	\subsection{M\'etaheuristiques}
	\label{subsec:tsp_meta}
	% ─────────────────────────────────────────────────────────────
	
	Les métaheuristiques sont des stratégies d'exploration de haut niveau
	conçues pour échapper aux optima locaux des heuristiques classiques, au
	prix d'un temps de calcul plus important.
	
	\textbf{Lin-Kernighan (LK).}
	L'heuristique de Lin-Kernighan généralise le principe du $k$-opt en
	autorisant des échanges de taille variable à chaque itération. À chaque
	étape, elle construit séquentiellement un ensemble d'arêtes à supprimer et
	à ajouter, en s'arrêtant dès qu'un gain positif est réalisable. Cette
	adaptabilité lui permet d'explorer des voisinages bien plus riches que le
	2-opt fixe, produisant des solutions dont le gap d'optimalité est
	typiquement inférieur à 1\% sur des instances euclidiennes standards.
	L'implémentation LKH de Helsgaun~\cite{helsgaun2000effective} est une référence performante parmi les heuristiques
	classiques du TSP ; sa qualité pratique ne constitue toutefois pas une garantie
	théorique d'optimalité.
	
	\textbf{Colonies de fourmis (ACO).}
	L'optimisation par colonies de fourmis simule le comportement collectif
	d'agents déposant une phéromone $\tau_{ij}$ sur les arêtes empruntées.
	La probabilité qu'une fourmi choisisse l'arête $(i,j)$ à l'étape $t$ est :
	$$
	p_{ij}(t) = \frac{[\tau_{ij}(t)]^\alpha \cdot [\eta_{ij}]^\beta}
	{\displaystyle\sum_{k \notin \text{visite}} [\tau_{ik}(t)]^\alpha
		\cdot [\eta_{ik}]^\beta}
	$$
	où $\eta_{ij} = 1/c_{ij}$ est une information heuristique (visibilité),
	et $\alpha, \beta > 0$ sont des paramètres contrôlant le compromis entre
	exploitation de la phéromone et exploration heuristique. Les phéromones
	s'évaporent à chaque itération selon un taux $\rho \in (0,1)$ :
	$$
	\tau_{ij}(t+1) = (1-\rho)\,\tau_{ij}(t) + \Delta\tau_{ij}(t)
	$$
	où $\Delta\tau_{ij}(t)$ est la contribution des fourmis ayant emprunté
	l'arête $(i,j)$ au cours de l'itération $t$.
	
	\begin{remark}{Bilan et transition vers le DRL}{bilan_methodes}
		L'ensemble de ces méthodes partage une limitation structurelle commune :
		chaque instance est résolue de manière indépendante, sans capitaliser sur
		l'expérience accumulée lors de la résolution d'instances précédentes. Le
		réglage des hyperparamètres (taux d'évaporation, paramètres $\alpha$ et
		$\beta$ pour ACO, longueur des chaînes LK) requiert une expertise
		manuelle spécifique à chaque famille de problèmes. C'est précisément
		cette incapacité à \textit{apprendre} qui motive le recours à
		l'Apprentissage par Renforcement Profond, étudié au Chapitre suivant.
	\end{remark}
	
	% ─────────────────────────────────────────────────────────────
	\section{Limites des approches classiques}
	\label{sec:tsp_limites}
	% ─────────────────────────────────────────────────────────────
	
L'analyse présentée dans ce chapitre donne une description systématique de certaines des approches classiques pour résoudre le problème du voyageur de commerce. Trois limitations majeures se dégagent, indépendamment de la famille de méthodes considérée.

\textbf{Limitation 1 --- Passage à l'échelle.}
Les algorithmes exacts (Bellman-Held-Karp, Branch-and-Cut) offrent une garantie d'optimalité, mais leur complexité temporelle augmente de façon exponentielle avec $n$. Ceci est visible dans le Tableau~\ref{tab:tsp_espace}, où il est clairement démontré qu'aucun algorithme exact ne sera capable de fonctionner efficacement au-delà de $n \approx 50$. Cependant, les heuristiques et les métaheuristiques offrent un moyen de surmonter cette barrière en temps polynomial, mais sans garantir l'optimalité. Selon le Théorème~\ref{thm:tsp_inapprox}, aucun algorithme en temps polynomial ne sera capable de fournir une approximation bornée pour le TSP général, sauf si $\mathcal{P} = \mathcal{NP}$.

\textbf{Limitation 2 --- Absence d'apprentissage.}
Quelle que soit la méthode utilisée --- le recuit simulé, l'optimisation par colonies de fourmis ou Lin-Kernighan --- il n'y a aucun moyen pour un algorithme d'utiliser les résolutions précédentes pour en apprendre. Chaque nouvelle instance doit être résolue à partir de zéro. Cette inefficacité est évidente : des milliers d'instances qui partagent la même structure sont traitées sans apporter la moindre amélioration au processus décisionnel de l'algorithme.

\textbf{Limitation 3 --- Conception manuelle et rigidité.}
Les méthodes heuristiques et métaheuristiques nécessitent des efforts importants de la part de l'expert pour concevoir et ajuster ces algorithmes (en choisissant les opérateurs de voisinage, ou en ajustant finement des paramètres comme $\alpha, \beta, \rho$ pour l'ACO, et $T_{\text{init}}, \alpha_{\text{cool}}$ pour le recuit simulé). Un tel processus d'ajustement est spécifique à certaines familles de problèmes et ne peut être généralisé automatiquement à d'autres variantes du TSP ou à d'autres problèmes combinatoires.
	
	\begin{table}[H]
		\centering
		\caption{Synthèse comparative des approches classiques pour le TSP.}
		\label{tab:synthese_classiques}
		\small
		\begin{tabular}{lcccc}
			\toprule
			\textbf{Méthode} & \textbf{Optimalité} & \textbf{Complexité}
			& \textbf{Passage à l'échelle} & \textbf{Apprentissage} \\
			\midrule
			Bellman-Held-Karp  & \checkmark & $\mathcal{O}(n^2 2^n)$      & $\times$   & $\times$ \\
			Branch-and-Cut     & \checkmark & exponentielle (pire cas)     & partiel    & $\times$ \\
			Nearest Neighbor   & $\times$   & $\mathcal{O}(n^2)$           & \checkmark & $\times$ \\
			2-opt              & $\times$   & $\mathcal{O}(n^2)$/iter      & \checkmark & $\times$ \\
			Lin-Kernighan      & $\times$   & variable & partiel    & $\times$ \\
			Recuit simulé      & $\times$   & paramétrique                 & partiel    & $\times$ \\
			ACO                & $\times$   & paramétrique                 & partiel    & $\times$ \\
			\bottomrule
		\end{tabular}
	\end{table}
	
\begin{remark}{Transition vers l'apprentissage par renforcement}{transition_drl}
	La colonne \og Apprentissage \fg{} du Tableau~\ref{tab:synthese_classiques} reste constamment à 0 pour toutes les techniques classiques. C'est précisément ce vide que cherche à combler l'apprentissage par renforcement profond. En apprenant une politique paramétrique $\pi_\theta$ sur un ensemble d'instances d'entraînement, un algorithme basé sur l'apprentissage par renforcement profond tente de surmonter les trois défis suivants :
	\begin{itemize}
		\item Produire une solution avec un coût d'inférence maîtrisé ; le coût dépend de l'architecture et de la stratégie de décodage et ne peut pas être supposé linéaire en $n$.
		\item Exploiter l'expérience accumulée entre les instances pour améliorer en moyenne la politique de décision, sans garantie de progression sur chaque instance.
		\item Apprendre la structure du problème sans avoir à concevoir d'opérateurs manuellement.
	\end{itemize}
	L'analyse de cette approche, partant de sa base mathématique pour se terminer par une mise en \oe{}uvre sur le problème du voyageur de commerce, constitue le thème des deux prochains chapitres.
\end{remark}
	
	
	% ============================================================
	%  CHAPITRE III : APPRENTISSAGE PAR RENFORCEMENT
	% ============================================================
	\chapter{Apprentissage par renforcement}
	\label{chap:apprentissage_renf}
	
	Ce chapitre détaille les fondements théoriques de l'Apprentissage par Renforcement Profond (DRL) et leur adaptation au Problème du Voyageur de Commerce (TSP).
	
	\section{Mod\`ele math\'ematique (MDP)}
	\label{sec:modele_mdp}
	
Pour mettre en \oe{}uvre l'apprentissage par renforcement dans un problème d'optimisation combinatoire statique comme le Problème du Voyageur de Commerce (TSP) il est nécessaire de modéliser la recherche de la solution comme un processus de prise de décision dans le temps.Ce type de modélisation peut être réalisé en utilisant la construction mathématique connue sous le nom de Processus de Décision de Markov (MDP).
	
	\begin{definition}{Processus de Décision Markovien (MDP)}{def_mdp}
		Un MDP est formellement défini par un tuple $(\mathcal{S}, \mathcal{A}, \mathcal{P}, \mathcal{R}, \gamma)$, où :
		\begin{itemize}
			\item $\mathcal{S}$ est l'espace des états,
			\item $\mathcal{A}$ est l'espace des actions,
			\item $\mathcal{P}(s_{t+1} \mid s_t, a_t)$ modélise la dynamique de transition,
			\item $\mathcal{R}(s_t, a_t)$ est la fonction de récompense,
			\item $\gamma \in [0, 1]$ est le facteur d'escompte (ou d'actualisation).
		\end{itemize}
	\end{definition}
	
	Dans le contexte spécifique de la construction d'une tournée hamiltonienne, ce tuple s'instancie de la manière suivante à chaque pas de temps $t$ :
	
	\begin{itemize}
		\item \textbf{L'espace des états ($\mathcal{S}$)} : L'état $s_t \in \mathcal{S}$ représente la situation de l'environnement à l'étape $t$. Pour le TSP, cet état doit encapsuler la structure statique du problème (le graphe complet contenant les coordonnées des villes) et l'information dynamique (la tournée partielle, c'est-à-dire la séquence des villes déjà visitées jusqu'à l'instant $t$).
		
		\item \textbf{L'espace des actions ($\mathcal{A}$)} : Une action $a_t \in \mathcal{A}$ correspond à la sélection du prochain n\oe{}ud à visiter. Pour garantir la validité de la solution (respect des contraintes d'élimination des sous-tours et de visite unique étudiées au Chapitre~\ref{chap:tsp}), l'espace des actions admissibles se restreint dynamiquement aux villes non encore visitées.
		
		\item \textbf{La dynamique de transition ($\mathcal{P}$)} : Dans notre modélisation du TSP, l'environnement est purement déterministe. Dès lors que l'agent sélectionne une ville valide, la probabilité que cette ville soit ajoutée à la tournée partielle pour former l'état $s_{t+1}$ est de $1$.
		
		\item \textbf{La fonction de récompense ($\mathcal{R}$)} : C'est le signal d'apprentissage central. Puisque l'objectif global est de minimiser la distance totale, la récompense immédiate $r_t$ associée au choix d'une ville est définie comme la distance négative entre la ville actuelle et la ville choisie : $r_t = -c_{\sigma(t), \sigma(t+1)}$. 
		
		\item \textbf{Le facteur d'escompte ($\gamma$)} : Le TSP est un problème épisodique à horizon fini. On fixe ici $\gamma = 1$ afin que toutes les composantes du coût de la tournée contribuent sans actualisation au retour total.
	\end{itemize}
	
	Ce formalisme permet à un agent apprenant (une politique paramétrée $\pi_\theta$) d'interagir avec le graphe en observant l'état actuel et en itérant pour construire une solution complète, comme l'illustre la Figure~\ref{fig:mdp_tsp}.
	
	\begin{figure}[H]
		\centering
		\begin{tikzpicture}[>=stealth, auto, node distance=3cm, font=\sffamily]
			
			% Noeuds principaux avec les couleurs du template (ThBlue et ThGreen)
			\node[draw=ThBlue, thick, fill=ThBlue!10, rounded corners, minimum width=4cm, minimum height=1.5cm, align=center] (agent) at (0,0) {\textbf{Agent (Réseau de Neurones)}\\Politique $\pi_\theta(a_t|s_t)$};
			
			\node[draw=ThGreen!80!black, thick, fill=ThGreen!10, rounded corners, minimum width=4cm, minimum height=1.5cm, align=center] (env) at (0,-4) {\textbf{Environnement (Graphe TSP)}\\Graphe complet, Villes à visiter};
			
			% Flèche Action
			\draw[->, thick, ThBlue] (agent.east) -- ++(2.5,0) |- node[near start, right, align=center] {Action $a_t$\\ \textit{(Choix de la prochaine ville)}} (env.east);
			
			% Flèche État / Récompense
			\draw[->, thick, ThGreen!80!black] (env.west) -- ++(-2.5,0) |- node[near start, left, align=center] {État $s_{t+1}$\\ \textit{(Tournée partielle)}\\ \\Récompense $r_{t+1}$\\ \textit{($-c_{\sigma(t), \sigma(t+1)}$)}} (agent.west);
		\end{tikzpicture}
		\caption{Modélisation du Problème du Voyageur de Commerce sous forme de Processus de Décision Markovien. À chaque étape, l'agent observe la tournée partielle et sélectionne la prochaine ville, recevant en retour une récompense basée sur la distance parcourue.}
		\label{fig:mdp_tsp}
	\end{figure}
	% ─────────────────────────────────────────────────────────────
	\section{Fonctions de valeur}
	\label{sec:fonctions_valeur}
	% ─────────────────────────────────────────────────────────────
	
Cependant, pour un processus de décision de Markov, se contenter de spécifier la dynamique du jeu (qui comprend les états, les actions et les récompenses) n'est pas suffisant pour résoudre le problème. Pour permettre à un agent d'évaluer la qualité de ses décisions, il faut lui fournir une méthode qui mesure les \og retours futurs attendus \fg{} à partir d'un certain état, sous une politique de décision donnée $\pi$.
	\begin{definition}{Fonctions de Valeur d'État et d'Action}{def_valeur}
		Soit un agent suivant une politique $\pi$. 
		\begin{itemize}
			\item \textbf{La fonction de valeur d'état $V^\pi(s)$} mesure l'espérance du retour cumulé si l'agent démarre dans l'état $s$ et suit la politique $\pi$ :
			\begin{equation}
				V^\pi(s) = \E_\pi \left[ \sum_{k=0}^{\infty} \gamma^k r_{t+k+1} \;\middle|\; s_t = s \right]
			\end{equation}
			\item \textbf{La fonction de valeur d'action $Q^\pi(s, a)$} (ou Q-valeur) mesure l'espérance du retour cumulé si l'agent démarre dans l'état $s$, prend immédiatement l'action $a$, puis suit la politique $\pi$ :
			\begin{equation}
				Q^\pi(s, a) = \E_\pi \left[ \sum_{k=0}^{\infty} \gamma^k r_{t+k+1} \;\middle|\; s_t = s, a_t = a \right]
			\end{equation}
		\end{itemize}
	\end{definition}
	
L'équation de Bellman est la véritable force qui fait fonctionner l'apprentissage par renforcement. L'idée de cette équation peut être exprimée comme suit : la valeur de l'état actuel ne dépend pas seulement de la récompense directe de cet état, mais aussi de l'estimation de l'état futur. Concernant notre tâche, l'objectif de l'agent est de trouver la meilleure politique $\pi^*$, c'est-à-dire celle qui maximise la récompense attendue pour chaque décision qu'il prend, afin de créer un itinéraire idéal.
	
	Les \textbf{Équations d'Optimalité de Bellman} s'écrivent ainsi :
	\begin{align}
		V^*(s) &= \max_{a \in \mathcal{A}} Q^*(s, a) \label{eq:bellman_v} \\
		Q^*(s, a) &= \E \left[ r_{t+1} + \gamma V^*(s_{t+1}) \;\middle|\; s_t = s, a_t = a \right] \label{eq:bellman_q}
	\end{align}
	
	\begin{remark}{Adaptation au TSP (Environnement Déterministe)}{rem_tsp_det}
		Dans le cadre du TSP, deux simplifications majeures s'opèrent sur l'équation~\eqref{eq:bellman_q} :
		\begin{enumerate}
			\item L'environnement étant purement déterministe, l'espérance $\E[\cdot]$ disparaît. L'action d'aller vers la ville $v_j$ mène avec une probabilité de $1$ à l'état suivant.
			\item L'horizon étant fini et toutes les décisions étant d'égale importance pour la tournée, le facteur d'actualisation $\gamma$ est fixé à $1$.
		\end{enumerate}
		La Q-valeur optimale se réduit alors à une simple addition entre le coût immédiat de l'arête et l'évaluation du reste du graphe : $Q^*(s_t, a_t) = -c(v_t, v_{t+1}) + V^*(s_{t+1})$.
	\end{remark}
	
	La Figure~\ref{fig:bellman_tree} illustre visuellement cette décomposition récursive : l'évaluation d'un état passe par l'exploration des actions admissibles (les villes non visitées), chacune menant à un nouvel état (une nouvelle tournée partielle).
	
	\begin{figure}[H]
		\centering
		\begin{tikzpicture}[>=stealth, thick, level distance=2.5cm, font=\sffamily\small]
			% Définition des styles en reprenant tes couleurs
			\tikzstyle{state}=[circle, draw=ThBlue, fill=ThBlue!10, minimum size=1.2cm, inner sep=0pt, font=\bfseries\color{ThBlue!80!black}]
			\tikzstyle{action}=[circle, draw=ThRed, fill=ThRed!10, minimum size=0.8cm, inner sep=0pt, font=\bfseries\color{ThRed!80!black}]
			\tikzstyle{nextstate}=[circle, draw=ThGreen!80!black, fill=ThGreen!10, minimum size=1.2cm, inner sep=0pt, font=\bfseries\color{ThGreen!80!black}]
			
			% Racine (État S)
			\node[state] (S) at (0, 0) {$s$};
			\node[above=0.2cm of S, color=ThBlue, font=\bfseries] {Valeur $V(s)$};
			
			% Actions (A1, A2, A3)
			\node[action] (A1) at (-3, -2.5) {$a_1$};
			\node[action] (A2) at (0, -2.5) {$a_2$};
			\node[action] (A3) at (3, -2.5) {$a_3$};
			
			% Flèches de l'état vers les actions (Politique)
			\draw[->, ThBlue] (S) -- (A1) node[midway, above left, font=\scriptsize, color=black] {$\pi(a_1|s)$};
			\draw[->, ThBlue] (S) -- (A2) node[midway, right, font=\scriptsize, color=black] {$\pi(a_2|s)$};
			\draw[->, ThBlue] (S) -- (A3) node[midway, above right, font=\scriptsize, color=black] {$\pi(a_3|s)$};
			
			% Texte pour les Q-valeurs
			\node[right=0.2cm of A3, color=ThRed, font=\bfseries] {Q-valeur $Q(s, a_3)$};
			
			% Nouveaux États (S')
			\node[nextstate] (S1) at (-3, -5) {$s'_{1}$};
			\node[nextstate] (S2) at (0, -5) {$s'_{2}$};
			\node[nextstate] (S3) at (3, -5) {$s'_{3}$};
			
			% Flèches des actions vers les nouveaux états (Dynamique)
			\draw[->, ThRed] (A1) -- (S1) node[midway, left, font=\scriptsize, color=black] {$r_1$};
			\draw[->, ThRed] (A2) -- (S2) node[midway, right, font=\scriptsize, color=black] {$r_2$};
			\draw[->, ThRed] (A3) -- (S3) node[midway, right, font=\scriptsize, color=black] {$r_3$};
			
			\node[below=0.2cm of S2, color=ThGreen!80!black, font=\bfseries] {Valeur $V(s')$};
			
			% Accolade explicative décalée vers la droite (x=5.5)
			\draw[decorate, decoration={brace, amplitude=10pt, mirror}, color=ThSilver] 
			(7.5, -5.2) -- (7.5, 0.2) node[midway, right=15pt, align=left] {Décomposition\\récursive\\de Bellman};
			
		\end{tikzpicture}
		\caption{Arbre de décision illustrant les relations de Bellman. La valeur d'un état $V(s)$ dépend des Q-valeurs des actions possibles $Q(s,a)$, qui elles-mêmes dépendent de la récompense immédiate $r$ et de la valeur du nouvel état atteint $V(s')$.}
		\label{fig:bellman_tree}
	\end{figure}
	
	% ─────────────────────────────────────────────────────────────
	\section{Algorithmes principaux}
	\label{sec:algos_principaux}
	% ─────────────────────────────────────────────────────────────
	
	Étant donné la taille énorme de l'espace des états dans le problème du voyageur de commerce (TSP), les algorithmes traditionnels d'apprentissage par renforcement comme le \textit{Q-Learning} deviennent inopérants pour ce problème. Cela nécessite la paramétrisation de la politique par un vecteur de poids $\theta$ (généralement un réseau de neurones) : $\pi_\theta(a|s)$.
	
	Lorsqu'on traite de problèmes d'optimisation combinatoire qui ont un espace d'action évolutif (le nombre de villes diminue à chaque itération), la littérature recommande fortement les approches de Gradient de Politique (\textit{Policy Gradient}). L'objectif de ces approches est de maximiser directement la récompense attendue de la politique paramétrée :
	\begin{equation}
		J(\theta) = \E_{\tau \sim \pi_\theta} \left[ R(\tau) \right]
	\end{equation}
	où $\tau$ est une trajectoire complète (une tournée) et $R(\tau)=-L(\tau)$, avec $L(\tau)$ la longueur totale de la tournée. Cette convention est celle de la maximisation d'une récompense ; le chapitre expérimental reformule l'apprentissage en minimisation du coût $L(\tau)$.
	
	\subsection{L'algorithme REINFORCE}
	
	\textbf{Théorème du gradient de politique.}
	
	Le gradient de la fonction objectif $J(\theta)$ peut être écrit sous forme analytique en utilisant les probabilités des actions sélectionnées :
	
	\begin{equation}
		\nabla_\theta J(\theta) = \E_{\tau \sim \pi_\theta} \left[ \sum_{t=0}^{n-1} \nabla_\theta \log \pi_\theta(a_t \mid s_t) R(\tau) \right]
	\end{equation}
	
	L'algorithme REINFORCE (Williams, 1992~\cite{williams1992simple}) s'appuie directement sur cette formule. Concrètement, l'agent construit une trajectoire $\tau$, calcule le retour $R(\tau)$ et modifie les poids $\theta$. Il augmente ainsi la probabilité des actions qui ont produit une \og bonne \fg{} trajectoire (plus courte) et diminue celle des actions ayant produit une \og mauvaise \fg{} trajectoire (plus longue).
	
	\begin{remark}{Le problème de la variance}{rem_variance_reinforce}
		Bien que théoriquement fondé, REINFORCE souffre d'une très forte variance empirique. La récompense $R(\tau)$ étant calculée à la toute fin de l'épisode (la tournée complète), il est difficile pour l'agent de déterminer \textit{laquelle} des $n$ décisions successives a été la cause principale du succès ou de l'échec. Ce phénomène est connu sous le nom de problème d'assignation du crédit (\textit{credit assignment problem}).
	\end{remark}
	
	Pour pallier cette variance, on introduit généralement une ligne de base (\textit{baseline}) $b(s_t)$. Le gradient devient proportionnel à $(R(\tau) - b(s_t))$. Si la tournée est meilleure que la ligne de base, les probabilités des actions sont renforcées.
	
	\subsection{L'architecture Acteur-Critique (Actor-Critic)}
	
	Pour stabiliser l'apprentissage et formaliser la ligne de base, l'architecture \textbf{Acteur-Critique} introduit un second modèle d'apprentissage. Au lieu d'attendre la fin de l'épisode ou d'utiliser une ligne de base statique, l'agent est scindé en deux entités concurrentes et complémentaires :
	
	\begin{itemize}
		\item \textbf{L'Acteur ($\pi_\theta$)} : C'est la politique proprement dite. Il observe l'état $s_t$ (le graphe et la tournée partielle) et décide de la prochaine action $a_t$ (la prochaine ville à visiter).
		\item \textbf{Le Critique ($V_\phi$)} : Paramétré par un second réseau de poids $\phi$, son rôle est d'estimer une valeur de référence pour l'état. Dans l'implémentation expérimentale, la baseline est apprise à partir de la longueur de la tournée complète afin de réduire la variance de l'estimateur de gradient.
	\end{itemize}
	
	L'erreur de prédiction du critique permet de calculer la \textbf{fonction d'Avantage} $A(s_t, a_t)$ :
	\begin{equation}
		A(s_t, a_t) = Q(s_t, a_t) - V_\phi(s_t) \approx r_{t+1} + V_\phi(s_{t+1}) - V_\phi(s_t)
	\end{equation}
	
	L'avantage indique si l'action choisie par l'Acteur s'est avérée meilleure ou pire que ce que le Critique avait prédit. Les deux réseaux s'améliorent mutuellement : l'Acteur met à jour ses probabilités en fonction de l'Avantage, et le Critique met à jour ses prédictions pour minimiser l'erreur quadratique moyenne avec la vraie récompense. Ce mécanisme est illustré par la Figure~\ref{fig:actor_critic}.
	
	\begin{figure}[H]
		\centering
		\begin{tikzpicture}[>=stealth, auto, node distance=2.5cm, font=\sffamily\small]
			% Couleurs
			\definecolor{actorCol}{RGB}{180, 0, 0} % ThRed
			\definecolor{criticCol}{RGB}{0, 51, 102} % ThBlue
			\definecolor{envCol}{RGB}{30, 130, 30} % ThGreen
			
			% Blocs principaux
			\node[draw=envCol, thick, fill=envCol!10, rounded corners, minimum width=3.5cm, minimum height=1.2cm, align=center] (env) at (0, -4.5) {\textbf{Environnement}\\(Graphe TSP)};
			
			\node[draw=actorCol, thick, fill=actorCol!10, rounded corners, minimum width=3cm, minimum height=1.2cm, align=center] (actor) at (-3.5, 0) {\textbf{Acteur $\pi_\theta(a|s)$}\\(Politique)};
			
			\node[draw=criticCol, thick, fill=criticCol!10, rounded corners, minimum width=3cm, minimum height=1.2cm, align=center] (critic) at (3.5, 0) {\textbf{Critique $V_\phi(s)$}\\(Évaluateur)};
			
			\node[draw=ThSilver, thick, fill=ThSilver!10, rectangle, minimum width=3cm, minimum height=1cm, align=center] (td) at (3.5, -2.5) {\textbf{Calcul de}\\\textbf{l'Avantage $A$}};
			
			% Flèches État (Noires) - Montent au centre et bifurquent sur les côtés
			% Vers l'Acteur (à gauche)
			\draw[->, thick] ([xshift=-0.5cm]env.north) -- ++(0, 3.6) -- node[above] {État $s_t$} ([yshift=-0.3cm]actor.east);
			% Vers le Critique (à droite)
			\draw[->, thick] ([xshift=0.5cm]env.north) -- ++(0, 3.6) -- node[above] {État $s_t$} ([yshift=-0.3cm]critic.west);
			
			% Flèche Action (Rouge)
			\draw[->, thick, actorCol] (actor.west) -- ++(-1,0) |- node[near start, left] {Action $a_t$} (env.west);
			
			% Flèche Récompense (Verte)
			\draw[->, thick, envCol] (env.east) -| node[pos=0.5, below] {Récompense $r_{t+1}$} (td.south);
			
			% Flèche Valeur V(s) du Critique vers TD (Bleue)
			\draw[->, thick, criticCol] ([xshift=0.8cm]critic.south) -- node[right] {$V_\phi(s_t)$} ([xshift=0.8cm]td.north);
			
			% Flèche Avantage (Mise à jour Acteur) - Le grand pont supérieur
			\draw[->, dashed, thick, ThSilver!80!black] 
			(td.west) -- ++(-0.5, 0) 
			|- node[pos=0.75, above] {Signal de mise à jour (Avantage)} 
			([yshift=1.2cm]actor.north) 
			-- (actor.north);
			
			% Flèche Erreur TD (Mise à jour Critique)
			\draw[->, dashed, thick, ThSilver!80!black] (td.east) -- ++(0.8,0) |- node[pos=0.25, right, align=center] {Erreur\\TD} (critic.east);
			
		\end{tikzpicture}
		\caption{Architecture Acteur-Critique. L'Acteur décide du prochain n\oe{}ud à visiter en s'appuyant sur l'état du graphe, tandis que le Critique évalue la qualité de cet état. La différence entre la récompense obtenue et l'évaluation du Critique sert de signal d'apprentissage pour stabiliser l'optimisation.}
		\label{fig:actor_critic}
	\end{figure}
	
	Cette architecture constitue la fondation algorithmique des modèles que nous implémenterons au Chapitre~\ref{chap:impl_resultats}. L'intégration de ces algorithmes avec des réseaux de neurones profonds, spécifiquement adaptés au traitement de graphes, fait l'objet des sections suivantes.
	
	% ─────────────────────────────────────────────────────────────
	\section{Deep Reinforcement Learning}
	\label{sec:drl}
	% ─────────────────────────────────────────────────────────────
	
	Les algorithmes présentés dans la section précédente reposent sur l'optimisation de fonctions paramétrées. Cependant, dans le cadre de l'optimisation combinatoire et particulièrement pour le TSP, la représentation de l'état $s_t$ (un graphe partiel et un ensemble de n\oe{}uds restants) est hautement complexe et de dimension variable.
	
	C'est ici qu'intervient l'Apprentissage par Renforcement Profond (\textit{Deep Reinforcement Learning} ou DRL). Le DRL fusionne la puissance d'abstraction du Deep Learning avec le cadre d'optimisation séquentielle du Reinforcement Learning. L'idée fondamentale est de remplacer les probabilités tabulaires ou les estimateurs linéaires par des Réseaux de Neurones Profonds (DNN), qui agissent comme des approximateurs de fonctions universels pour contrer l'explosion combinatoire de l'espace des états.
	
	\begin{definition}{Approximation par Réseaux de Neurones}{def_approx_nn}
		Dans le cadre du DRL appliqué à l'architecture Acteur-Critique, les fonctions mathématiques sont remplacées par deux réseaux de neurones distincts :
		\begin{itemize}
			\item \textbf{Le réseau Acteur (Politique)} : paramétré par les poids $\theta$, il prend en entrée l'état $s_t$ et retourne une distribution de probabilité sur l'espace d'action : $\pi_\theta(a_t \mid s_t)$.
			\item \textbf{Le réseau Critique (Valeur)} : paramétré par les poids $\phi$, il prend en entrée l'état $s_t$ et retourne une estimation scalaire de la valeur de cet état : $V_\phi(s_t)$.
		\end{itemize}
	\end{definition}
	
	\textbf{Le défi de la représentation de l'état.}
	L'un des défis majeurs du DRL pour le TSP est la nature de la donnée en entrée. Contrairement au traitement d'images (des pixels figés dans une grille fixe) ou au traitement du langage (des mots discrets), le TSP manipule des nuages de points (les villes) qui n'ont pas d'ordre intrinsèque et dont la taille (le nombre $n$) varie selon l'instance.
	
	Pour qu'un réseau de neurones puisse traiter cet état brut, il est nécessaire de projeter les coordonnées géographiques (généralement en 2D, $(x,y)$) dans un espace latent de plus grande dimension (par exemple $d = 128$). Cette opération, appelée \textbf{Plongement} (ou \textit{Embedding}), permet au modèle de créer des représentations vectorielles riches pour chaque ville. Ces vecteurs capturent non seulement la position spatiale locale, mais sont ensuite mis à jour pour intégrer la relation structurelle globale avec toutes les autres villes du graphe (voir Figure~\ref{fig:drl_approx}).
	
	\begin{figure}[H]
		\centering
		\begin{tikzpicture}[>=stealth, auto, node distance=2.5cm, font=\sffamily\small]
			
			% Styles
			\tikzstyle{tensor}=[draw=ThSilver, thick, fill=ThSilver!10, rounded corners, minimum width=2cm, minimum height=1.5cm, align=center]
			\tikzstyle{network}=[draw=ThBlue, thick, fill=ThBlue!10, rounded corners, minimum width=2.5cm, minimum height=1.5cm, align=center]
			\tikzstyle{output}=[draw=ThRed, thick, fill=ThRed!10, rounded corners, minimum width=2.5cm, minimum height=1.5cm, align=center]
			
			% Noeuds
			\node[tensor] (input) at (0,0) {\textbf{État Brut $s_t$}\\Coordonnées $(x,y)$\\Masque de visite};
			
			\node[network] (embedding) at (4.5,0) {\textbf{Couche de}\\\textbf{Plongement}\\(Projection dim. $d$)};
			
			\node[network] (dnn) at (9,0) {\textbf{Réseau de}\\\textbf{Neurones Profond}\\(Paramètres $\theta$)};
			
			\node[output] (action) at (13.5,0) {\textbf{Distribution}\\\textbf{d'Actions}\\$\pi_\theta(\cdot|s_t)$};
			
			% Flèches
			\draw[->, thick, ThBlue] (input) -- (embedding);
			\draw[->, thick, ThBlue] (embedding) -- (dnn);
			\draw[->, thick, ThRed] (dnn) -- (action);
			
			% Accolade de feature extraction
			\draw[decorate, decoration={brace, amplitude=6pt}, color=ThSilver!80!black, thick]
			(2.5, 1.2) -- (10.5, 1.2) node[midway, above=10pt] {Extraction et traitement des caractéristiques spatiales};
			
		\end{tikzpicture}
		\caption{Flux de traitement dans un agent DRL. L'état brut est d'abord projeté dans un espace latent (Embedding) de dimension $d$ avant d'être traité par les couches profondes pour produire la politique de décision.}
		\label{fig:drl_approx}
	\end{figure}
	
	La nature exacte de l'architecture profonde utilisée dans le bloc central est déterminante pour les performances. L'utilisation de réseaux denses classiques (Perceptrons Multicouches) ou convolutifs (CNN) s'avère inefficace pour capturer les invariances des graphes. C'est pourquoi la modélisation neuronale du TSP fait massivement appel à des architectures dites \textit{Sequence-to-Sequence} enrichies de mécanismes d'Attention, qui font l'objet de la section suivante.
	
	% ─────────────────────────────────────────────────────────────
	\section{Modélisation du TSP et Architectures Neuronales}
	\label{sec:modelisation_tsp}
	% ─────────────────────────────────────────────────────────────
	
	L'application du Deep Reinforcement Learning au TSP requiert une architecture capable de traiter des graphes de taille variable et de produire une séquence ordonnée de n\oe{}uds, tout en respectant les contraintes du problème (ne visiter chaque ville qu'une seule fois).
	
	Les implémentations de référence dans la littérature récente (telles que celles étudiées dans les dépôts \texttt{learning-tsp} et \texttt{TSP\_DRL\_PtrNet}) suivent un pipeline de résolution unifié en cinq étapes distinctes, illustré par la Figure~\ref{fig:pipeline_drl}.
	
	\begin{figure}[H]
		\centering
		\resizebox{\textwidth}{!}{%
			\begin{tikzpicture}[>=stealth, auto, node distance=0.8cm, font=\sffamily\small]
				% Définition des styles inspirés de ta charte
				\tikzstyle{pipeBlock}=[draw, thick, rounded corners=4pt, minimum width=2.8cm, minimum height=1.5cm, align=center, text=black]
				
				% Création des 5 blocs du pipeline
				\node[pipeBlock, fill=ThSilver!20, draw=ThSilver!80!black] (b1) at (0,0) {\textbf{Définition}\\\textbf{du Problème}};
				
				\node[pipeBlock, fill=ThBlue!15, draw=ThBlue] (b2) [right=of b1] {\textbf{Plongement}\\\textbf{du Graphe}};
				
				\node[pipeBlock, fill=violet!15, draw=violet!80!black] (b3) [right=of b2] {\textbf{Décodage de}\\\textbf{la Solution}};
				
				\node[pipeBlock, fill=ThRed!15, draw=ThRed] (b4) [right=of b3] {\textbf{Recherche de}\\\textbf{la Solution}};
				
				\node[pipeBlock, fill=orange!20, draw=orange!80!black] (b5) [right=of b4] {\textbf{Apprentissage}\\\textbf{de la Politique}};
				
				% Labels alphabétiques sous les blocs
				\node[below=0.2cm of b1] {(a)};
				\node[below=0.2cm of b2] {(b)};
				\node[below=0.2cm of b3] {(c)};
				\node[below=0.2cm of b4] {(d)};
				\node[below=0.2cm of b5] {(e)};
				
				% Flèches de connexion
				\draw[->, thick, ThSilver!80!black] (b1) -- (b2);
				\draw[->, thick, ThSilver!80!black] (b2) -- (b3);
				\draw[->, thick, ThSilver!80!black] (b3) -- (b4);
				\draw[->, thick, ThSilver!80!black] (b4) -- (b5);
				
			\end{tikzpicture}%
		}
		\caption{Pipeline général de l'Optimisation Combinatoire Neurale. Le problème brut (a) est transformé en vecteurs latents (b). Un décodeur calcule les probabilités de transition (c) exploitées par un algorithme de recherche (d) pour extraire une tournée. La qualité de cette tournée sert de signal pour optimiser les poids du modèle (e).}
		\label{fig:pipeline_drl}
	\end{figure}
	
	\subsection{Analyse du Pipeline de Résolution}
	
	Ce pipeline structure le flux de données depuis les coordonnées bidimensionnelles jusqu'à la mise à jour des paramètres du réseau $\theta$.
	
	\begin{enumerate}[label=\textbf{(\alph*)}, leftmargin=*]
		\item \textbf{Définition du Problème (Problem Definition) :} Le TSP est défini par un ensemble de n\oe{}uds en 2D. En entrée, l'instance est une matrice $X \in \R^{n \times 2}$ contenant les coordonnées spatiales des villes.
		
		\item \textbf{Plongement du Graphe (Graph Embedding) :} Les coordonnées brutes $(x,y)$ sont insuffisantes pour capturer la topologie complexe d'une tournée. Chaque ville est projetée dans un espace vectoriel de haute dimension $d$ (souvent $d=128$) via une transformation linéaire ou un Réseau de Neurones sur Graphes (GNN). Ce plongement permet à chaque n\oe{}ud d'agréger des informations sur son voisinage, créant une représentation globale du graphe encodée dans une matrice $H \in \R^{n \times d}$.
		
		\item \textbf{Décodage de la Solution (Solution Decoding) :} C'est le c\oe{}ur de l'architecture. Contrairement aux modèles de traduction (Seq2Seq) qui génèrent des mots depuis un vocabulaire fixe, le décodeur du TSP doit pointer vers un élément de la séquence d'entrée.
		L'architecture \textbf{Pointer Network} (Vinyals et al., 2015) utilise un mécanisme d'\textit{Attention} pour calculer une distribution de probabilité sur les n\oe{}uds restants. À l'étape $t$, le décodeur compare l'état de la tournée partielle avec les \textit{embeddings} des villes non visitées, produisant un vecteur de probabilités $\pi_\theta(a_t \mid s_t)$.
		
		\item \textbf{Recherche de la Solution (Solution Search) :} Une fois les probabilités de transition calculées par le réseau, une stratégie de décodage est appliquée pour extraire la tournée finale. Deux approches dominent :
		\begin{itemize}
			\item \textit{Décodage Glouton (Greedy Search)} : Sélection déterministe de la ville ayant la probabilité maximale (utilisée en inférence pour minimiser le temps de calcul).
			\item \textit{Échantillonnage (Sampling)} : Sélection stochastique basée sur la distribution de probabilités (utilisée pendant l'entraînement pour encourager l'exploration).
		\end{itemize}
		
		\item \textbf{Apprentissage de la Politique (Policy Learning) :} La tournée générée à l'étape précédente est évaluée. Son coût $L(\pi)$ (la distance euclidienne totale) est renvoyé à l'algorithme d'apprentissage (REINFORCE ou Actor-Critic, comme vu en Section~\ref{sec:algos_principaux}). L'erreur TD ou l'Avantage permet de rétropropager le gradient et d'ajuster les poids $\theta$ du Pointer Network.
	\end{enumerate}
	
	Cette modélisation élégante clôt le chapitre théorique. Elle démontre comment un problème d'optimisation purement discret peut être rendu différentiable et résolu par des réseaux de neurones. L'implémentation de cette architecture, son entraînement sur des milliers d'instances, et la comparaison de ses performances empiriques avec les heuristiques classiques font l'objet du Chapitre~\ref{chap:impl_resultats}.
	
	
	% ============================================================
	%  CHAPITRE IV : IMPLEMENTATION ET RESULTATS
	% ============================================================

\chapter{Implémentation et résultats}
\label{chap:impl_resultats}

\section{Vue d'ensemble du chapitre}
\label{sec:ch4-overview}

Ce chapitre présente une approche par apprentissage par renforcement profond
(Deep RL) pour la résolution du problème du voyageur de commerce (TSP), en
adaptant l'implémentation \texttt{TSP\_DRL\_PtrNet}~\cite{rintarooo_repo} de
l'architecture proposée par Bello \textit{et al.}~\cite{bello2016neural},
elle-même construite sur le Pointer Network de Vinyals \textit{et al.}~\cite{vinyals2015pointer}.
L'objectif est double : (i) formaliser le TSP comme un problème d'apprentissage
séquentiel entraîné sans exemples étiquetés, et (ii) comparer les performances
de cette approche (qualité des tournées et temps de résolution) à une sélection
de méthodes classiques étudiées au Chapitre~\ref{chap:fondements} et retenues
comme références dans le protocole expérimental principal : le plus proche
voisin (PPV), le PPV amélioré par 2-opt (PPV+2-opt) et l'optimisation par
colonies de fourmis (ACO). Les autres méthodes présentées dans le cadre
théorique, notamment le recuit simulé, les algorithmes génétiques et
Bellman--Held--Karp, ne sont pas incluses dans cette comparaison expérimentale
principale.

Le chapitre est organisé comme suit : la section~\ref{sec:ch4-theorie} formalise
le cadre théorique ; la section~\ref{sec:ch4-implementation} décrit
l'implémentation et son adaptation à l'environnement Google Colab ; la
section~\ref{sec:ch4-resultats} présente les résultats comparatifs obtenus
expérimentalement ; la section~\ref{sec:ch4-discussion} discute les limites
de l'approche et les résultats obtenus.

% ================================================================
\section{Formalisation théorique}
\label{sec:ch4-theorie}
% ================================================================

\subsection{Le TSP comme processus de décision markovien}
\label{sec:ch4-mdp}

On formalise la construction séquentielle d'un tour comme un processus de
décision markovien (MDP) $(\mathcal{S}, \mathcal{A}, P, R)$.

\begin{definition}{Formulation MDP du TSP}{mdp-tsp}
	Étant donné une instance de $n$ villes de coordonnées $\{x_1, \dots, x_n\}$ :
	\begin{itemize}
		\item \textbf{État} $s_i$ : la ville courante et l'ensemble des villes
		déjà visitées (encodé via un masque binaire) ;
		\item \textbf{Action} $a_i$ : le choix de la prochaine ville à visiter
		parmi celles non encore visitées ;
		\item \textbf{Récompense} : dans l'implémentation évaluée, elle est donnée en fin d'épisode,
		$R(\pi)=-L(\pi)$, où $L(\pi)$ est la longueur totale du tour $\pi$ ;
		\item \textbf{Politique} $p_\theta(\pi \mid s)$ : distribution de
		probabilité sur les permutations, paramétrée par un réseau de neurones
		(le Pointer Network).
	\end{itemize}
\end{definition}

\begin{remark}{Récompense terminale}{reward-terminal}
	Le modèle expérimental attribue une récompense terminale $R(\pi)=-L(\pi)$ : la qualité de la tournée n'est connue qu'une fois le cycle complet construit. Cela justifie le recours à une estimation de retour de type Monte-Carlo (REINFORCE). La récompense théorique par transition décrite au chapitre précédent est une autre convention ; avec $\gamma=1$ et en incluant toutes les arêtes du cycle, sa somme vaut également $-L(\pi)$, mais la distribution temporelle du signal d'apprentissage diffère.
\end{remark}

\subsection{Le Pointer Network comme politique}
\label{sec:ch4-ptrnet}

Un encodeur récurrent lit la séquence de coordonnées et produit les états
cachés $e_1, \dots, e_n$. À chaque étape $i$ de décodage, un score de
compatibilité est calculé entre l'état du décodeur $d_i$ et chaque état de
l'encodeur :

\begin{equation}
	u_i^j = v^\top \tanh(W_1 e_j + W_2 d_i), \qquad j = 1, \dots, n
	\label{eq:ptr-score}
\end{equation}

La politique du modèle est obtenue par normalisation softmax de ces scores,
après masquage des villes déjà visitées :

\begin{equation}
	p_\theta(a_i = j \mid s_i) =
	\begin{cases}
		\dfrac{\exp(u_i^j)}{\sum_{k \notin \text{visité}} \exp(u_i^k)} & \text{si } j \notin \text{visité} \\[6pt]
		0 & \text{sinon}
	\end{cases}
	\label{eq:ptr-softmax}
\end{equation}

\begin{remark}{Différence avec l'attention classique}{ptr-vs-attention}
	Dans un mécanisme d'attention standard, les poids issus du softmax servent à
	calculer une moyenne pondérée injectée dans une couche de sortie sur
	vocabulaire fixe. Le Pointer Network court-circuite cette étape : la
	distribution $p_\theta(a_i \mid s_i)$ \emph{est} directement la sortie du
	modèle, ce qui permet de traiter des instances de taille arbitraire.
\end{remark}

\begin{figure}[H]
	\centering
	\begin{tikzpicture}[
		enc/.style={block, fill=ThBlue!12},
		dec/.style={block, fill=ThGreen!12, draw=ThGreen!70!black},
		crit/.style={block, fill=ThRed!10, draw=ThRed},
		node distance=0.9cm and 1.4cm
		]
		\node[enc] (x1) {$x_1$};
		\node[enc, right=of x1] (x2) {$x_2$};
		\node[enc, right=of x2] (xn) {$\dots\ x_n$};
		\node[enc, below=1.1cm of x2, minimum width=5.2cm] (encoder) {Encodeur LSTM};
		\draw[arrow] (x1) -- (encoder.north -| x1);
		\draw[arrow] (x2) -- (encoder.north -| x2);
		\draw[arrow] (xn) -- (encoder.north -| xn);
		
		\node[dec, below=1.4cm of encoder] (decoder) {Décodeur LSTM + Pointer};
		\draw[arrow] (encoder) -- node[right, font=\scriptsize]{$e_1,\dots,e_n$} (decoder);
		
		\node[dec, below left=1.0cm and -0.3cm of decoder] (mask) {Masque\\(villes visitées)};
		\draw[dasharrow] (mask) -- (decoder);
		
		\node[dec, below=1.0cm of decoder] (policy) {$p_\theta(a_i \mid s_i)$};
		\draw[arrow] (decoder) -- (policy);
		
		\node[crit, right=2.0cm of decoder] (critic) {Critic $b_\phi(s)$};
		\draw[arrow] (encoder.east) -- ++(0.8,0) |- (critic.west);
		\node[crit, below=1.0cm of critic] (baseline) {$(b_\phi(s) - L(\pi))^2$};
		\draw[arrow] (critic) -- (baseline);
		
		\node[below=0.6cm of policy, font=\scriptsize\itshape] {tour échantillonné $\pi \sim p_\theta$};
	\end{tikzpicture}
	\caption{Architecture Encodeur--Décodeur du Pointer Network et branche critic associée. L'encodeur produit les représentations $e_1,\dots,e_n$ des villes ; à chaque pas de décodage, le score de compatibilité~\eqref{eq:ptr-score} est calculé entre l'état du décodeur et chaque $e_j$ non masqué, puis normalisé par softmax~\eqref{eq:ptr-softmax} pour former $p_\theta(a_i \mid s_i)$. Le critic, branché sur la même représentation encodée, estime séparément la baseline $b_\phi(s)$ utilisée dans l'équation~\eqref{eq:reinforce-baseline}.}
	\label{fig:ptrnet-archi}
\end{figure}

\subsection{Apprentissage par REINFORCE avec baseline (Actor-Critic)}
\label{sec:ch4-reinforce}

L'entraînement supervisé du Pointer Network original nécessite des tours
optimaux étiquetés, impraticable à grande échelle. Bello \textit{et al.}~\cite{bello2016neural}
proposent un entraînement par renforcement, ne nécessitant que la capacité à
évaluer la longueur d'un tour donné. Par le théorème du gradient de
politique :

\begin{equation}
	\nabla_\theta \mathcal{C}(\theta) = \mathbb{E}\big[(L(\pi) - b_\phi(s))\, \nabla_\theta \log p_\theta(\pi \mid s)\big]
	\label{eq:reinforce-baseline}
\end{equation}

où $\mathcal{C}(\theta)=\mathbb{E}_{\pi\sim p_\theta}[L(\pi)]$ est le coût moyen à minimiser et $b_\phi(s)$ une baseline estimée par le critique, entraîné par régression à minimiser $(b_\phi(s)-L(\pi))^2$. Les paramètres de l'acteur sont mis à jour par descente de gradient, $\theta\leftarrow\theta-\eta\nabla_\theta\mathcal{C}(\theta)$.

\begin{algorithm}[H]
	\KwEntree{instances TSP échantillonnées, taux d'apprentissage $\eta$}
	\KwInit{paramètres actor $\theta$, paramètres critic $\phi$}
	\Pour{chaque batch d'instances}{
		échantillonner un tour $\pi \sim p_\theta(\cdot \mid s)$ \Comment{actor, mode stochastique}\;
		calculer $L(\pi)$ \Comment{longueur réelle du tour}\;
		estimer $b_\phi(s)$ \Comment{critic}\;
		mettre à jour $\theta$ selon l'équation~\eqref{eq:reinforce-baseline}\;
		mettre à jour $\phi$ pour minimiser $(b_\phi(s) - L(\pi))^2$\;
	}
	\caption{Entraînement Actor-Critic / REINFORCE (Bello et al., 2016)}
	\label{alg:actor-critic}
\end{algorithm}

\begin{remark}{Terminologie}{terminologie-ac}
	Il ne s'agit pas d'Actor-Critic au sens classique (avec bootstrap TD), mais
	de REINFORCE avec baseline apprise --- terme souvent désigné par abus de
	langage comme "Actor-Critic" dans la littérature du RL combinatoire.
\end{remark}

% ================================================================
\section{Implémentation}
\label{sec:ch4-implementation}
% ================================================================

\subsection{Environnement et adaptation du dépôt}

Les expériences décrites dans ce chapitre proviennent de deux campagnes qu'il
convient de distinguer. Les premiers essais constituent une campagne pilote,
conservée pour documenter les difficultés rencontrées pendant le développement.
Les conclusions principales reposent plutôt sur l'évaluation indépendante
v1.1, conduite après un entraînement propre et gelé du modèle.

\paragraph{Campagne principale : entraînement propre et évaluation v1.1.}
Le modèle retenu pour l'analyse principale provient du run
\texttt{clean\_restart\_seed20261013\_v1}, initialisé aléatoirement et entraîné
sur des instances de $n=20$ villes. Le checkpoint final
\texttt{train20\_1013\_17\_14\_step12999\_act.pt} correspond à
$13\,000$ mises à jour de l'optimiseur ; le suffixe \texttt{12999} reflète
l'indexation des pas à partir de zéro. Le protocole et les artefacts consignent
le commit \texttt{33636039c37da190ca} et l'empreinte SHA-256 du checkpoint.
Cette campagne a été exécutée sur un GPU Tesla T4 (14,56 Gio), avec Python
3.13.15, PyTorch 2.11+cu130 et CUDA 13.0. Le choix du checkpoint et le test final
frais sont décrits à la section~\ref{sec:eval-v1-1}.

Le protocole de cette campagne documente les adaptations nécessaires à une
réplication contrôlée : fixation de la graine d'entraînement
\texttt{20261013}, neutralisation de l'arrêt anticipé d'origine afin d'effectuer
le budget prévu, et sauvegarde périodique de points de contrôle de l'acteur pour
l'évaluation. Le protocole v1.1 constitue un redémarrage propre ; il ne s'agit
pas d'une continuation des points de contrôle pilotes décrits ci-dessous.

\paragraph{Campagne pilote historique.}
Les premiers essais ont également été exécutés sur Google Colab, avec les
spécifications consignées à l'époque : GPU Tesla T4 (CUDA 12.8), PyTorch
2.11.0+cu128 et Python 3.12.13. Les poids et journaux étaient conservés sur
Google Drive, dans le répertoire \texttt{PFE\_TSP/TSP\_DRL\_PtrNet}. Deux
adaptations propres à cette campagne ont été apportées : (i) ajout de
\texttt{torch.manual\_seed(42)} avant l'instanciation du générateur de données
pour fixer son flux aléatoire ; (ii) ajout dans \texttt{train.py} d'une reprise
d'entraînement chargeant les poids de l'acteur lorsque
\texttt{cfg.act\_model\_path} est renseigné. Ces modifications décrivent les
essais pilotes et ne doivent pas être confondues avec le protocole du run
principal v1.1.

\begin{remark}{Indépendance des jeux d'entraînement et de test}{no-leakage}
Dans la campagne pilote, les instances d'entraînement provenaient du générateur
interne du dépôt, initialisé par \texttt{torch.manual\_seed(42)}, tandis que les
instances d'évaluation étaient construites séparément avec
\texttt{generate\_shared\_instance} et des flux \texttt{numpy.random} dédiés.
Pour l'évaluation principale v1.1, une validation indépendante de 100 instances
à $n=20$ (graine \texttt{7302026}) a servi uniquement à sélectionner le
checkpoint ; le test final frais a ensuite été généré après gel du choix, à
raison de 100 instances pour chacune des tailles $n=20$, $50$ et $100$.
Les graines effectives et la provenance des instances finales sont consignées
dans le manifeste de test, présenté à la section~\ref{sec:eval-v1-1-protocole}.
Les deux protocoles sont donc décrits séparément, sans assimiler les petits jeux
de test pilotes au test final indépendant.
\end{remark}

\begin{remark}{Champ de graine dans la configuration pilote}{cfg-seed}
Le fichier historique \texttt{cfg.pkl} contient un champ \texttt{seed} fixé à
$1$ par défaut dans le dépôt d'origine. Dans la campagne pilote, ce champ ne
pilotait pas le flux utilisé pour générer les données d'entraînement, lequel
était contrôlé par le \texttt{torch.manual\_seed(42)} ajouté manuellement.
Cette précision concerne le fichier de configuration pilote ; elle ne remplace
pas la graine explicitement documentée pour le run principal v1.1.
\end{remark}

\subsection{Architecture et hyperparamètres}

Les hyperparamètres retenus (identiques pour l'\textit{actor} et le
\textit{critic}, sauf mention contraire) sont résumés au
Tableau~\ref{tab:hyperparams}.

\begin{table}[H]
	\centering
	\begin{tabular}{ll}
		\toprule
		Hyperparamètre & Valeur \\
		\midrule
		Dimension d'embedding & 128 \\
		Dimension cachée (LSTM) & 128 \\
		Nombre de glimpses & 1 \\
		Nombre de passes de traitement (\textit{n\_process}) & 3 \\
		Écrêtage des logits (\textit{clip\_logits}) & 10 \\
		Température softmax & 1.0 \\
		Optimiseur & Adam \\
		Taux d'apprentissage initial & 0.001 \\
		Décroissance du taux d'apprentissage & 0.96 tous les 5000 pas \\
		Initialisation des poids & $\mathcal{U}(-0.08,\ 0.08)$ \\
		Mode de décodage (entraînement) & échantillonnage stochastique \\
		Taille de ville de l'instance ($n$) & 20 \\
		Taille de batch & 512 \\
		\bottomrule
	\end{tabular}
	\caption{Hyperparamètres du modèle Pointer Network (actor + critic).}
	\label{tab:hyperparams}
\end{table}

\begin{remark}{Décroissance du taux d'apprentissage et campagne pilote}{lr-decay-noop}
La configuration prévoit une décroissance du taux d'apprentissage par un
facteur $0.96$ tous les $5\,000$ pas. Les points de contrôle pilotes aux pas
200 et 1\,550 (ce dernier après une reprise depuis le pas 1\,999, soit environ
3\,550 pas cumulés) ont été atteints avant le premier palier ; la décroissance
n'a donc pas influencé ces résultats exploratoires. Cette observation ne vaut
pas pour le run principal v1.1, qui a atteint $13\,000$ mises à jour et franchi
les paliers programmés à $5\,000$ et $10\,000$ pas.
\end{remark}

\subsection{Protocole d'entraînement}
\label{sec:ch4-protocole}

Le protocole d'entraînement retenu pour l'analyse principale est celui du run
propre \texttt{clean\_restart\_seed20261013\_v1}. Le modèle a été initialisé
aléatoirement et entraîné uniquement sur des instances de $n=20$ villes, avec
un batch de 512, une dimension cachée de 128, l'optimiseur Adam et un taux
d'apprentissage initial de $0.001$. Le budget réalisé est de
$13\,000$ mises à jour de l'optimiseur. Les snapshots intermédiaires suffixés
\texttt{1000}, \texttt{5000} et \texttt{10000} sont enregistrés après le
nombre correspondant de mises à jour. Le checkpoint final porte le suffixe
\texttt{12999}, car la sauvegarde finale du dépôt utilise l'indice de boucle
commençant à zéro ($i=0,\ldots,12\,999$). Les quatre checkpoints ont été évalués
sur une validation indépendante ; celui au suffixe \texttt{12999} a été retenu
selon la règle pré-spécifiée du minimum de la longueur moyenne en décodage
glouton. Le test final frais a été constitué après cette sélection et n'a pas
servi au choix du checkpoint (détails en section~\ref{sec:eval-v1-1}).

\begin{remark}{Historique des entraînements pilotes}{entrainement-partiel}
Avant ce redémarrage propre, une première campagne avait comporté un entraînement
initial de 2\,000 pas, puis des reprises incomplètes. Un essai s'est arrêté au
point de contrôle du pas 200 de la reprise
(\texttt{train20\_0811\_08\_13\_step200\_act.pt}, pertes actor/critic
approximativement $-4.3/8.3$). Un autre essai, repris depuis le point de départ
\texttt{step1999}, s'est arrêté au pas 1\,550 de reprise
(\texttt{train20\_0815\_14\_27\_step1550\_act.pt}), soit environ 3\,550 pas
cumulés, avec des pertes stabilisées mais sans preuve de convergence globale.
Ces checkpoints sont conservés dans les tableaux exploratoires pour retracer
les essais initiaux ; ils ne constituent pas le modèle retenu pour l'évaluation
indépendante finale. Le poids \texttt{train20\_1113\_12\_12\_step14999\_act.pt}
fourni avec le dépôt est un modèle externe, et non un modèle entraîné dans ce
PFE.
\end{remark}

\begin{figure}[H]
\centering
\includegraphics[width=\textwidth]{images/convergence_step1550.pdf}
\caption{Courbes d'apprentissage d'un essai pilote interrompu au pas 1\,550 de
reprise (depuis le pas 1\,999). La perte de l'Actor passe de $-12{,}68$ à
$-0{,}14$ : sa valeur numérique augmente, tandis que son amplitude absolue
diminue. La perte du Critic diminue de $16{,}55$ à $0{,}30$. La moyenne cumulée
des longueurs des tournées d'entraînement passe de $4{,}2757$ à $4{,}1433$, soit
une baisse d'environ $3{,}1~\%$ ; son minimum observé dans les journaux apparaît
autour du pas 1\,530. Ces métriques décrivent cet essai historique et ne
constituent ni une évaluation indépendante de la généralisation ni une preuve
de convergence globale. Elles ne concernent pas le run principal v1.1.}
\label{fig:convergence-step1550}
\end{figure}

% ================================================================
\section{Résultats expérimentaux}
\label{sec:ch4-resultats}
% ================================================================


\subsection{Évaluation indépendante du modèle entraîné : protocole v1.1}
\label{sec:eval-v1-1}

Cette sous-section présente l'évaluation indépendante \textbf{v1.1}, conduite \emph{après gel de l'entraînement} du modèle \texttt{clean\_restart\_seed20261013\_v1} (aucun ré-entraînement). Elle lève les limites du pilote identifiées dans \path{research/review/TSP_DRL_Audit_Review.md} : validation séparée, test final frais, protocole de temps standardisé et provenance complète. Les résultats sont descriptifs pour \emph{ce seul seed d'entraînement} (\texttt{20261013}) et ne quantifient pas la variance inter-seeds. Les résultats numériques des expériences antérieures (checkpoints 200 et 1\,550, ainsi que poids externe au pas $14\,999$) sont conservés à titre historique et présentés séparément.

\subsubsection{Protocole expérimental et plan de test indépendant}
\label{sec:eval-v1-1-protocole}

\paragraph{Modèle évalué.} Checkpoint final \texttt{train20\_1013\_17\_14\_step12999\_act.pt}, soit \textbf{13~000 mises à jour d'optimiseur} (\emph{suffixe fichier 12999} dû à l'indexation $i=0\ldots12999$ ; empreinte \texttt{SHA-256} \path{87610c76c1239c32d2bdba58...}, commit \path{33636039c37da190ca}). Le modèle a été entraîné uniquement à $n=20$, \texttt{batch} 512, \texttt{hidden} 128, Adam \texttt{lr} $0{,}001$.

\paragraph{Jeux de données.} Un jeu de \textbf{validation indépendant} à $n=20$ (100 instances, graine \texttt{7302026}, coordonnées uniformes $[0,1]^2$, stockage \texttt{float32}) a servi \emph{uniquement} à choisir le checkpoint selon la règle pré-spécifiée du minimum de la longueur moyenne en décodage \emph{glouton}. Un jeu de \textbf{test final frais}, généré \emph{après gel} de ce choix, couvre $n=20$, $50$ et $100$ à raison de 100 instances / taille (\texttt{graine de base} $8302026$, règle \texttt{effective\_seed} $=$ \texttt{base}+$n$, soit \texttt{8302046} / \texttt{8302076} / \texttt{8302126} ; cf. \path{final_test_instances/final_test_instances_manifest.json}). Les 100 instances d'une même taille sont \textbf{appariées} entre méthodes (même \texttt{instance\_id}). La colonne \texttt{seed} de \texttt{final\_test\_per\_instance.csv} enregistre des \textbf{graines méthode-spécifiques} (p. ex.\ $n=20$, \texttt{iid} 0 : $10302026$ pour classiques, $19302026$ pour DRL) : \emph{elle ne doit pas être interprétée comme graine de génération des coordonnées}, dont la source de vérité reste le manifest.

\paragraph{Méthodes comparées sur test final.} Cinq configurations évaluées sur chaque instance : \texttt{NN\_start0} (plus proche voisin), \texttt{NN\_start0\_plus\_2opt} (PPV + recherche locale 2-opt jusqu'à optimum local), \texttt{ACO\_fixed\_parameters} (colonies de fourmis, paramètres fixes : 15 fourmis, 30 itérations, $\alpha=1$, $\beta=3$, $\rho=0{,}1$, $q=1$), \texttt{DRL\_greedy} (une passe gloutonne) et \texttt{DRL\_best\_of\_512} (meilleur-sur-512). Soit $1\,500$ tournées au total ($100 \times 5 \times 3$ tailles). Le gap est calculé \emph{par instance} contre le \textbf{meilleur observé parmi les cinq méthodes comparées} :
\[
\textit{gap}_{m}(i)=100\cdot\left(\frac{L_{m}(i)}{\min_{k}L_{k}(i)}-1\right)\;[\%].
\]
Un gap de $0\,\%$ signifie \og meilleur des cinq ici \fg{}, non optimal. À $n=100$, \texttt{NN+2-opt} obtient le gap moyen $0\,\%$ sur $100/100$ instances \emph{par construction} --- ce qui ne prouve pas l'optimalité exacte (aucun solveur exact n'a été exécuté). Voir \path{research/final_evaluation_v1_1/README_DATA.md} \S\S~5--6.

\paragraph{Inférence statistique.} Différences appariées $\Delta_{m,b}(i)=L_{m}(i)-L_{b}(i)$ pour $m\in\{\text{DRL}\}\times b\in\{\text{classiques}\}$ sur les $100$ instances appariées / taille, moyenne appariée et \textbf{intervalle de confiance bootstrap apparié à 95~\%} par ré-échantillonnage non paramétrique des instances ($B=10\,000$ répliques, \textbf{graine RNG} $8802026$, percentile $2{,}5$--$97{,}5$), sans correction de multiplicité (18 comparaisons). Un intervalle croisant zéro décrit l'absence de différence moyenne claire à 95~\%, non une preuve d'égalité (voir \S~4).

\paragraph{Temps de résolution.} Mesures \textbf{répétées} sur un sous-ensemble de $10$ instances / taille, $1$ échauffement + $5$ répétitions $\to$ $50$ mesures / $(n,\text{méthode})$, $750$ au total (\texttt{runtime\_profile\_repeated.csv}). Chargement du modèle exclu (synchronisation CUDA). Pour le DRL : transfert entrée + forward + score des candidats + récupération résultat ; pour CPU : exécution solveur + calcul longueur. Matériel : Tesla~T4 (14,56~Gio), Python~3.13.15, PyTorch~2.11+cu130, CUDA~13.0. Les temps mono-mesure figurent aussi dans \texttt{final\_test\_per\_instance.csv} mais la synthèse utilise les mesures répétées.

\subsubsection{Sélection du point de contrôle par validation indépendante}
\label{sec:eval-v1-1-validation}

\begin{table}[H]
\centering
\small
\begin{tabular}{rccc}
\toprule
Suffixe fichier & Mises à jour & Longueur moyenne & Écart-type ($\sigma$) \\
 & effectuées & (validation $n=20$) & (échantillon) \\
\midrule
1000  &  1\,000 & 7,595 & 1,180 \\
5000  &  5\,000 & 4,126 & 0,326 \\
10000 & 10\,000 & 4,037 & 0,346 \\
\textbf{12999} & \textbf{13\,000} & \textbf{4,016} & \textbf{0,321} \\
\bottomrule
\end{tabular}
\caption[Sélection du point de contrôle par validation indépendante]{Sélection du point de contrôle par validation indépendante ($n=20$, 100 instances, graine $7302026$, décodage glouton uniquement). La règle pré-spécifiée est le minimum de la longueur moyenne ; retenu : \texttt{train20\_1013\_17\_14\_step12999\_act.pt} --- 13\,000 mises à jour (\emph{suffixe 12999} dû à l'index $i=0\ldots12999$). Source : \texttt{validation/checkpoint\_selection\_validation\_summary.csv}.}
\label{tab:validation-selection-v1-1}
\end{table}

\begin{figure}[H]
\centering
\includegraphics[width=0.78\textwidth]{images/validation_selection_v1_1.pdf}
\caption{Reconstitution graphique à partir des moyennes et écarts-types arrondis du Tableau~\ref{tab:validation-selection-v1-1} : sélection par validation indépendante : longueur moyenne gloutonne ($\pm\sigma$ échantillon) selon le checkpoint. Le minimum à 13\,000 mises à jour justifie le choix du checkpoint final. Le gap n'est pas représenté ici (validation sans référence multi-méthode).}
\label{fig:validation-selection-v1-1}
\end{figure}

La Figure~\ref{fig:validation-selection-v1-1} illustre la décroissance monotone de la longueur moyenne gloutonne avec le nombre de mises à jour sur la validation indépendante, le minimum étant atteint au checkpoint final retenu pour le test final frais.

\subsubsection{Résultats sur test final frais ($n=20$, $50$, $100$)}
\label{sec:eval-v1-1-resultats}

\begin{table}[H]
\centering
\small
\begin{tabular}{llccccc}
\toprule
$n$ & Méthode & \multicolumn{2}{c}{Longueur de tournée} & Médiane & \multicolumn{2}{c}{Gap au meilleur observé (\%)} \\
 & & Moyenne & $\sigma$ (éch.) & & Moyenne & $\sigma$ (éch.) \\
\midrule
\multirow{5}{*}{$20$}
 & ACO (param.~fixes) & 3,878 & 0,273 & 3,876 & 0,34 & 0,76 \\
 & PPV + 2-opt & 3,942 & 0,300 & 3,935 & 2,01 & 3,02 \\
 & DRL meilleur-sur-512 & 3,959 & 0,307 & 3,983 & 2,43 & 2,93 \\
 & DRL glouton & 4,037 & 0,325 & 4,028 & 4,44 & 3,97 \\
 & PPV & 4,496 & 0,473 & 4,503 & 16,42 & 10,40 \\
\midrule
\multirow{5}{*}{$50$}
 & PPV + 2-opt & 5,971 & 0,293 & 5,945 & 0,60 & 1,33 \\
 & ACO (param.~fixes) & 6,157 & 0,315 & 6,151 & 3,75 & 3,17 \\
 & DRL meilleur-sur-512 & 6,347 & 0,361 & 6,339 & 6,97 & 4,80 \\
 & DRL glouton & 6,672 & 0,437 & 6,638 & 12,46 & 6,25 \\
 & PPV & 6,944 & 0,533 & 6,900 & 17,05 & 8,04 \\
\midrule
\multirow{5}{*}{$100$}
 & PPV + 2-opt & 8,271 & 0,304 & 8,304 & 0,00 & 0,00 \\
 & ACO (param.~fixes) & 9,161 & 0,371 & 9,100 & 10,81 & 3,76 \\
 & PPV & 9,739 & 0,530 & 9,717 & 17,81 & 5,92 \\
 & DRL meilleur-sur-512 & 9,928 & 0,392 & 9,939 & 20,13 & 4,80 \\
 & DRL glouton & 10,468 & 0,460 & 10,470 & 26,66 & 5,70 \\
\bottomrule
\end{tabular}
\caption[Résultats sur test final frais v1.1]{Test final frais v1.1 (100 instances / taille, graines effectives $8302046$ / $8302076$ / $8302126$) du checkpoint final (13\,000 mises à jour, suffixe \texttt{12999}). Longueur euclidienne moyenne $\pm$ écart-type \emph{échantillon} (ddof=1), médiane, et gap moyen au meilleur observé parmi les cinq méthodes (non optimum exact) $\pm\sigma$. DRL : \texttt{DRL\_greedy} (1 passe) et \texttt{DRL\_best\_of\_512}. Source : \texttt{final\_test\_summary.csv} ; les valeurs reproduisent les CSV à $10^{-12}$ (moyenne) et $6{\times}10^{-16}$ ($\sigma$). Méthodes triées par longueur moyenne croissante au sein de chaque taille.}
\label{tab:resultats-final-v1-1}
\end{table}

\begin{figure}[H]
\centering
\includegraphics[width=0.95\textwidth]{images/tour_lengths_final_test_v1_1.pdf}
\caption{Reconstitution graphique à partir des moyennes et écarts-types arrondis du Tableau~\ref{tab:resultats-final-v1-1} : longueurs moyennes $\pm\sigma$ (échantillon) sur test final frais v1.1 par taille et méthode. Checkpoint final 13\,000 mises à jour (suffixe 12999). Les barres d'erreur sont des $\sigma$ échantillon, non des intervalles de confiance de la moyenne.}
\label{fig:tour-lengths-final-v1-1}
\end{figure}

\begin{figure}[H]
\centering
\includegraphics[width=0.95\textwidth]{images/gaps_final_test_v1_1.pdf}
\caption{Reconstitution graphique à partir des statistiques arrondies du Tableau~\ref{tab:resultats-final-v1-1} : gaps moyens (\%) au meilleur observé parmi les cinq méthodes sur test final frais v1.1. Un gap de $0\,\%$ signifie \og meilleur des cinq ici \fg{} (non optimum exact). À $n=100$, PPV+2-opt obtient 0~\% sur 100/100 instances \emph{par construction}.}
\label{fig:gaps-final-v1-1}
\end{figure}

À $n=20$ sur ce test final apparié (100 instances, graines $8302046$), l'ACO obtient la longueur moyenne la plus faible ($3{,}878\pm0{,}273$), suivi de PPV+2-opt ($3{,}942\pm0{,}300$) puis du DRL meilleur-sur-512 ($3{,}959\pm0{,}307$) et du DRL glouton ($4{,}037\pm0{,}325$), PPV seul étant le plus long ($4{,}496\pm0{,}473$). L'écart DRL meilleur-sur-512 vs.\ PPV+2-opt est examiné de façon appariée au \S~suivant : l'intervalle à 95~\% croise zéro et ne permet pas d'établir une différence moyenne claire.

À $n=50$ (graine $8302076$), les deux modes DRL améliorent le PPV seul en longueur moyenne mais restent supérieurs (plus longs) à l'ACO ($6{,}157\pm0{,}315$) et surtout à PPV+2-opt ($5{,}971\pm0{,}293$, meilleur observé). À $n=100$ (graine $8302126$), les deux modes DRL sont supérieurs à \emph{toutes} les trois références classiques ; PPV+2-opt domine nettement ($8{,}271\pm0{,}304$) et représente le minimum observé sur $100/100$ instances --- ce qui, en l'absence d'optimum exact, ne prouve pas l'optimalité mais constitue le meilleur des cinq ici.

\subsubsection{Comparaisons appariées et incertitude}
\label{sec:eval-v1-1-paires}

\begin{table}[H]
\centering
\scriptsize
\begin{tabular}{cllccc}
\toprule
$n$ & DRL & Référence & $\Delta = L_{\text{DRL}}-L_{\text{ref}}$ & IC 95~\% bootstrap & Croise $0$ ? \\
 & & & moyenne appariée & (percentile, $B{=}10\,000$, RNG $8802026$) & \\
\midrule
\multirow{6}{*}{$20$}
 & glouton & PPV & $-0{,}459591$ & $[-0{,}536047\,;\,-0{,}383334]$ & non \\
 & glouton & PPV+2-opt & $+0{,}094120$ & $[+0{,}055884\,;\,+0{,}134430]$ & non \\
 & glouton & ACO & $+0{,}158482$ & $[+0{,}129191\,;\,+0{,}190806]$ & non \\
 & meilleur-sur-512 & PPV & $-0{,}536787$ & $[-0{,}613731\,;\,-0{,}459909]$ & non \\
 & meilleur-sur-512 & PPV+2-opt & $+0{,}016923$ & $[-0{,}017012\,;\,+0{,}049661]$ & \textbf{oui} \\
 & meilleur-sur-512 & ACO & $+0{,}081285$ & $[+0{,}059004\,;\,+0{,}104979]$ & non \\
\midrule
\multirow{6}{*}{$50$}
 & glouton & PPV & $-0{,}272310$ & $[-0{,}385350\,;\,-0{,}161928]$ & non \\
 & glouton & PPV+2-opt & $+0{,}700817$ & $[+0{,}628162\,;\,+0{,}776003]$ & non \\
 & glouton & ACO & $+0{,}515168$ & $[+0{,}441707\,;\,+0{,}589007]$ & non \\
 & meilleur-sur-512 & PPV & $-0{,}597777$ & $[-0{,}702518\,;\,-0{,}493301]$ & non \\
 & meilleur-sur-512 & PPV+2-opt & $+0{,}375350$ & $[+0{,}319424\,;\,+0{,}431710]$ & non \\
 & meilleur-sur-512 & ACO & $+0{,}189701$ & $[+0{,}131827\,;\,+0{,}248989]$ & non \\
\midrule
\multirow{6}{*}{$100$}
 & glouton & PPV & $+0{,}729301$ & $[+0{,}614851\,;\,+0{,}844684]$ & non \\
 & glouton & PPV+2-opt & $+2{,}197607$ & $[+2{,}111417\,;\,+2{,}282387]$ & non \\
 & glouton & ACO & $+1{,}307510$ & $[+1{,}214654\,;\,+1{,}404016]$ & non \\
 & meilleur-sur-512 & PPV & $+0{,}189396$ & $[+0{,}082336\,;\,+0{,}296867]$ & non \\
 & meilleur-sur-512 & PPV+2-opt & $+1{,}657701$ & $[+1{,}583517\,;\,+1{,}731271]$ & non \\
 & meilleur-sur-512 & ACO & $+0{,}767604$ & $[+0{,}686702\,;\,+0{,}851001]$ & non \\
\bottomrule
\end{tabular}
\caption[Comparaisons appariées et intervalles de confiance bootstrap]{Comparaisons appariées : différences moyennes $L_{\text{DRL}}-L_{\text{référence}}$ sur $100$ instances appariées / taille et intervalles de confiance bootstrap appariés à 95~\% ($B{=}10\,000$, graine $8802026$, sans correction de multiplicité ; 18 comparaisons). Valeur négative favorise le DRL (tour plus courte). Un intervalle croisant zéro décrit l'absence de différence moyenne claire à 95~\%, non une preuve d'égalité. Source : \texttt{paired\_comparisons\_bootstrap.csv} (moyennes recalculées $\le 1{,}4{\times}10^{-17}$).}
\label{tab:comparaisons-appariees-v1-1}
\end{table}

La convention de signe est $L_{\text{DRL}}-L_{\text{référence}}$ : le négatif favorise le DRL (tour plus courte). \emph{À $n=20$}, le DRL glouton améliore le PPV seul mais perd face à PPV+2-opt et ACO ; le DRL meilleur-sur-512 améliore le PPV seul, \textbf{ne se distingue pas clairement de PPV+2-opt} ($+0{,}016923$, IC $[-0{,}017012\,;\,+0{,}049661]$ \textbf{croisant zéro}, donc pas de différence moyenne claire à 95~\% sans constituer une preuve d'égalité), et perd face à l'ACO. \emph{À $n=50$}, les deux modes DRL améliorent le PPV seul (glouton $-0{,}272310$ IC $[-0{,}385350\,;\,-0{,}161928]$, meilleur-sur-512 $-0{,}597777$ IC $[-0{,}702518\,;\,-0{,}493301]$) mais perdent face à PPV+2-opt et ACO. \emph{À $n=100$}, les deux modes DRL perdent face aux \emph{trois} références classiques : les six IC sont strictement positifs (p. ex.\ meilleur-sur-512 vs.\ PPV+2-opt $+1{,}657701$ IC $[+1{,}583517\,;\,+1{,}731271]$, vs.\ ACO $+0{,}767604$ IC $[+0{,}686702\,;\,+0{,}851001]$, et même vs.\ PPV seul $+0{,}189396$ IC $[+0{,}082336\,;\,+0{,}296867]$).

\subsubsection{Résultats de temps de résolution}
\label{sec:eval-v1-1-temps}

\begin{table}[H]
\centering
\scriptsize
\begin{tabular}{llccccc}
\toprule
$n$ & Méthode & $N$ inst. & $M$ mesures & Moyenne (s) & $\sigma$ (éch.) & Médiane (s) \\
\midrule
\multirow{5}{*}{$20$}
 & PPV & 10 & 50 & 0,00013 & 0,00006 & 0,00012 \\
 & PPV+2-opt & 10 & 50 & 0,00107 & 0,00036 & 0,00097 \\
 & DRL glouton & 10 & 50 & 0,02422 & 0,00165 & 0,02373 \\
 & DRL meilleur-sur-512 & 10 & 50 & 0,05499 & 0,00855 & 0,05302 \\
 & ACO (param.~fixes) & 10 & 50 & 0,60588 & 0,02287 & 0,60301 \\
\midrule
\multirow{5}{*}{$50$}
 & PPV & 10 & 50 & 0,00047 & 0,00007 & 0,00044 \\
 & PPV+2-opt & 10 & 50 & 0,01660 & 0,00631 & 0,01487 \\
 & DRL glouton & 10 & 50 & 0,05767 & 0,00379 & 0,05611 \\
 & DRL meilleur-sur-512 & 10 & 50 & 0,18584 & 0,01480 & 0,18384 \\
 & ACO (param.~fixes) & 10 & 50 & 1,59835 & 0,14994 & 1,56377 \\
\midrule
\multirow{5}{*}{$100$}
 & PPV & 10 & 50 & 0,00197 & 0,00055 & 0,00176 \\
 & PPV+2-opt & 10 & 50 & 0,13877 & 0,04145 & 0,13114 \\
 & DRL glouton & 10 & 50 & 0,12009 & 0,00914 & 0,11637 \\
 & DRL meilleur-sur-512 & 10 & 50 & 0,58797 & 0,00907 & 0,58888 \\
 & ACO (param.~fixes) & 10 & 50 & 3,21503 & 0,10536 & 3,21214 \\
\bottomrule
\end{tabular}
\caption[Temps de résolution : protocole répété v1.1]{Temps de résolution : protocole répété v1.1 ($1$ échauffement + $5$ répétitions, $10$ instances / taille, soit $50$ mesures / configuration ; $750$ au total). Chargement du modèle exclu. DRL : transfert entrée + inférence forward + score des candidats + récupération résultat (synchronisé CUDA). CPU : exécution solveur + calcul longueur. Matériel : Tesla~T4 (14,56~Gio). Les moyennes recalculées depuis \texttt{runtime\_profile\_repeated.csv} coïncident à $<10^{-12}$.}
\label{tab:temps-v1-1}
\end{table}

\begin{figure}[H]
\centering
\includegraphics[width=0.92\textwidth]{images/runtime_final_test_v1_1.pdf}
\caption{Reconstitution graphique des moyennes et écarts-types arrondis du Tableau~\ref{tab:temps-v1-1} : temps moyens de résolution $\pm\sigma$ (échelle logarithmique) --- protocole répété v1.1 : 10 instances / taille $\times$ 5 répétitions (1 échauffement). Le DRL glouton ($\sim 0{,}02$--$0{,}12$~s) est plus lent que PPV+2-opt aux petites tailles mais plus rapide que l'ACO (0,6--3,2~s) ; le meilleur-sur-512 double à triple ce coût. Les barres représentent des moyennes et écarts-types, non les distributions des mesures individuelles. Chargement exclu ; frontières DRL/CPU non identiques : il s'agit du \emph{temps mesuré solveur/passe avant} de cette implémentation, non d'une latence end-to-end.}
\label{fig:runtime-final-v1-1}
\end{figure}

Les temps médians (50 mesures / configuration) placent le DRL glouton à $\sim 0{,}024$~s ($n=20$), $0{,}058$~s ($n=50$), $0{,}120$~s ($n=100$) et le DRL meilleur-sur-512 à $0{,}055$--$0{,}186$--$0{,}588$~s, à comparer à PPV+2-opt ($0{,}001$--$0{,}017$--$0{,}139$~s) et ACO ($0{,}606$--$1{,}598$--$3{,}215$~s). Ces mesures décrivent le \emph{temps solveur/passe avant} de cette implémentation, non une latence end-to-end (chargement exclu, frontières de chronométrage non identiques).

\subsubsection{Interprétation et limites}
\label{sec:eval-v1-1-limites}

Cette évaluation indépendante confirme et précise les lectures du pilote : \textbf{à $n=20$} l'ACO reste la meilleure longueur moyenne parmi les cinq configurations testées sur ce test frais ; \textbf{DRL meilleur-sur-512 et PPV+2-opt} ne se distinguent pas clairement en moyenne ($+0{,}016923$ IC croisant zéro). \textbf{À $n=50$}, le DRL améliore le PPV seul mais \emph{perd} face à PPV+2-opt et ACO ; \textbf{à $n=100$}, les deux modes DRL perdent face aux \emph{trois} références classiques, PPV+2-opt dominant nettement (meilleur sur $100/100$ instances). Il s'agit d'une dégradation de la généralisation pour \emph{cette configuration d'entraînement} (uniquement à $n=20$, $13\,000$ mises à jour, validation à $n=20$ seul), non d'une conclusion sur le DRL en général : attribuer ces écarts aux limites du protocole actuel plutôt qu'à une borne théorique.

\paragraph{Portée historique des limites v1.1.}
Les limites ci-dessous décrivent l'état de la première évaluation, conduite avec
la seule graine d'entraînement $20261013$ ; elles ne doivent pas être lues comme
la description du protocole final après extension. L'extension v1.2 présentée
à la sous-section suivante ajoute deux entraînements indépendants et porte à
54 le nombre de comparaisons bootstrap conditionnelles. Elle répond donc
partiellement à la limite initiale relative à l'absence de répétitions
inter-graines, sans supprimer l'incertitude liée au faible nombre de trois
entraînements. Restent valables pour les deux campagnes : (1) la sélection du
checkpoint repose uniquement sur une validation gloutonne à $n=20$, sans
validation ni adaptation à $n=50$/$100$ ; (2) dans v1.1, les gaps étaient
calculés par rapport au meilleur observé parmi cinq configurations, et non à
l'optimum exact ; l'extension v1.2 redéfinit le référentiel sur les neuf
configurations comparées, sans qu'aucune résolution exacte soit effectuée ;
(3) les intervalles bootstrap sont descriptifs, conditionnels au modèle et
sans correction de multiplicité ; un intervalle croisant zéro ne prouve pas
l'égalité ; (4) les temps reposent sur un seul profil répété v1.1, avec un
échauffement, cinq répétitions et dix instances par taille, chargement du
modèle exclu et frontières DRL/CPU différentes. Ils caractérisent donc le
temps mesuré du solveur/de la passe avant de cette implémentation, non une
latence bout-en-bout, et ne mesurent aucune dispersion inter-graines des
temps. Enfin, la colonne \texttt{seed} de
\texttt{final\_test\_per\_instance.csv} contient des graines
méthode-spécifiques ; elle ne doit pas être confondue avec les graines
utilisées pour générer les coordonnées des instances.

% ================================================================


\subsection{Robustesse inter-graines : extension expérimentale v1.2}
\label{sec:eval-multiseed}

\paragraph{Protocole et indépendance des entraînements.}
L'évaluation v1.1 précédente portait sur le seul entraînement de graine
\texttt{20261013}. Une extension v1.2 ajoute deux entraînements initialisés
indépendamment, de graines \texttt{20261014} et \texttt{20261015}, sans
réutilisation des poids d'un autre run. Les trois exécutions reposent sur le
même commit du dépôt (\texttt{33636039c37da190cadb5bc136162c1e3f4daf46}) et
le même contrat expérimental : instances d'entraînement à $n=20$, batch de
512, dimensions d'embedding et d'état caché égales à 128, Adam avec un taux
initial de $0{,}001$, même calendrier de taux d'apprentissage et
$13\,000$ mises à jour de l'optimiseur. Les trois modèles ont donc été
entraînés à la même taille, puis testés sans adaptation à $n=50$ et $n=100$.

Pour chacun des trois entraînements, le checkpoint est choisi séparément
parmi les suffixes de fichier \texttt{1000}, \texttt{5000}, \texttt{10000}
et \texttt{12999}, en minimisant la longueur moyenne gloutonne sur le même
jeu de validation indépendant de 100 instances à $n=20$ (graine
\texttt{7302026}). Le checkpoint retenu est le suffixe \texttt{12999}, soit
$13\,000$ mises à jour, pour les trois graines. Les longueurs moyennes de
validation correspondantes sont respectivement $4{,}0163$, $4{,}0453$ et
$4{,}0582$. Les fichiers de sélection et leurs empreintes sont consignés dans
\texttt{multiseed\_selected\_checkpoints.csv} et dans les manifests propres
à chaque évaluation.

Le test final est strictement commun aux trois modèles : 100 instances
inédites pour chacune des tailles $n\in\{20,50,100\}$, avec la règle de
graine effective $8302026+n$ (soit $8302046$, $8302076$ et $8302126$).
L'identité des coordonnées entre les trois évaluations a été vérifiée par
empreintes SHA-256. Les trois baselines classiques (PPV, PPV+2-opt et ACO)
ont également été vérifiées comme identiques entre les évaluations et ne sont
comptées qu'une fois dans le jeu agrégé. Pour chaque instance, le tableau
combiné contient ainsi six résultats DRL (trois modèles fois deux modes de
décodage) et trois résultats classiques, soit neuf observations par instance
et $2\,700$ lignes au total.

\paragraph{Règle d'agrégation.}
Pour chaque taille et chaque mode DRL, on calcule d'abord la longueur moyenne
sur les 100 instances pour chacun des trois modèles. La moyenne rapportée
ci-dessous est la moyenne de ces trois moyennes par modèle ; l'écart-type
associé est l'écart-type empirique d'échantillon entre les trois moyennes
(\texttt{ddof}=1). Il ne s'agit ni de l'écart-type entre les 300 résultats
modèle-instance ni d'un intervalle de confiance. Les baselines étant uniques
et communes aux trois évaluations, leur moyenne et leur écart-type sont
calculés sur les 100 instances appariées. Ces deux types de dispersion sont
présentés séparément afin de ne pas confondre variabilité inter-graines et
variabilité inter-instances. Les résultats proviennent de
\path{multiseed_across_seed_summary.csv} et
\path{multiseed_per_seed_summary.csv}.

\begin{table}[H]
\centering
\small
\begin{tabular}{crrr}
\toprule
$n$ & Configuration DRL & Moyenne des trois moyennes & Écart-type inter-graines \\
\midrule
20  & DRL glouton          & 4,0355  & 0,0052 \\
20  & DRL meilleur-sur-512 & 3,9584  & 0,0054 \\
50  & DRL glouton          & 6,6403  & 0,0281 \\
50  & DRL meilleur-sur-512 & 6,3294  & 0,0209 \\
100 & DRL glouton          & 10,4373 & 0,0364 \\
100 & DRL meilleur-sur-512 & 9,8946  & 0,0293 \\
\bottomrule
\end{tabular}
\caption{Résultats DRL agrégés sur trois graines. La première statistique est la moyenne des trois longueurs moyennes par modèle ; la seconde est l'écart-type d'échantillon entre ces trois moyennes.}
\label{tab:multiseed-drl}
\end{table}

\begin{table}[H]
\centering
\small
\begin{tabular}{crrr}
\toprule
$n$ & Méthode classique & Moyenne sur 100 instances & Écart-type entre instances \\
\midrule
20  & PPV          & 4,4961 & 0,4729 \\
20  & PPV+2-opt    & 3,9424 & 0,3003 \\
20  & ACO          & 3,8780 & 0,2735 \\
\midrule
50  & PPV          & 6,9443 & 0,5326 \\
50  & PPV+2-opt    & 5,9712 & 0,2932 \\
50  & ACO          & 6,1568 & 0,3150 \\
\midrule
100 & PPV          & 9,7388 & 0,5301 \\
100 & PPV+2-opt    & 8,2705 & 0,3038 \\
100 & ACO          & 9,1606 & 0,3711 \\
\bottomrule
\end{tabular}
\caption{Résultats des baselines classiques sur le jeu de test commun. Contrairement au tableau~\ref{tab:multiseed-drl}, l'écart-type est calculé entre les 100 instances et non entre graines d'entraînement.}
\label{tab:multiseed-baselines}
\end{table}

\paragraph{Comparaisons appariées et lecture des résultats.}
Pour chaque graine, chaque taille et chacun des deux modes DRL, les différences
$L_{\mathrm{DRL}}-L_{\text{référence}}$ sont calculées sur les mêmes 100
instances pour les trois baselines. Les 54 intervalles bootstrap à 95~\%
(3 graines $\times$ 3 tailles $\times$ 2 modes DRL $\times$ 3 références)
utilisent 10\,000 ré-échantillonnages chacun. Ils caractérisent l'incertitude
liée aux instances conditionnellement à chaque modèle fixé ; ils ne sont pas
fusionnés en un échantillon artificiel de 300 instances et ne quantifient pas
la variabilité de l'entraînement. Aucune correction de multiplicité n'est
appliquée. Les 54 différences et leurs bornes sont données en
Annexe~\ref{app:multiseed-bootstrap}.

À $n=20$, ACO obtient la plus faible moyenne ($3{,}8780$), suivi de PPV+2-opt
($3{,}9424$) et de DRL meilleur-sur-512 ($3{,}9584$). Pour ce dernier mode,
les différences appariées avec PPV+2-opt sont positives pour les trois
modèles ($+0{,}016923$, $+0{,}010224$ et $+0{,}020807$), mais chacun des trois
intervalles bootstrap correspondants croise zéro. On ne met donc pas en
évidence de différence moyenne claire à 95~\% dans ces comparaisons
conditionnelles, sans démontrer l'équivalence des méthodes. En revanche, les
trois intervalles comparant DRL meilleur-sur-512 à ACO sont strictement
positifs : ACO obtient une longueur moyenne plus faible dans ces essais.

À $n=50$, les deux modes DRL restent meilleurs que PPV seul en longueur
moyenne ; les intervalles appariés conditionnels aux trois modèles sont tous
strictement négatifs pour ces comparaisons. Ils restent toutefois derrière
PPV+2-opt et ACO : les intervalles correspondants sont strictement positifs
pour les trois graines. À $n=100$, les deux modes DRL sont dépassés par les
trois baselines pour chacun des modèles ; les 18 intervalles associés sont
strictement positifs. PPV+2-opt obtient la longueur moyenne la plus faible
($8{,}2705$) et le meilleur coût observé sur les 100 instances, par
construction du référentiel empirique à neuf configurations.

Le mode meilleur-sur-512 réduit la longueur moyenne par rapport au décodage
glouton pour les trois tailles : l'amélioration moyenne est d'environ
$0{,}0771$, $0{,}3109$ et $0{,}5428$ pour $n=20$, $50$ et $100$. Cette
amélioration s'accompagne d'un budget d'inférence plus élevé, comme le montre
le profil de temps répété v1.1 ; ce profil n'a pas été remesuré pour les deux
entraînements supplémentaires et ne permet donc pas d'estimer la variabilité
des temps entre graines.

\paragraph{Environnement logiciel.}
Les trois évaluations utilisent PyTorch \texttt{2.11.0+cu130}, CUDA
\texttt{13.0} et un GPU Tesla~T4. La version mineure corrective de Python
n'est pas strictement identique : \texttt{3.13.15} pour le premier run et
\texttt{3.13.16} pour les deux suivants. Cette différence est documentée
comme une limite de reproductibilité logicielle ; les versions de PyTorch,
CUDA et le matériel d'inférence sont, eux, identiques. Le profil de temps
présenté dans la sous-section v1.1 est réutilisé à l'identique et reste
spécifique à cette première campagne.


\subsection{Protocole de comparaison des essais pilotes}

Les tableaux historiques qui suivent documentent les premiers essais sur
$10$ instances par configuration. Ils servent à retracer le développement du
modèle, mais ne constituent pas le protocole comparatif principal ; celui-ci
est l'évaluation indépendante v1.1 présentée dans les sous-sections précédentes.

Pour les essais pilotes, le Pointer Network a été évalué selon deux modes :
\begin{itemize}
\item \textbf{une passe stochastique} : une seule tournée générée par instance
avec \texttt{decode\_type=sampling} ; il ne s'agit donc pas d'un décodage glouton ;
\item \textbf{échantillonnage meilleur-sur-512} : 512 tournées candidates
échantillonnées par instance, dont la plus courte est conservée.
\end{itemize}

Dans ces tableaux pilotes, le \emph{gap} est calculé par rapport au meilleur
coût observé parmi les méthodes comparées sur l'instance considérée, et non par
rapport à l'optimum exact (Bellman--Held--Karp devient rapidement coûteux,
cf. Annexe~\ref{ann:bhk}).

\begin{remark}{Nature du gap rapporté}{nature-gap}
Un gap de $0\,\%$ signifie uniquement « meilleur parmi les méthodes testées ici »,
non « optimal ». La composition du groupe comparé varie dans les tableaux
historiques ; pour l'évaluation principale v1.1, les cinq méthodes sont évaluées
sur les mêmes instances et le dénominateur est défini dans la
section~\ref{sec:eval-v1-1-protocole}.
\end{remark}

\begin{remark}{Décodage de la campagne pilote}{decode-stochastic}
Dans l'ancienne configuration pilote, le paramètre \texttt{decode\_type=sampling}
restait actif à l'inférence, y compris pour le mode « une passe ». Les résultats
correspondants étaient donc stochastiques, avec une variance possible entre
exécutions. Cette limite ne doit pas être transférée au protocole v1.1, qui
comprend explicitement un décodage \texttt{DRL\_greedy} et un mode
\texttt{DRL\_best\_of\_512}.
\end{remark}

\begin{remark}{Taille des échantillons de la campagne pilote}{sample-size}
Chaque configuration présentée dans les tableaux historiques de cette partie
est évaluée sur $10$ instances indépendantes. Cet effectif réduit limite la
précision des statistiques descriptives ; il ne décrit pas la taille du test
final v1.1, qui comprend $100$ instances appariées par taille. Les mesures de
temps répétées du protocole v1.1 utilisent, quant à elles, un sous-ensemble de
$10$ instances par taille avec un échauffement et cinq répétitions.
\end{remark}

\subsection{Résultats exploratoires historiques : point de contrôle au pas 200}
\label{sec:res-step200}

\subsubsection{Résultats sur $n=20$ (taille d'entraînement)}

\begin{table}[H]
	\centering
	\begin{tabular}{lccc}
		\toprule
		Méthode & Gap moyen (\%) & Écart-type (\%) & Temps/instance (s) \\
		\midrule
		DRL (1 passe)                & 12.66 & 11.20 & 0.0064 \\
		DRL (\textit{sampling}, 512) & 5.56  & 5.81  & 0.1540 \\
		Nearest Neighbor             & 17.86 & 9.87  & 0.0001 \\
		NN + 2-opt                   & 1.17  & 1.64  & 0.0008 \\
		ACO                          & 0.09  & 0.28  & 0.2277 \\
		\bottomrule
	\end{tabular}
	\caption{Comparaison sur 10 instances à $n=20$ villes, modèle entraîné non convergé (pas 200).}
	\label{tab:resultats-n20-batch}
\end{table}

\subsubsection{Généralisation à des tailles non vues ($n=50$, $n=100$)}

\begin{table}[H]
	\centering
	\begin{tabular}{llccc}
		\toprule
		Taille & Méthode & Gap moyen (\%) & Écart-type (\%) & Temps/instance (s) \\
		\midrule
		\multirow{3}{*}{$n=50$}
		& DRL (1 passe)                & 20.24 & 6.09 & 0.0107 \\
		& DRL (\textit{sampling}, 512) & 9.42  & 4.66 & 0.2040 \\
		& Nearest Neighbor             & 15.21 & 6.04 & 0.0011 \\
		& NN + 2-opt                   & 0.02  & 0.06 & 0.0126 \\
		\midrule
		\multirow{3}{*}{$n=100$}
		& DRL (1 passe)                & 41.69 & 7.34 & 0.0204 \\
		& DRL (\textit{sampling}, 512) & 28.95 & 4.88 & 0.5820 \\
		& Nearest Neighbor             & 18.26 & 3.40 & 0.0021 \\
		& NN + 2-opt                   & 0.00  & 0.00 & 0.0295 \\
		\bottomrule
	\end{tabular}
	\caption{Généralisation à $n=50$ et $n=100$, modèle entraîné non convergé (pas 200), sur $n=20$ (ACO non testé à ces tailles).}
	\label{tab:resultats-generalisation}
\end{table}

\subsection{Résultats exploratoires historiques : point de contrôle au pas 1550}
\label{sec:res-step1550}

Le point de contrôle \texttt{train20\_0815\_14\_27\_step1550\_act.pt}
(cf. Remarque~\ref{rem:entrainement-partiel}), obtenu après 1\,550 pas de reprise
à partir du point de contrôle du pas 1\,999 (soit environ 3\,550 pas cumulés),
présente des pertes stabilisées sur la séquence observée. Il est donc qualifié ici
de point de contrôle \emph{à pertes stabilisées}, plutôt que de modèle dont la
convergence globale serait démontrée. Il a été évalué à $n=20$, puis en
généralisation à $n=50$ et $n=100$.

\begin{table}[H]
	\centering
	\begin{tabular}{llccc}
		\toprule
		Taille & Méthode & Gap moyen (\%) & Écart-type (\%) & Temps/instance (s) \\
		\midrule
		\multirow{4}{*}{$n=20$}
		& DRL (\textit{sampling}, 512) & 5.68  & 4.87 & 0.0724 \\
		& Nearest Neighbor             & 17.88 & 9.88 & 0.0001 \\
		& NN + 2-opt                   & 1.19  & 1.62 & 0.0008 \\
		& ACO                          & 0.05  & 0.15 & 0.1667 \\
		\midrule
		\multirow{4}{*}{$n=50$}
		& DRL (1 passe)     & 20.49 & 5.24 & 0.0077 \\
		& Nearest Neighbor  & 16.90 & 5.98 & 0.0006 \\
		& NN + 2-opt        & 1.49  & 1.83 & 0.0072 \\
		& ACO               & 0.85  & 1.31 & 1.2347 \\
		\midrule
		\multirow{4}{*}{$n=100$}
		& DRL (1 passe)     & 36.96 & 4.63 & 0.0154 \\
		& Nearest Neighbor  & 18.27 & 3.39 & 0.0023 \\
		& NN + 2-opt        & 0.01  & 0.03 & 0.0303 \\
		& ACO               & 4.86  & 3.11 & 4.2970 \\
		\bottomrule
	\end{tabular}
	\caption{Évaluation complète (tailles $n=20, 50, 100$) du modèle entraîné dans ce travail, au point de contrôle du pas 1550.}
	\label{tab:resultats-step1550}
\end{table}

\subsection[Comparateur externe historique : poids du dépôt (pas 14 999)]{Comparateur externe historique : poids du dépôt (pas 14\,999)}
\label{sec:res-pretrained}

\begin{remark}{Modèle non entraîné dans ce travail}{ref-model}
	Les résultats de cette sous-section utilisent
	\texttt{train20\_1113\_12\_12\_step14999\_act.pt}, fourni par l'auteur du
	dépôt \texttt{TSP\_DRL\_PtrNet}. Ils ne reflètent pas un entraînement
	réalisé dans ce PFE, mais servent de repère externe indiquant le niveau de
	performance atteignable une fois le modèle de référence au pas 14\,999.
\end{remark}

\begin{table}[H]
	\centering
	\begin{tabular}{lccc}
		\toprule
		Méthode & Gap moyen (\%) & Écart-type (\%) & Temps/instance (s) \\
		\midrule
		DRL (\textit{sampling}, 512) & 2.53 & 2.08 & 0.0779 \\
		\bottomrule
	\end{tabular}
	\caption{$n=20$, modèle de référence (pas 14\,999), mode \textit{sampling} uniquement (pas de test 1 passe effectué à cette taille). Gap calculé contre les mêmes coûts NN/2-opt/ACO que le Tableau~\ref{tab:resultats-n20-batch}.}
\end{table}

\begin{table}[H]
	\centering
	\begin{tabular}{llccc}
		\toprule
		Taille & Méthode & Gap moyen (\%) & Écart-type (\%) & Temps/instance (s) \\
		\midrule
		\multirow{4}{*}{$n=50$}
		& DRL (1 passe)     & 13.02 & 4.21 & 0.0078 \\
		& Nearest Neighbor  & 17.86 & 5.64 & 0.0007 \\
		& NN + 2-opt        & 2.36  & 2.84 & 0.0074 \\
		& ACO               & 0.82  & 1.37 & 1.1283 \\
		\midrule
		\multirow{4}{*}{$n=100$}
		& DRL (1 passe)     & 24.76 & 4.18 & 0.0154 \\
		& Nearest Neighbor  & 18.26 & 3.40 & 0.0023 \\
		& NN + 2-opt        & 0.00  & 0.00 & 0.0329 \\
		& ACO               & 4.37  & 3.29 & 4.2344 \\
		\bottomrule
	\end{tabular}
	\caption{Généralisation à $n=50$ et $n=100$, modèle de référence (pas 14\,999), avec ACO.}
	\label{tab:resultats-generalisation-ref}
\end{table}

\subsection{Synthèse des évaluations exploratoires}

La Figure~\ref{fig:gaps-comparaison} synthétise les résultats de la campagne
pilote sur $10$ instances par configuration : checkpoint pilote au pas 1\,550,
heuristique PPV+2-opt, ACO et poids externes fournis avec le dépôt (pas 14\,999).
Elle doit être lue comme un bilan exploratoire historique. Les conclusions
principales de ce mémoire reposent sur le test final frais v1.1 décrit dans
les sous-sections précédentes.

\begin{figure}[H]
\centering
\includegraphics[width=0.85\textwidth]{images/gaps_comparaison.pdf}
\caption{Gap moyen (\%) dans la campagne exploratoire pilote par méthode et par
taille d'instance : checkpoint du pas 1\,550, heuristique PPV+2-opt, ACO et
modèle externe du dépôt (pas 14\,999). Résultats historiques des
Tableaux~\ref{tab:resultats-step1550} et~\ref{tab:resultats-generalisation-ref} ;
ils ne remplacent pas l'évaluation indépendante v1.1.}
\label{fig:gaps-comparaison}
\end{figure}

\section{Discussion}
\label{sec:ch4-discussion}
% ================================================================

\subsection{Analyse des résultats}

L'interprétation principale repose sur l'évaluation indépendante v1.1 et son
extension v1.2 à trois graines d'entraînement indépendantes
($20261013$, $20261014$ et $20261015$). Les trois modèles ont été entraînés
avec le même protocole, sélectionnés sur le même jeu de validation indépendant
et évalués sur les mêmes coordonnées de test à $n=20$, $50$ et $100$. Les
résultats DRL rapportés dans cette discussion sont les moyennes des trois
moyennes obtenues par modèle ; l'écart-type associé décrit la dispersion entre
graines, et non entre les instances.

À $n=20$, ACO obtient la plus faible longueur moyenne ($3{,}8780$), devant
PPV+2-opt ($3{,}9424$), DRL meilleur-sur-512 ($3{,}9584\pm0{,}0054$) et DRL
glouton ($4{,}0355\pm0{,}0052$), tandis que PPV seul atteint $4{,}4961$.
Pour DRL meilleur-sur-512 face à PPV+2-opt, les différences appariées par
modèle sont faibles ($+0{,}016923$, $+0{,}010224$ et $+0{,}020807$) et les
trois intervalles bootstrap à 95~\% croisent zéro. Ces résultats ne mettent
pas en évidence de différence moyenne claire conditionnellement à chacun des
modèles, mais ne démontrent pas l'équivalence. Les comparaisons avec ACO
favorisent ce dernier dans les trois entraînements.

À $n=50$, DRL glouton ($6{,}6403\pm0{,}0281$) et DRL meilleur-sur-512
($6{,}3294\pm0{,}0209$) obtiennent de meilleurs résultats moyens que PPV seul
($6{,}9443$), mais restent derrière PPV+2-opt ($5{,}9712$) et ACO ($6{,}1568$).
Les intervalles bootstrap conditionnels confirment la même direction pour les
trois graines : les différences DRL--PPV sont négatives, tandis que les
différences DRL--PPV+2-opt et DRL--ACO sont positives.

À $n=100$, DRL glouton ($10{,}4373\pm0{,}0364$) et DRL meilleur-sur-512
($9{,}8946\pm0{,}0293$) sont dépassés en longueur moyenne par PPV seul
($9{,}7388$), ACO ($9{,}1606$) et PPV+2-opt ($8{,}2705$). Les 18 intervalles
bootstrap conditionnels à un modèle pour les deux modes DRL et les trois
baselines sont strictement positifs. PPV+2-opt produit le meilleur coût
observé sur les 100 instances ; ce résultat reste relatif aux neuf
configurations évaluées et ne prouve pas l'optimalité exacte.

La dispersion observée entre les trois modèles reste limitée pour cette
configuration : l'écart-type inter-graines de DRL meilleur-sur-512 est de
$0{,}0054$, $0{,}0209$ et $0{,}0293$ pour $n=20$, $50$ et $100$. Ces valeurs
sont descriptives et leur précision demeure faible avec seulement trois
entraînements. Elles ne justifient ni une généralisation à toutes les
initialisations possibles ni une conclusion sur l'ensemble des méthodes DRL.

Le mode meilleur-sur-512 produit des tournées plus courtes que le décodage
glouton en moyenne, mais demande davantage de calcul. Dans le profil de temps
répété v1.1, les temps médians du DRL glouton sont d'environ $0{,}024$,
$0{,}058$ et $0{,}120$ seconde pour $n=20$, $50$ et $100$, contre environ
$0{,}055$, $0{,}186$ et $0{,}588$ seconde pour le meilleur-sur-512. PPV et
PPV+2-opt sont plus rapides aux petites tailles, tandis que l'ACO a le temps
d'exécution le plus élevé dans cette implémentation. Ces chiffres restent
liés au matériel et aux frontières de chronométrage décrites dans le protocole
v1.1 : chargement du modèle exclu, synchronisation CUDA pour le DRL et mesures
CPU comprenant l'exécution du solveur et le calcul de la longueur. Les temps
n'ont pas été remesurés pour les deux graines supplémentaires ; ils ne
quantifient donc pas une variabilité inter-graines de la latence et ne
constituent pas une comparaison universelle end-to-end.

Les tableaux aux checkpoints 200 et 1\,550 ainsi que le modèle externe au pas
14\,999 correspondent à une campagne pilote historique, fondée sur seulement
$10$ instances par configuration et, dans le cas du mode « une passe », sur un
décodage encore stochastique. Ils documentent le développement du modèle,
mais ne constituent pas la base quantitative principale des conclusions,
qui repose sur les évaluations indépendantes v1.1 et v1.2.

\subsection{Limites de l'approche}
\label{sec:ch4-limites}

\begin{itemize}
\item L'extension v1.2 porte le nombre d'entraînements indépendants à trois
($20261013$, $20261014$ et $20261015$). L'écart-type inter-graines est donc
calculé sur seulement trois moyennes de modèle ; il fournit une description
initiale de la variabilité observée, mais reste trop peu précis pour
caractériser la distribution des performances de tous les entraînements
possibles.

\item Les trois modèles ont été entraînés uniquement sur des instances de
$n=20$ villes ; la sélection du checkpoint a également été réalisée par
validation gloutonne à $n=20$. Les résultats à $n=50$ et $n=100$ mesurent le
transfert vers des tailles non vues pendant l'entraînement, sans validation
ni adaptation à ces tailles.

\item Le gap multi-graines est relatif au meilleur coût observé parmi les
neuf configurations (six sorties DRL issues des trois modèles et trois
baselines) sur chaque instance, non à l'optimum exact. Un gap de $0\,\%$
indique uniquement le meilleur résultat empirique parmi ces configurations ;
il ne prouve pas l'optimalité globale.

\item Les 54 intervalles de confiance bootstrap appariés sont calculés
séparément pour chaque modèle fixé, sur 100 instances par taille et avec
10\,000 ré-échantillonnages. Ils décrivent l'incertitude liée aux instances,
ne remplacent pas l'analyse inter-graines et ne sont pas corrigés pour la
multiplicité. Un intervalle croisant zéro ne constitue pas une preuve
d'équivalence.

\item Les temps dépendent de l'implémentation, du matériel et des frontières
de mesure du protocole v1.1. Le chargement du modèle est exclu ; le calcul
chronométré n'est pas identique côté DRL et côté CPU. Le temps GPU cumulé de
l'entraînement n'a pas été mesuré. De plus, les fichiers de temps répété v1.1
ont été réutilisés pour v1.2 et n'ont pas été remesurés pour les deux graines
supplémentaires ; aucune dispersion inter-graines des temps n'est donc
estimée.

\item Les trois évaluations utilisent PyTorch \texttt{2.11.0+cu130}, CUDA
\texttt{13.0} et Tesla~T4, mais la version corrective de Python diffère entre
la première exécution (\texttt{3.13.15}) et les deux suivantes
(\texttt{3.13.16}). Cette variation logicielle mineure est documentée et
constitue une limite de reproductibilité stricte.

\item Les tableaux pilotes antérieurs reposent sur $10$ instances par
configuration et présentent une incertitude plus élevée. Ils sont conservés
comme historique, mais ne sont pas utilisés comme estimation principale de
la performance du modèle.

\item Les hyperparamètres sont principalement issus de la configuration du
dépôt et n'ont pas fait l'objet d'une recherche systématique propre à ce
travail. Les résultats caractérisent donc cette configuration précise, et
non l'ensemble des méthodes d'apprentissage par renforcement profond
applicables au TSP.
\end{itemize}

\subsection{Perspectives}

\begin{itemize}
\item Étendre l'analyse inter-graines au-delà de trois entraînements
indépendants, afin de mieux estimer la variabilité de l'initialisation et de
l'ordre des données, tout en conservant des jeux de validation et de test
communs et indépendants.
\item Étudier des budgets d'entraînement supérieurs à $13\,000$ mises à jour
et comparer plusieurs tailles d'entraînement, notamment $n=20$, $50$ et
$100$, avec une sélection de checkpoint adaptée mais sans contamination du
test final.
\item Mesurer le coût GPU total de l'entraînement et le mettre en regard du
coût d'inférence répété pour plusieurs graines, afin de fournir une
comparaison temps/qualité plus complète avec les méthodes classiques.
\item Évaluer une approche hybride associant les tournées produites par le
Pointer Network à un post-traitement 2-opt, puis la comparer aux baselines
avec un protocole apparié et un budget de calcul explicite.
\item Étudier d'autres architectures, notamment les Transformers, avec des
protocoles d'entraînement et d'évaluation contrôlés ; toute comparaison avec
PPO ou d'autres algorithmes de renforcement devra reposer sur des expériences
spécifiques et non sur une extrapolation des résultats présents.
\end{itemize}

%  CONCLUSION GENERALE
	% ============================================================
	\chapter*{Conclusion G\'en\'erale}
	\addcontentsline{toc}{chapter}{Conclusion G\'en\'erale}
	\markboth{Conclusion G\'en\'erale}{Conclusion G\'en\'erale}

Ce mémoire avait pour objectif d'étudier la place de l'apprentissage par
renforcement profond dans la résolution du problème du voyageur de commerce,
en reliant sa formulation mathématique à une mise en œuvre expérimentale.
Le travail a d'abord présenté le TSP comme un problème d'optimisation
combinatoire sur des cycles hamiltoniens, puis rappelé les principales
familles de méthodes exactes, heuristiques et métaheuristiques. Il a ensuite
décrit la modélisation séquentielle par processus de décision markovien et le
principe de l'apprentissage de politique.

Sur le plan pratique, l'étude a adapté une implémentation PyTorch de Pointer
Network fondée sur un encodeur-décodeur LSTM et l'algorithme REINFORCE avec
baseline apprise. L'évaluation indépendante a été étendue à trois entraînements
initialisés indépendamment (graines $20261013$, $20261014$ et $20261015$), chacun
comportant $13\,000$ mises à jour de l'optimiseur. Pour chaque modèle, le
checkpoint a été sélectionné par validation gloutonne indépendante à $n=20$,
puis évalué sur les mêmes 100 instances inédites pour chacune des tailles
$n=20$, $50$ et $100$. Les méthodes comparées sont le plus proche voisin,
PPV+2-opt, ACO à paramètres fixes, DRL glouton et DRL meilleur-sur-512.

Les résultats agrégés montrent que la compétitivité du modèle dépend de la
taille des instances. À $n=20$, ACO obtient la meilleure moyenne ($3{,}8780$),
devant PPV+2-opt ($3{,}9424$) et DRL meilleur-sur-512 ($3{,}9584$). Pour ce
dernier mode, les trois intervalles bootstrap appariés par modèle face à
PPV+2-opt croisent zéro ; les données ne mettent pas en évidence de différence
moyenne claire conditionnellement aux modèles, sans prouver leur équivalence.
À $n=50$, les deux modes DRL surpassent PPV seul mais restent derrière
PPV+2-opt ($5{,}9712$) et ACO ($6{,}1568$). À $n=100$, les deux modes DRL sont
dépassés par les trois références classiques ; PPV+2-opt atteint la plus
faible longueur moyenne ($8{,}2705$) et le meilleur coût observé sur les
100 instances. Le mode meilleur-sur-512 produit des tournées plus courtes que
le décodage glouton, mais au prix d'un budget d'inférence plus élevé.

La dispersion des moyennes de modèle entre les trois graines reste limitée
pour cette configuration (écart-type inter-graines de DRL meilleur-sur-512
égal à $0{,}0054$, $0{,}0209$ et $0{,}0293$ pour $n=20$, $50$ et $100$). Elle
reste toutefois descriptive compte tenu du faible nombre de répétitions.
Les intervalles bootstrap appariés sont calculés séparément pour chaque modèle
sur les mêmes 100 instances ; ils décrivent l'incertitude inter-instances
conditionnellement au modèle et ne doivent pas être interprétés comme une
analyse groupée de 300 instances indépendantes. Tous les gaps sont relatifs
au meilleur coût observé parmi les neuf configurations comparées pour chaque
instance, et non à l'optimum exact ; aucun solveur exact n'a été exécuté sur
l'ensemble du test final.

Les tableaux aux checkpoints 200 et 1\,550 et le poids externe au pas 14\,999
proviennent d'une campagne pilote antérieure, fondée sur 10 instances par
configuration. Ils sont conservés comme historique des essais initiaux, mais
ne constituent pas la base principale des conclusions. Les temps de résolution
rapportés proviennent du profil répété v1.1 : ils décrivent cette implémentation
sur Tesla~T4, avec chargement du modèle exclu et des frontières de chronométrage
DRL/CPU différentes ; ils n'ont pas été remesurés pour les deux graines
supplémentaires. Le coût total entraînement plus inférence reste donc à
étudier.

Ces conclusions restent limitées à l'architecture, aux hyperparamètres, au
budget de $13\,000$ mises à jour et à la distribution d'instances étudiés.
Les trois modèles ont été entraînés et validés uniquement à $n=20$ ; les
résultats à $n=50$ et $n=100$ mesurent un transfert vers des tailles non vues
pendant l'entraînement. La variabilité inter-graines est estimée à partir de
trois entraînements seulement. Enfin, les évaluations partagent PyTorch
\texttt{2.11.0+cu130}, CUDA~\texttt{13.0} et Tesla~T4, mais diffèrent d'une
version corrective de Python (\texttt{3.13.15} contre \texttt{3.13.16}),
variation consignée dans les manifests.

Les prolongements les plus utiles sont donc l'augmentation du nombre de
graines indépendantes, l'étude de modèles entraînés à plusieurs tailles,
la mesure séparée du coût complet d'entraînement et d'inférence, ainsi que
l'évaluation d'une hybridation du Pointer Network avec 2-opt. D'autres
architectures, notamment les Transformers, et d'autres algorithmes tels que
PPO~\cite{schulman2017proximal} devront faire l'objet d'expériences dédiées,
avec les mêmes exigences de validation, de test indépendant et de comparaison
appariée.

En définitive, ce travail constitue une étude algorithmique expérimentale de
l'apprentissage par renforcement appliqué au TSP. Sa contribution réside dans
l'implémentation documentée, la construction d'un protocole reproductible,
l'évaluation sur trois entraînements indépendants et l'analyse critique de
leurs résultats. Ceux-ci illustrent l'intérêt de l'apprentissage de politiques
pour construire des solutions, mais ne montrent pas une supériorité générale
face aux heuristiques et métaheuristiques classiques testées.

%  ANNEXE : IMPLÉMENTATIONS PYTHON
	% ============================================================
	\appendix
% Appendix tables are numbered by appendix letter (A.1, B.1, ...).
\renewcommand{\thetable}{\thechapter.\arabic{table}}
	\chapter{Implémentations Python des algorithmes de résolution du TSP}
	\label{ann:implementations}
	\addcontentsline{toc}{chapter}{Annexe : Implémentations Python}
	
	Les implémentations présentées dans cette annexe ont été utilisées pour
	produire les résultats comparatifs du Chapitre~\ref{chap:impl_resultats}.
	Elles constituent la référence expérimentale contre laquelle les agents
	d'apprentissage par renforcement profond sont évalués.
	
	% ─────────────────────────────────────────────────────────────
	\section{Utilitaires communs}
	\label{ann:utils}
	% ─────────────────────────────────────────────────────────────
	
	\begin{lstlisting}[caption={Génération d'instances euclidiennes et calcul du coût d'une tournée.}, label={lst:utils}]
		import math, random
		
		def euclidean(p1, p2):
		return math.sqrt((p1[0]-p2[0])**2 + (p1[1]-p2[1])**2)
		
		def tour_cost(tour, dist):
		return sum(dist[tour[i]][tour[(i+1) % len(tour)]]
		for i in range(len(tour)))
		
		def generate_instance(n, seed=42):
		random.seed(seed)
		cities = [(random.uniform(0, 100),
		random.uniform(0, 100)) for _ in range(n)]
		dist   = [[euclidean(cities[i], cities[j])
		for j in range(n)] for i in range(n)]
		return cities, dist
	\end{lstlisting}
	
	% ─────────────────────────────────────────────────────────────
	\section{Croissance de l'espace de recherche}
	\label{ann:espace}
	% ─────────────────────────────────────────────────────────────
	
	\begin{lstlisting}[caption={Vérification numérique du Tableau~2.1 : nombre de tournées $(n-1)!/2$.}, label={lst:espace}]
		import math
		
		print(f"{'n':>6} | {'(n-1)!/2':>30}")
		print("-" * 40)
		for n in [5, 10, 20, 50, 100]:
		result = math.factorial(n - 1) // 2
		if n <= 10:
		print(f"{n:>6} | {result:>30,}")
		else:
		print(f"{n:>6} | {result:>30.2e}")
	\end{lstlisting}
	
	% ─────────────────────────────────────────────────────────────
	\section{Tableau de complexité Bellman-Held-Karp}
	\label{ann:bhk_table}
	% ─────────────────────────────────────────────────────────────
	
	\begin{lstlisting}[caption={Calcul automatique des temps et mémoires du Tableau~2.2.}, label={lst:bhk_table}]
		import math
		
		def format_time(n):
		seconds = (n**2) * (2**n) / 1e9
		if seconds < 0.001:      return "< 1 ms"
		elif seconds < 1:        return f"~{seconds:.2f} s"
		elif seconds < 60:       return f"~{seconds:.0f} s"
		elif seconds < 3600*24*365:
		return f"~{seconds/(3600*24*365):.1f} ans"
		else:
		return f"~10^{math.floor(math.log10(seconds/31557600))} ans"
		
		def format_memory(n):
		b = n * (2**n) * 8
		if b < 1024**2:   return f"~{b/1024:.0f} Ko"
		elif b < 1024**3: return f"~{b/1024**2:.0f} Mo"
		elif b < 1024**4: return f"~{b/1024**3:.0f} Go"
		else:             return "hors d'atteinte"
		
		for n in [10, 20, 30, 50, 100]:
		ops = (n**2) * (2**n)
		print(f"n={n:>3} | ops={ops:.2e} | "
		f"temps={format_time(n):>22} | "
		f"mem={format_memory(n)}")
	\end{lstlisting}
	
	% ─────────────────────────────────────────────────────────────
	\section{Méthode exacte --- Bellman-Held-Karp}
	\label{ann:bhk}
	% ─────────────────────────────────────────────────────────────
	
	\begin{lstlisting}[caption={Programmation dynamique de Bellman-Held-Karp. Complexité $\mathcal{O}(n^2 2^n)$ en temps, $\mathcal{O}(n\,2^n)$ en espace.}, label={lst:bhk}]
		def bellman_held_karp(dist):
		n   = len(dist)
		INF = float('inf')
		dp     = [[INF] * n for _ in range(1 << n)]
		parent = [[-1]  * n for _ in range(1 << n)]
		dp[1][0] = 0
		
		for S in range(1, 1 << n):
		for u in range(n):
		if not (S >> u & 1) or dp[S][u] == INF:
		continue
		for v in range(n):
		if S >> v & 1:
		continue
		S_new    = S | (1 << v)
		new_cost = dp[S][u] + dist[u][v]
		if new_cost < dp[S_new][v]:
		dp[S_new][v]     = new_cost
		parent[S_new][v] = u
		
		S_full = (1 << n) - 1
		cost, last = INF, -1
		for u in range(1, n):
		candidate = dp[S_full][u] + dist[u][0]
		if candidate < cost:
		cost, last = candidate, u
		
		tour, S, curr = [], S_full, last
		while curr != -1:
		tour.append(curr)
		prev = parent[S][curr]
		S    = S ^ (1 << curr)
		curr = prev
		tour.reverse()
		tour.append(tour[0])
		return cost, tour
	\end{lstlisting}
	
	% ─────────────────────────────────────────────────────────────
	\section{Heuristiques}
	\label{ann:heuristiques}
	% ─────────────────────────────────────────────────────────────
	
	\begin{lstlisting}[caption={Plus proche voisin (Nearest Neighbor). Complexité $\mathcal{O}(n^2)$.}, label={lst:nn}]
		def nearest_neighbor(dist, start=0):
		n       = len(dist)
		visited = [False] * n
		tour    = [start]
		visited[start] = True
		for _ in range(n - 1):
		current = tour[-1]
		nearest = min(
		(j for j in range(n) if not visited[j]),
		key=lambda j: dist[current][j]
		)
		tour.append(nearest)
		visited[nearest] = True
		return tour
	\end{lstlisting}
	
	\begin{lstlisting}[caption={Amélioration 2-opt. Complexité $\mathcal{O}(n^2)$ par itération.}, label={lst:2opt}]
		def two_opt(tour, dist):
		n, best, improved = len(tour), tour[:], True
		while improved:
		improved = False
		for i in range(1, n - 1):
		for j in range(i + 1, n):
		before = (dist[best[i-1]][best[i]]
		+ dist[best[j]][best[(j+1) % n]])
		after  = (dist[best[i-1]][best[j]]
		+ dist[best[i]][best[(j+1) % n]])
		if after < before - 1e-10:
		best[i:j+1] = best[i:j+1][::-1]
		improved = True
		return best
	\end{lstlisting}
	
	\begin{lstlisting}[caption={Recherche locale 3-opt simplifiée. Complexité cubique pour l'examen des triplets, par balayage.}, label={lst:3opt}]
		def three_opt_segment(tour, dist):
		n, best, improved = len(tour), tour[:], True
		while improved:
		improved = False
		for i in range(n - 2):
		for j in range(i + 1, n - 1):
		for k in range(j + 1, n):
		A, B = best[i],      best[(i+1)%n]
		C, D = best[j],      best[(j+1)%n]
		E, F = best[k],      best[(k+1)%n]
		d0 = dist[A][B]+dist[C][D]+dist[E][F]
		d1 = dist[A][C]+dist[B][D]+dist[E][F]
		d2 = dist[A][B]+dist[C][E]+dist[D][F]
		d3 = dist[A][D]+dist[E][B]+dist[C][F]
		if d1 < d0 - 1e-10:
		best[i+1:j+1] = best[i+1:j+1][::-1]
		improved = True
		elif d2 < d0 - 1e-10:
		best[j+1:k+1] = best[j+1:k+1][::-1]
		improved = True
		elif d3 < d0 - 1e-10:
		best[i+1:j+1] = best[i+1:j+1][::-1]
		best[j+1:k+1] = best[j+1:k+1][::-1]
		improved = True
		return best
	\end{lstlisting}
	
	% ─────────────────────────────────────────────────────────────
	\section{Métaheuristiques}
	\label{ann:metaheuristiques}
	% ─────────────────────────────────────────────────────────────
	
	\begin{lstlisting}[caption={Recuit simulé (Simulated Annealing).}, label={lst:sa}]
		def simulated_annealing(dist, T_init=1000.0, T_min=1e-3,
		alpha=0.995, n_iter_per_temp=100):
		n            = len(dist)
		current_tour = nearest_neighbor(dist)
		current_cost = tour_cost(current_tour, dist)
		best_tour, best_cost = current_tour[:], current_cost
		T = T_init
		while T > T_min:
		for _ in range(n_iter_per_temp):
		i, j     = sorted(random.sample(range(n), 2))
		neighbor = current_tour[:]
		neighbor[i:j+1] = neighbor[i:j+1][::-1]
		delta = tour_cost(neighbor, dist) - current_cost
		if delta < 0 or random.random() < math.exp(-delta/T):
		current_tour = neighbor
		current_cost = tour_cost(neighbor, dist)
		if current_cost < best_cost:
		best_cost = current_cost
		best_tour = current_tour[:]
		T *= alpha
		return best_tour, best_cost
	\end{lstlisting}
	
	\begin{lstlisting}[caption={Algorithme génétique avec croisement OX1 et sélection par tournoi.}, label={lst:ga}]
		def genetic_algorithm(dist, pop_size=100, n_gen=500,
		mutation_rate=0.02, tournament_k=5):
		n          = len(dist)
		population = [random.sample(range(n), n)
		for _ in range(pop_size)]
		
		def tournament(pop):
		return min(random.sample(pop, tournament_k),
		key=lambda t: tour_cost(t, dist))
		
		def ox1_crossover(p1, p2):
		a, b  = sorted(random.sample(range(n), 2))
		child = [None] * n
		child[a:b+1] = p1[a:b+1]
		fill  = [v for v in p2 if v not in set(p1[a:b+1])]
		idx   = 0
		for k in range(n):
		if child[k] is None:
		child[k] = fill[idx]; idx += 1
		return child
		
		best_tour = min(population,
		key=lambda t: tour_cost(t, dist))
		best_cost = tour_cost(best_tour, dist)
		for _ in range(n_gen):
		new_pop = [best_tour[:]]
		while len(new_pop) < pop_size:
		child = ox1_crossover(tournament(population),
		tournament(population))
		if random.random() < mutation_rate:
		i, j = random.sample(range(n), 2)
		child[i], child[j] = child[j], child[i]
		new_pop.append(child)
		population = new_pop
		gen_best   = min(population,
		key=lambda t: tour_cost(t, dist))
		if tour_cost(gen_best, dist) < best_cost:
		best_cost = tour_cost(gen_best, dist)
		best_tour = gen_best[:]
		return best_tour, best_cost
	\end{lstlisting}
	
	\begin{lstlisting}[caption={Optimisation par colonies de fourmis (ACO).}, label={lst:aco}]
		def ant_colony(dist, n_ants=20, n_iter=100,
		alpha=1.0, beta=2.0, rho=0.5):
		n   = len(dist)
		eta = [[0 if i==j else 1/dist[i][j] for j in range(n)]
		for i in range(n)]
		tau = [[1.0]*n for _ in range(n)]
		best_tour, best_cost = None, float('inf')
		for _ in range(n_iter):
		tours, costs = [], []
		for _ in range(n_ants):
		start   = random.randint(0, n-1)
		visited = [False]*n
		tour    = [start]
		visited[start] = True
		for _ in range(n-1):
		cur   = tour[-1]
		probs = [0 if visited[j] else
		(tau[cur][j]**alpha)*(eta[cur][j]**beta)
		for j in range(n)]
		total = sum(probs)
		probs = [p/total for p in probs]
		r, cumul, nxt = random.random(), 0.0, n-1
		for j, p in enumerate(probs):
		cumul += p
		if r <= cumul:
		nxt = j; break
		tour.append(nxt); visited[nxt] = True
		c = tour_cost(tour, dist)
		tours.append(tour); costs.append(c)
		if c < best_cost:
		best_cost, best_tour = c, tour[:]
		for i in range(n):
		for j in range(n):
		tau[i][j] *= (1-rho)
		for tour, cost in zip(tours, costs):
		d = 1.0/cost
		for k in range(n):
		i, j = tour[k], tour[(k+1)%n]
		tau[i][j] += d; tau[j][i] += d
		return best_tour, best_cost
	\end{lstlisting}
	
	% ─────────────────────────────────────────────────────────────
	\section{Comparaison globale des méthodes}
	\label{ann:comparaison}
	% ─────────────────────────────────────────────────────────────
	
	\begin{lstlisting}[caption={Tableau comparatif de toutes les méthodes sur une instance euclidienne à $n$ villes.}, label={lst:comparaison}]
		import time
		
		def optimality_gap(cost, optimal):
		return 100 * (cost - optimal) / optimal
		
		n = 15
		cities, dist = generate_instance(n, seed=42)
		cost_opt, _  = bellman_held_karp(dist)
		results = {}
		
		for label, func in [
		("Nearest Neighbor", lambda: (nearest_neighbor(dist), None)),
		]:
		t0 = time.time()
		tour_nn = nearest_neighbor(dist)
		results["Nearest Neighbor"] = (tour_cost(tour_nn, dist),
		time.time()-t0)
		t0 = time.time()
		results["NN + 2-opt"] = (tour_cost(two_opt(tour_nn[:], dist),
		dist), time.time()-t0)
		t0 = time.time()
		results["NN + 3-opt"] = (tour_cost(
		three_opt_segment(tour_nn[:], dist), dist), time.time()-t0)
		t0 = time.time()
		_, c = simulated_annealing(dist)
		results["Recuit simule"] = (c, time.time()-t0)
		t0 = time.time()
		_, c = genetic_algorithm(dist)
		results["Algo. genetique"] = (c, time.time()-t0)
		t0 = time.time()
		_, c = ant_colony(dist)
		results["ACO"] = (c, time.time()-t0)
		
		print(f"n={n} | Optimal BHK : {cost_opt:.4f}\n")
		print(f"{'Methode':<22}|{'Cout':>10}|{'Gap(%)':>8}|{'Temps':>10}")
		print("-"*55)
		for m, (cost, t) in results.items():
		g = optimality_gap(cost, cost_opt)
		print(f"{m:<22}|{cost:>10.4f}|{g:>7.2f}%|{t:>10.4f}")
	\end{lstlisting}
	
	% ─────────────────────────────────────────────────────────────
	\section{Génération de la figure de complexité}
	\label{ann:figures}
	% ─────────────────────────────────────────────────────────────
	
	\begin{lstlisting}[caption={Script de génération de la Figure~\ref{fig:complexite_comparative}.}, label={lst:figures}]
		import numpy as np
		import matplotlib.pyplot as plt, math
		
		n_vals    = np.arange(1, 25)
		n_log_n   = n_vals * np.log2(np.maximum(n_vals, 1))
		n_squared = n_vals ** 2
		bhk       = (n_vals**2) * (2.0**n_vals)
		factorial = np.array([math.factorial(int(n))
		for n in n_vals], dtype=float)
		
		fig, axes = plt.subplots(1, 2, figsize=(12, 4.5))
		ax = axes[0]
		ax.semilogy(n_vals, n_log_n,   color='green',     lw=2,
		label=r'$\mathcal{O}(n\log n)$')
		ax.semilogy(n_vals, n_squared, color='steelblue',  lw=2,
		label=r'$\mathcal{O}(n^2)$')
		ax.semilogy(n_vals, bhk,       color='darkorange', lw=2,
		label=r'$\mathcal{O}(n^2 2^n)$ -- BHK')
		ax.semilogy(n_vals[:15], factorial[:15], color='crimson',
		lw=2, linestyle='--',
		label=r'$\mathcal{O}(n!)$ -- Force brute')
		ax.set_xlabel("Taille $n$"); ax.legend(fontsize=9)
		ax.grid(True, which='both', linestyle='--', alpha=0.4)
		
		n_zoom = np.arange(5, 21)
		ax2 = axes[1]
		ax2.semilogy(n_zoom, (n_zoom**2)*(2.0**n_zoom),
		color='darkorange', lw=2, label='BHK')
		ax2.semilogy(n_zoom,
		np.array([math.factorial(int(n))
		for n in n_zoom], dtype=float),
		color='crimson', lw=2, linestyle='--',
		label='Force brute')
		ax2.axvline(x=20, color='gray', linestyle=':', lw=1.5)
		ax2.set_title("Zoom : BHK vs Force brute")
		ax2.legend(fontsize=9)
		ax2.grid(True, which='both', linestyle='--', alpha=0.4)
		
		plt.tight_layout()
		plt.savefig('images/complexite_comparative.pdf',
		bbox_inches='tight', dpi=300)
	\end{lstlisting}
	
	
	
	
	
	
	
	% ============================================================
	%  BIBLIOGRAPHIE
	% ============================================================
\chapter{Comparaisons appariées par bootstrap selon la graine d'entraînement}
\label{app:multiseed-bootstrap}

Le tableau suivant rassemble les 54 comparaisons appariées de l'évaluation
v1.2. Chaque différence est définie par
$\Delta=L_{\mathrm{DRL}}-L_{\text{référence}}$ sur les mêmes 100 instances
pour une taille, une graine d'entraînement et un mode de décodage donnés.
Une valeur négative favorise le DRL. Les intervalles sont des intervalles
percentiles bootstrap à 95~\% calculés avec 10\,000 ré-échantillonnages par
comparaison. Ils sont conditionnels au modèle évalué et ne quantifient pas
la variabilité entre graines. Aucune correction de multiplicité n'est
appliquée.

\scriptsize
\begin{longtable}{r r l l r r r}
\caption{Différences appariées DRL--référence et intervalles bootstrap à 95~\% par modèle.}\label{tab:multiseed-paired-bootstrap}\\
\toprule
$n$ & Graine & Mode DRL & Référence & $\Delta$ & IC inf. & IC sup. \\
\midrule
\endfirsthead
\multicolumn{7}{c}{\tablename\ \thetable{} -- suite}\\
\toprule
$n$ & Graine & Mode DRL & Référence & $\Delta$ & IC inf. & IC sup. \\
\midrule
\endhead
\midrule
\multicolumn{7}{r}{Suite à la page suivante}\\
\endfoot
\bottomrule
\endlastfoot
20 & 20261013 & meilleur-sur-512 & ACO & +0,081285 & +0,059004 & +0,104979 \\
20 & 20261013 & meilleur-sur-512 & PPV & -0,536787 & -0,613731 & -0,459909 \\
20 & 20261013 & meilleur-sur-512 & PPV+2-opt & +0,016923 & -0,017012 & +0,049661 \\
20 & 20261013 & glouton & ACO & +0,158482 & +0,129191 & +0,190806 \\
20 & 20261013 & glouton & PPV & -0,459591 & -0,536047 & -0,383334 \\
20 & 20261013 & glouton & PPV+2-opt & +0,094120 & +0,055884 & +0,134430 \\
20 & 20261014 & meilleur-sur-512 & ACO & +0,074586 & +0,053218 & +0,098076 \\
20 & 20261014 & meilleur-sur-512 & PPV & -0,543487 & -0,622364 & -0,468003 \\
20 & 20261014 & meilleur-sur-512 & PPV+2-opt & +0,010224 & -0,021014 & +0,041364 \\
20 & 20261014 & glouton & ACO & +0,162067 & +0,130185 & +0,195114 \\
20 & 20261014 & glouton & PPV & -0,456006 & -0,538579 & -0,371562 \\
20 & 20261014 & glouton & PPV+2-opt & +0,097705 & +0,060080 & +0,136566 \\
20 & 20261015 & meilleur-sur-512 & ACO & +0,085169 & +0,060763 & +0,110810 \\
20 & 20261015 & meilleur-sur-512 & PPV & -0,532904 & -0,609507 & -0,455226 \\
20 & 20261015 & meilleur-sur-512 & PPV+2-opt & +0,020807 & -0,012391 & +0,053558 \\
20 & 20261015 & glouton & ACO & +0,151902 & +0,122139 & +0,183281 \\
20 & 20261015 & glouton & PPV & -0,466171 & -0,548174 & -0,387399 \\
20 & 20261015 & glouton & PPV+2-opt & +0,087540 & +0,051356 & +0,124573 \\
50 & 20261013 & meilleur-sur-512 & ACO & +0,189701 & +0,131827 & +0,248989 \\
50 & 20261013 & meilleur-sur-512 & PPV & -0,597777 & -0,702518 & -0,493301 \\
50 & 20261013 & meilleur-sur-512 & PPV+2-opt & +0,375350 & +0,319424 & +0,431710 \\
50 & 20261013 & glouton & ACO & +0,515168 & +0,441707 & +0,589007 \\
50 & 20261013 & glouton & PPV & -0,272310 & -0,385350 & -0,161928 \\
50 & 20261013 & glouton & PPV+2-opt & +0,700817 & +0,628162 & +0,776003 \\
50 & 20261014 & meilleur-sur-512 & ACO & +0,149255 & +0,095867 & +0,202205 \\
50 & 20261014 & meilleur-sur-512 & PPV & -0,638223 & -0,733688 & -0,542195 \\
50 & 20261014 & meilleur-sur-512 & PPV+2-opt & +0,334904 & +0,279100 & +0,391240 \\
50 & 20261014 & glouton & ACO & +0,461323 & +0,398326 & +0,523360 \\
50 & 20261014 & glouton & PPV & -0,326155 & -0,431095 & -0,225027 \\
50 & 20261014 & glouton & PPV+2-opt & +0,646972 & +0,582731 & +0,713520 \\
50 & 20261015 & meilleur-sur-512 & ACO & +0,178886 & +0,125552 & +0,233041 \\
50 & 20261015 & meilleur-sur-512 & PPV & -0,608592 & -0,703563 & -0,514731 \\
50 & 20261015 & meilleur-sur-512 & PPV+2-opt & +0,364535 & +0,308658 & +0,421228 \\
50 & 20261015 & glouton & ACO & +0,474028 & +0,410211 & +0,537145 \\
50 & 20261015 & glouton & PPV & -0,313450 & -0,409695 & -0,218313 \\
50 & 20261015 & glouton & PPV+2-opt & +0,659677 & +0,594836 & +0,724125 \\
100 & 20261013 & meilleur-sur-512 & ACO & +0,767604 & +0,686702 & +0,851001 \\
100 & 20261013 & meilleur-sur-512 & PPV & +0,189396 & +0,082336 & +0,296867 \\
100 & 20261013 & meilleur-sur-512 & PPV+2-opt & +1,657701 & +1,583517 & +1,731271 \\
100 & 20261013 & glouton & ACO & +1,307510 & +1,214654 & +1,404016 \\
100 & 20261013 & glouton & PPV & +0,729301 & +0,614851 & +0,844684 \\
100 & 20261013 & glouton & PPV+2-opt & +2,197607 & +2,111417 & +2,282387 \\
100 & 20261014 & meilleur-sur-512 & ACO & +0,714292 & +0,639715 & +0,790425 \\
100 & 20261014 & meilleur-sur-512 & PPV & +0,136084 & +0,031274 & +0,245024 \\
100 & 20261014 & meilleur-sur-512 & PPV+2-opt & +1,604390 & +1,537374 & +1,673266 \\
100 & 20261014 & glouton & ACO & +1,286152 & +1,199639 & +1,373863 \\
100 & 20261014 & glouton & PPV & +0,707944 & +0,592178 & +0,822329 \\
100 & 20261014 & glouton & PPV+2-opt & +2,176250 & +2,097338 & +2,257752 \\
100 & 20261015 & meilleur-sur-512 & ACO & +0,719907 & +0,644000 & +0,796924 \\
100 & 20261015 & meilleur-sur-512 & PPV & +0,141699 & +0,033925 & +0,246811 \\
100 & 20261015 & meilleur-sur-512 & PPV+2-opt & +1,610005 & +1,541183 & +1,677590 \\
100 & 20261015 & glouton & ACO & +1,236475 & +1,148335 & +1,324218 \\
100 & 20261015 & glouton & PPV & +0,658267 & +0,539794 & +0,774696 \\
100 & 20261015 & glouton & PPV+2-opt & +2,126573 & +2,045353 & +2,207810 \\
\end{longtable}
\normalsize

	\cleardoublepage
	\phantomsection
	\addcontentsline{toc}{chapter}{Bibliographie}
	\nocite{*}
	\bibliography{references}
	
\end{document}
% ================================================================
%  FIN
% ================================================================